% Interaction électrostatique et champ électrique — Physique Tle D — Cours signé OnBuch+
% Programme officiel MINESEC. Compiler avec : tectonic P03-interaction-electrostatique-champ-electrique.tex
\newcommand{\DOCMATIERE}{Physique}
\newcommand{\DOCNIVEAU}{Tle D}
\newcommand{\DOCMODULE}{Module 2 : Mouvements et interactions — évolutions temporelles des systèmes mécaniques}
\newcommand{\DOCLECON}{P03}
\newcommand{\DOCTITRE}{Interaction électrostatique et champ électrique}
% OnBuch+ — préambule « cours signés » ENRICHI (compilé avec tectonic / XeLaTeX)
% Système visuel « Pop » d'OnBuch+ : fond crème (jamais blanc), bordures encre
% épaisses, ombres « dures » décalées, titres Archivo Black en capitales.
% Version MATHÉMATIQUES : graphes (pgfplots/tikz), boîtes pédagogiques
% (définition, propriété, méthode, attention, le savais-tu, exercice résolu,
% exercice, TP, bilan), page de garde et signature OnBuch+ ancrée en bas.
%
% Macros attendues dans main.tex AVANT \input{preamble} :
%   \DOCMATIERE  \DOCNIVEAU  \DOCMODULE  \DOCLECON (numéro)  \DOCTITRE
\documentclass[11pt]{article}

\usepackage{fontspec}
\usepackage[french]{babel}
\usepackage[a4paper,margin=2cm,top=3.4cm,bottom=2.8cm,headheight=1.4cm,footskip=1.3cm]{geometry}
\usepackage{amsmath,amssymb,amsfonts}
\usepackage{mathtools} % \xmapsto, \coloneqq... souvent utilisés par les modèles
\usepackage{cancel} % simplifications barrées (bilans de forces)
% \abs{z} (module, valeur absolue) et \norm{u} (norme), utilisés par les modèles
\providecommand{\abs}{}\let\abs\relax\DeclarePairedDelimiter{\abs}{\lvert}{\rvert}
\providecommand{\norm}{}\let\norm\relax\DeclarePairedDelimiter{\norm}{\lVert}{\rVert}
\usepackage[table]{xcolor} % [table] : fournit aussi \rowcolors (alternance de lignes)
\usepackage{pagecolor}
\usepackage{tikz}
\usetikzlibrary{arrows.meta,positioning,calc,shapes.geometric,shapes.misc,shapes.arrows,decorations.pathmorphing,decorations.markings,patterns,shadows,mindmap,trees,fit,backgrounds,matrix,intersections}
\usepackage{pgfplots}
\usepackage[version=4]{mhchem} % équations nucléaires \ce{^{235}_{92}U}
\usepackage[european,siunitx]{circuitikz} % schémas électriques
\pgfplotsset{compat=1.18}
\usepgfplotslibrary{fillbetween} % aires entre courbes (calcul intégral)
\usepackage{enumitem}
\usepackage{fancyhdr}
% [most] charge toutes les bibliothèques tcolorbox (listings, minted, raster,
% documentation...) et gonfle inutilement la mémoire TeX sur un document long
% (~300 boîtes) : on ne charge que ce qui est réellement utilisé.
\usepackage[skins,breakable]{tcolorbox}
\usepackage{graphicx}
\usepackage{colortbl}
\usepackage{booktabs}
\usepackage{tabularx}
\usepackage{diagbox} % case d'en-tête diagonale des tableaux à double entrée
% Colonnes extensibles centrée / gauche / droite (C, L, R), souvent utilisées par les modèles
\newcolumntype{C}{>{\centering\arraybackslash}X}
\newcolumntype{L}{>{\raggedright\arraybackslash}X}
\newcolumntype{R}{>{\raggedleft\arraybackslash}X}
\usepackage{array}
\usepackage{multicol}
\usepackage{siunitx}
% parse-numbers=false : accepte aussi \SI{2,5 \times 10^{-3}}{...}
\sisetup{locale=FR,output-decimal-marker={,},group-minimum-digits=4,parse-numbers=false}
\DeclareSIUnit\division{div} % réglages d'oscilloscope (ms/div, V/div)
\usepackage{qrcode}
% \degree : commande fréquemment utilisée par les modèles pour un angle ou une
% température en dehors de \SI{}{} (non fournie par les paquets chargés).
\providecommand{\degree}{^{\circ}}
% \qed (fin de démonstration), utilisé par les modèles sans amsthm
\providecommand{\qed}{\hfill$\square$}
% \barre : double barre des valeurs interdites dans un tableau de variations (tkz-tab)
\providecommand{\barre}{\ensuremath{\|}}
% environnement proof (amsthm non chargé) utilisé par les modèles
\ifdefined\proof\else\newenvironment{proof}[1][Démonstration]{\par\noindent\textit{#1.}\ }{\qed\par}\fi
\usepackage{lastpage}
\usepackage{hyperref}
\hypersetup{hidelinks}

% ============================================================
% Palette « Pop » OnBuch+ — mêmes hex que la classe P (plus_ui.dart)
% ============================================================
\definecolor{poporange}{HTML}{F59321}
\definecolor{popdark}{HTML}{E07A0C}
\definecolor{popcream}{HTML}{FFF6E8}
\definecolor{popink}{HTML}{1C1714}
\definecolor{poppurple}{HTML}{7A5AE0}
\definecolor{popgreen}{HTML}{2E9E6A}
\definecolor{poppink}{HTML}{E0457B}
\definecolor{popblue}{HTML}{2F7DE1}
\definecolor{popgold}{HTML}{C98A12}
\definecolor{popmuted}{HTML}{8A7F76}
\definecolor{popline}{HTML}{EADFD2}
\definecolor{popgray}{HTML}{8A7F76}
\colorlet{popgrey}{popgray}
% Alias tolérants (le modèle nomme parfois les couleurs différemment)
\colorlet{popwhite}{white}
\colorlet{popblack}{popink}
\colorlet{popred}{poppink}
\colorlet{popyellow}{popgold}
% Teintes claires pour fonds de tableaux / zones
\colorlet{popblueL}{popblue!10!white}
\colorlet{poporangeL}{poporange!14!white}
\colorlet{popgreenL}{popgreen!12!white}
\colorlet{poppurpleL}{poppurple!12!white}
\colorlet{poppinkL}{poppink!12!white}
\colorlet{popgoldL}{popgold!14!white}

\pagecolor{popcream}

% ============================================================
% Typographie
% ============================================================
\defaultfontfeatures{Ligatures=TeX}
\setmainfont{PlusJakartaSans-Medium.ttf}[Path=../../../fonts/,BoldFont=PlusJakartaSans-Bold.ttf]
% Maths en sans-serif (Fira Math) avec chiffres et lettres droites de Plus
% Jakarta, pour que les formules se fondent dans le texte.
\usepackage[math-style=ISO,bold-style=ISO]{unicode-math}
\setmathfont{FiraMath-Regular.otf}[Path=../../../fonts/]
\setmathfont{PlusJakartaSans-Medium.ttf}[Path=../../../fonts/,range={up/num,up/{Latin,latin},\mathup}]
\usepackage{icomma} % virgule décimale sans espace en mode math
% Caractères mathématiques Unicode absents des polices de texte (Plus Jakarta,
% Archivo Black) : rendus par leur équivalent mathématique, en texte comme en
% formule — sinon ils s'affichent en carré vide (ex. « Arithmétique dans ℤ »).
\usepackage{newunicodechar}
\newunicodechar{ℝ}{\ensuremath{\mathbb{R}}}
\newunicodechar{ℤ}{\ensuremath{\mathbb{Z}}}
\newunicodechar{ℂ}{\ensuremath{\mathbb{C}}}
\newunicodechar{ℕ}{\ensuremath{\mathbb{N}}}
\newunicodechar{ℚ}{\ensuremath{\mathbb{Q}}}
\newunicodechar{𝔻}{\ensuremath{\mathbb{D}}}
\newunicodechar{α}{\ensuremath{\alpha}}
\newunicodechar{β}{\ensuremath{\beta}}
\newunicodechar{γ}{\ensuremath{\gamma}}
\newunicodechar{δ}{\ensuremath{\delta}}
\newunicodechar{ε}{\ensuremath{\varepsilon}}
\newunicodechar{θ}{\ensuremath{\theta}}
\newunicodechar{λ}{\ensuremath{\lambda}}
\newunicodechar{μ}{\ensuremath{\mu}}
\newunicodechar{σ}{\ensuremath{\sigma}}
\newunicodechar{τ}{\ensuremath{\tau}}
\newunicodechar{φ}{\ensuremath{\varphi}}
\newunicodechar{ψ}{\ensuremath{\psi}}
\newunicodechar{ω}{\ensuremath{\omega}}
\newunicodechar{Σ}{\ensuremath{\Sigma}}
% majuscules produites par \MakeUppercase dans les titres (α-aminés -> Α)
\newunicodechar{Α}{\ensuremath{\alpha}}
\newunicodechar{Β}{\ensuremath{\beta}}
\newunicodechar{Γ}{\ensuremath{\Gamma}}
\newunicodechar{Δ}{\ensuremath{\Delta}}
\newunicodechar{Ε}{\ensuremath{\varepsilon}}
\newunicodechar{Θ}{\ensuremath{\Theta}}
\newunicodechar{Λ}{\ensuremath{\Lambda}}
\newunicodechar{Μ}{\ensuremath{\mu}}
\newunicodechar{Τ}{\ensuremath{\tau}}
\newunicodechar{Φ}{\ensuremath{\Phi}}
\newunicodechar{Ψ}{\ensuremath{\Psi}}
\newunicodechar{Ω}{\ensuremath{\Omega}}
\newunicodechar{↦}{\ensuremath{\mapsto}}
\newunicodechar{∈}{\ensuremath{\in}}
\newunicodechar{∉}{\ensuremath{\notin}}
\newunicodechar{∧}{\ensuremath{\wedge}}
\newunicodechar{⇔}{\ensuremath{\Leftrightarrow}}
\newunicodechar{⟺}{\ensuremath{\Longleftrightarrow}}
\newunicodechar{⇒}{\ensuremath{\Rightarrow}}
\newunicodechar{⟹}{\ensuremath{\Longrightarrow}}
\newunicodechar{∘}{\ensuremath{\circ}}
\newunicodechar{∪}{\ensuremath{\cup}}
\newunicodechar{∩}{\ensuremath{\cap}}
\newunicodechar{≡}{\ensuremath{\equiv}}
\newunicodechar{⊂}{\ensuremath{\subset}}
\newunicodechar{⊥}{\ensuremath{\perp}}
\newunicodechar{∃}{\ensuremath{\exists}}
\newunicodechar{∀}{\ensuremath{\forall}}
\newunicodechar{∛}{\ensuremath{\sqrt[3]{\,}}}
\newunicodechar{∜}{\ensuremath{\sqrt[4]{\,}}}
\newunicodechar{⃗}{\ensuremath{{}^{\to}}}
\newunicodechar{ⁿ}{\ifmmode^{n}\else\textsuperscript{\ensuremath{n}}\fi}
\newunicodechar{ˣ}{\ifmmode^{x}\else\textsuperscript{\ensuremath{x}}\fi}
\newunicodechar{ᵇ}{\ifmmode^{b}\else\textsuperscript{\ensuremath{b}}\fi}
\newunicodechar{ʳ}{\ifmmode^{r}\else\textsuperscript{\ensuremath{r}}\fi}
\newunicodechar{ᵘ}{\ifmmode^{u}\else\textsuperscript{\ensuremath{u}}\fi}
\newunicodechar{ʸ}{\ifmmode^{y}\else\textsuperscript{\ensuremath{y}}\fi}
\newunicodechar{ᵃ}{\ifmmode^{a}\else\textsuperscript{\ensuremath{a}}\fi}
\newunicodechar{ᵉ}{\ifmmode^{e}\else\textsuperscript{\ensuremath{e}}\fi}
\newunicodechar{ᵏ}{\ifmmode^{k}\else\textsuperscript{\ensuremath{k}}\fi}
\newunicodechar{ᵖ}{\ifmmode^{p}\else\textsuperscript{\ensuremath{p}}\fi}
\newunicodechar{ᴺ}{\ifmmode^{N}\else\textsuperscript{\ensuremath{N}}\fi}
\newunicodechar{ᶜ}{\ifmmode^{c}\else\textsuperscript{\ensuremath{c}}\fi}
\newunicodechar{ʰ}{\ifmmode^{h}\else\textsuperscript{\ensuremath{h}}\fi}
\newunicodechar{ᵠ}{\ifmmode^{q}\else\textsuperscript{\ensuremath{q}}\fi}
\newunicodechar{ₙ}{\ifmmode_{n}\else\textsubscript{\ensuremath{n}}\fi}
\newunicodechar{ₐ}{\ifmmode_{a}\else\textsubscript{\ensuremath{a}}\fi}
\newunicodechar{ᵢ}{\ifmmode_{i}\else\textsubscript{\ensuremath{i}}\fi}
\newunicodechar{ᵣ}{\ifmmode_{r}\else\textsubscript{\ensuremath{r}}\fi}
\newunicodechar{ₖ}{\ifmmode_{k}\else\textsubscript{\ensuremath{k}}\fi}
\newunicodechar{ᵦ}{\ifmmode_{\beta}\else\textsubscript{\ensuremath{\beta}}\fi}
\newunicodechar{ₓ}{\ifmmode_{x}\else\textsubscript{\ensuremath{x}}\fi}
\newfontfamily\popdisplay{ArchivoBlack-Regular.ttf}[Path=../../../fonts/]
\newfontfamily\poplogo{SpaceGrotesk-Bold.ttf}[Path=../../../fonts/]
\newfontfamily\popsemi{PlusJakartaSans-SemiBold.ttf}[Path=../../../fonts/]
\newfontfamily\popxbold{PlusJakartaSans-ExtraBold.ttf}[Path=../../../fonts/]
\newfontfamily\popxboldls{PlusJakartaSans-ExtraBold.ttf}[Path=../../../fonts/,LetterSpace=3.0]

\setlist[enumerate]{leftmargin=1.7em,itemsep=5pt,topsep=4pt}
\setlist[itemize]{leftmargin=1.7em,itemsep=4pt,topsep=4pt,label={\color{poporange}\textbullet}}
\setlength{\parindent}{0pt}
\setlength{\parskip}{5pt}
\renewcommand{\arraystretch}{1.25}

% pgfplots : style Pop par défaut
\pgfplotsset{
  every axis/.append style={axis line style={popink,line width=1pt},tick style={popink},
    grid style={popline,line width=0.5pt},label style={font=\small},tick label style={font=\footnotesize}},
  popcurve/.style={poporange,line width=1.8pt,no markers,smooth},
}
% Même style disponible dans un \draw TikZ simple (hors pgfplots)
\tikzset{popcurve/.style={poporange,line width=1.8pt}}
\tikzset{popgrid/.style={popline,very thin}}

% ============================================================
% Pill / squiggle
% ============================================================
\newcommand{\poppill}[3]{%
  \tikz[baseline=-0.55ex]\node[draw=popink,line width=1.2pt,rounded corners=2.6pt,%
    fill=#2,inner xsep=6.5pt,inner ysep=3pt]{%
    \popxboldls\fontsize{7}{7}\selectfont\color{#3}\MakeUppercase{#1}};%
}
\newcommand{\popsquiggle}[1]{%
  \tikz[baseline=-0.4ex]{
    \draw[#1,line width=2.6pt,line cap=round]
      plot [smooth] coordinates {(0,0.09) (0.34,0.01) (0.68,0.21) (1.02,0.01) (1.36,0.10)};
  }%
}
% Mot-clé surligné (marqueur orange sous le texte)
\newcommand{\cle}[1]{{\popxbold\color{popdark}#1}}

% ============================================================
% En-tête / pied
% ============================================================
\pagestyle{fancy}
\fancyhf{}
\renewcommand{\headrulewidth}{0pt}
\renewcommand{\footrulewidth}{0pt}
\lhead{\raisebox{-0.15ex}{\poplogo\fontsize{13}{13}\selectfont\color{popink}OnBuch\color{poporange}+}}
\rhead{\poppill{\DOCMATIERE\ \textbullet\ \DOCNIVEAU\ \textbullet\ Leçon \DOCLECON}{white}{popink}}
\lfoot{\popxbold\fontsize{7.6}{7.6}\selectfont\color{popmuted}\MakeUppercase{Cours signé OnBuch+}\ \textbullet\ {\popsemi\DOCTITRE}}
\rfoot{\tikz[baseline=-0.6ex]\node[rounded corners=8.5pt,fill=popink,minimum height=17pt,minimum width=17pt,inner xsep=5pt,inner sep=0]{\popxbold\fontsize{8}{8}\selectfont\color{popcream}\thepage\,/\,\pageref*{LastPage}};}

% ============================================================
% Titres
% ============================================================
\newcommand{\chaptitre}[2]{%
  \begin{tcolorbox}[enhanced,colback=poporange,colframe=popink,boxrule=2.6pt,arc=11pt,
      drop shadow=popink,left=15pt,right=15pt,top=12pt,bottom=11pt,before skip=3pt,after skip=18pt]
    \poppill{#2}{popink}{popcream}\par\vspace{9pt}
    {\raggedright\hyphenpenalty=10000\exhyphenpenalty=10000\popdisplay\fontsize{22}{24}\selectfont\color{popcream}\MakeUppercase{#1}\par}
  \end{tcolorbox}
}
% Section numérotée : pastille orange avec le numéro + titre Archivo
\newcounter{popsec}
\newcommand{\coursec}[1]{%
  \stepcounter{popsec}\setcounter{popsub}{0}%
  \par\vspace{16pt}\noindent
  \begin{minipage}{\linewidth}
  \tikz[baseline=-0.7ex]\node[draw=popink,line width=1.6pt,fill=poporange,rounded corners=4pt,minimum size=22pt,inner sep=0]{\popdisplay\fontsize{12}{12}\selectfont\color{popcream}\thepopsec};%
  \hspace{8pt}{\popdisplay\fontsize{15}{17}\selectfont\color{popink}\MakeUppercase{#1}}\par
  \vspace{4pt}\noindent{\color{poporange}\rule{36pt}{3.6pt}}
  \end{minipage}\par\nopagebreak\vspace{8pt}
}
\newcounter{popsub}
\newcommand{\courssub}[1]{%
  \stepcounter{popsub}%
  \par\vspace{9pt}\noindent{\popxbold\fontsize{11.5}{13}\selectfont\color{popink}{\color{poporange}\thepopsec.\thepopsub}\hspace{6pt}#1}\par\nopagebreak\vspace{4pt}
}

% ============================================================
% Boîtes pédagogiques — même recette : fond blanc, bordure épaisse colorée,
% ombre dure encre, badge plein. Titre optionnel [#1].
% ============================================================
\tcbset{popbase/.style={breakable,enhanced,colback=white,boxrule=2.3pt,arc=9pt,
  shadow={3pt}{-3pt}{0pt}{popink},left=14pt,right=14pt,top=11pt,bottom=11pt,before skip=11pt,after skip=15pt,
  fontupper=\popsemi\fontsize{10.6}{14.5}\selectfont\color{popink},
  fontlower=\popsemi\fontsize{10.6}{14.5}\selectfont\color{popink}}}
% Coche dessinée (le glyphe ✓ n'existe pas dans les polices du document)
\newcommand{\popcheck}[1]{\tikz[baseline=-0.5ex]\draw[#1,line width=1.6pt,line cap=round,line join=round](-0.13,0.0)--(-0.04,-0.09)--(0.13,0.1);}
\providecommand{\coche}{\popcheck{popgreen}}
\newcommand{\popboxhead}[3]{\poppill{#1}{#2}{white}\ifx\relax#3\relax\else\hspace{6pt}{\popxbold\fontsize{10.6}{12}\selectfont\color{popink}#3}\fi\par\vspace{7pt}}

\newtcolorbox{aretenir}[1][]{popbase,colframe=popgreen,before upper={\popboxhead{À retenir}{popgreen}{#1}}}
\newtcolorbox{exemplebox}[1][]{popbase,colframe=poporange,before upper={\popboxhead{Exemple}{poporange}{#1}}}
\newtcolorbox{definition}[1][]{popbase,colframe=popblue,before upper={\popboxhead{Définition}{popblue}{#1}}}
\newtcolorbox{propriete}[1][]{popbase,colframe=poppurple,before upper={\popboxhead{Propriété}{poppurple}{#1}}}
\newtcolorbox{methode}[1][]{popbase,colframe=popgold,before upper={\popboxhead{Méthode}{popgold}{#1}}}
\newtcolorbox{attention}[1][]{popbase,colframe=poppink,before upper={\popboxhead{Attention}{poppink}{#1}}}
\newtcolorbox{experience}[1][]{popbase,colframe=popblue,colback=popblueL,before upper={\popboxhead{Expérience}{popblue}{#1}}}
% « Le savais-tu ? » : carte violette pleine (contexte, culture, Cameroun)
\newtcolorbox{savaistu}[1][]{popbase,colback=poppurple,colframe=popink,
  fontupper=\popsemi\fontsize{10.4}{14.2}\selectfont\color{white},
  before upper={\poppill{Le savais-tu\,?}{white}{poppurple}\ifx\relax#1\relax\else\hspace{6pt}{\popxbold\color{white}#1}\fi\par\vspace{7pt}}}
% Exercice résolu : énoncé en haut, \tcblower puis la solution sur fond crème
\newcounter{popexo}\newcounter{popexr}
\newtcolorbox{exoresolu}[1][]{popbase,colframe=poporange,colbacklower=popcream,
  segmentation style={popink,line width=1.2pt,dashed},
  before upper={\stepcounter{popexr}\popboxhead{Exercice résolu \thepopexr}{poporange}{#1}},
  before lower={\poppill{Solution}{popink}{popcream}\par\vspace{6pt}}}
% Exercice (non résolu en place) — la correction est dans la partie Corrigés
\newtcolorbox{exercice}[1][]{popbase,colframe=popink,
  before upper={\stepcounter{popexo}\popboxhead{Exercice \thepopexo}{popink}{#1}}}
% Corrigé
\newtcolorbox{corrige}[1][]{popbase,colframe=popgreen,colback=popgreenL,
  before upper={\popboxhead{Corrigé}{popgreen}{#1}}}
% Objectifs / prérequis / bilan
\newtcolorbox{objectifs}{popbase,colframe=popink,colback=poporangeL,
  before upper={\popboxhead{Objectifs de la leçon}{popink}{}}}
\newtcolorbox{prerequis}{popbase,colframe=popink,colback=popblueL,
  before upper={\popboxhead{Prérequis}{popblue}{}}}
\newtcolorbox{bilan}{popbase,colframe=popink,colback=white,boxrule=2.8pt,
  before upper={\popboxhead{Fiche bilan}{poporange}{L'essentiel à connaître pour l'examen}}}
% Figure encadrée avec légende
\newtcolorbox{popfigure}[1][]{enhanced,breakable=false,colback=white,colframe=popline,boxrule=1.4pt,arc=8pt,
  left=10pt,right=10pt,top=10pt,bottom=8pt,before skip=10pt,after skip=12pt,halign=center,
  lower separated=false,halign lower=center,
  fontlower=\popsemi\fontsize{9}{11.5}\selectfont\color{popmuted},
  #1}
\newcommand{\legende}[1]{\par\vspace{6pt}{\popsemi\fontsize{9}{11.5}\selectfont\color{popmuted}#1}}

% ============================================================
% Signature OnBuch+
% ============================================================
\newcommand{\poplogoinline}[1]{{\poplogo\fontsize{#1}{#1}\selectfont\color{popink}OnBuch\color{poporange}+}}
\newcommand{\signatureonbuch}{%
  \begin{tcolorbox}[enhanced,colback=popink,colframe=popink,boxrule=2.2pt,arc=10pt,
      drop shadow=poporange,left=14pt,right=14pt,top=10pt,bottom=10pt,before skip=0pt,after skip=0pt]
    \tikz[baseline=-0.7ex]\node[circle,fill=popgreen,draw=popcream,line width=1.3pt,minimum size=22pt,inner sep=0]{\popcheck{white}};%
    \hspace{9pt}\begin{minipage}[c]{0.62\linewidth}
      {\popxbold\fontsize{12}{13}\selectfont\color{popcream}Signé\ \textbullet\ {\poplogo OnBuch\color{poporange}+}}\par\vspace{2pt}
      {\raggedright\popsemi\fontsize{8}{10.5}\selectfont\color{popline}\MakeUppercase{Équipe pédagogique OnBuch+}\\\MakeUppercase{Conforme au programme officiel MINESEC}\par}
    \end{minipage}\hfill
    {\poplogo\fontsize{22}{22}\selectfont\color{popcream}OnBuch\color{poporange}+}
  \end{tcolorbox}%
}

% ============================================================
% Page de garde
% #1 titre · #2 matière · #3 niveau · #4 module · #5 n° leçon · #6 durée ·
% #7 description / objectifs (texte ou itemize)
% ============================================================
\newcommand{\pagegarde}[7]{%
  \thispagestyle{empty}
  \newgeometry{a4paper,margin=1.7cm,top=1.6cm,bottom=1.4cm}%
  \begin{tikzpicture}[remember picture,overlay]
    \fill[poporange!18] ([xshift=-3.2cm,yshift=-3.6cm]current page.north east) circle (4.6cm);
    \fill[poppurple!14] ([xshift=2.4cm,yshift=7.6cm]current page.south west) circle (3.8cm);
  \end{tikzpicture}%
  {\poplogo\fontsize{30}{30}\selectfont\color{popink}OnBuch\color{poporange}+}\hfill
  \raisebox{0.6ex}{\poppill{Cours signé \textbullet\ Premium}{popink}{popcream}}\par
  \vspace{1pt}\popsquiggle{poppurple}\par
  \vspace{8pt}
  \begin{tcolorbox}[enhanced,colback=poporange,colframe=popink,boxrule=2.8pt,arc=13pt,
      drop shadow=popink,left=18pt,right=18pt,top=12pt,bottom=13pt,after skip=10pt]
    \poppill{#2 \textbullet\ #3}{popink}{popcream}\hspace{5pt}\poppill{#4}{white}{popink}\par\vspace{8pt}
    {\popdisplay\fontsize{14}{14}\selectfont\color{popink}LEÇON #5}\par\vspace{4pt}
    {\raggedright\hyphenpenalty=10000\exhyphenpenalty=10000\popdisplay\fontsize{25}{27}\selectfont\color{popcream}\MakeUppercase{#1}\par}
    \vspace{8pt}
    {\popxbold\fontsize{9.5}{9.5}\selectfont\color{popink}\MakeUppercase{Durée conseillée : #6}}
  \end{tcolorbox}
  \begin{tcolorbox}[enhanced,colback=white,colframe=popink,boxrule=2.2pt,arc=10pt,
      drop shadow=popink,left=16pt,right=16pt,top=10pt,bottom=10pt,after skip=8pt]
    \poppill{Dans cette leçon}{poppurple}{white}\par\vspace{5pt}
    \popsemi\fontsize{9.3}{12.6}\selectfont\color{popink}#7
  \end{tcolorbox}
  \noindent\begin{minipage}[t]{0.487\linewidth}
    \begin{tcolorbox}[enhanced,colback=popgreen,colframe=popink,boxrule=2.2pt,arc=10pt,
        drop shadow=popink,left=11pt,right=11pt,top=10pt,bottom=10pt,height=3.1cm,valign=center]
      \tcbox[colback=white,colframe=popink,boxrule=1.3pt,arc=5pt,boxsep=0pt,nobeforeafter,
        left=3pt,right=3pt,top=3pt,bottom=3pt,valign=center]{\qrcode[height=1.9cm]{https://play.google.com/store/apps/details?id=com.luvvix.onbuch&hl=fr}}%
      \hspace{8pt}\begin{minipage}[c]{0.5\linewidth}
        {\popdisplay\fontsize{11}{12.5}\selectfont\color{white}\MakeUppercase{Télécharge OnBuch}}\par\vspace{4pt}
        {\raggedright\popsemi\fontsize{8.4}{11}\selectfont\color{white}Gratuit \textbullet\ Google Play\\Tuteur IA, annales, concours, résultats\par}
      \end{minipage}
    \end{tcolorbox}
  \end{minipage}\hfill
  \begin{minipage}[t]{0.487\linewidth}
    \begin{tcolorbox}[enhanced,colback=poppurple,colframe=popink,boxrule=2.2pt,arc=10pt,
        drop shadow=popink,left=11pt,right=11pt,top=10pt,bottom=10pt,height=3.1cm,valign=center]
      \tikz[baseline=-0.7ex]\node[circle,fill=white,draw=popink,line width=1.3pt,minimum size=1.6cm,inner sep=0]{\popdisplay\fontsize{24}{24}\selectfont\color{poppurple}+};%
      \hspace{8pt}\begin{minipage}[c]{0.6\linewidth}
        {\popdisplay\fontsize{11}{12.5}\selectfont\color{white}\MakeUppercase{Passe à OnBuch+}}\par\vspace{4pt}
        {\raggedright\popsemi\fontsize{8.4}{11}\selectfont\color{white}Tous les cours signés, Tuteur IA illimité, fiches PDF — onglet OnBuch+ de l'app\par}
      \end{minipage}
    \end{tcolorbox}
  \end{minipage}
  \vspace{8pt}
  \popcontenu
  \vfill
  \signatureonbuch
  \restoregeometry
  \newpage
}

% Rangée « Ce cours contient » (page de garde)
\newcommand{\poptile}[3]{% #1 couleur #2 gros texte #3 légende
  \begin{tcolorbox}[enhanced,colback=#1,colframe=popink,boxrule=2pt,arc=9pt,shadow={2.5pt}{-2.5pt}{0pt}{popink},
      left=6pt,right=6pt,top=9pt,bottom=9pt,halign=center,valign=center,height=1.75cm,width=\linewidth,nobeforeafter]
    {\popdisplay\fontsize{12.5}{13}\selectfont\color{white}\MakeUppercase{#2}}\par\vspace{3pt}
    {\centering\popsemi\fontsize{7.8}{9.5}\selectfont\color{white}#3\par}
  \end{tcolorbox}}
\newcommand{\popcontenu}{%
  \par\noindent{\popxboldls\fontsize{7.5}{7.5}\selectfont\color{popmuted}CE COURS SIGNÉ CONTIENT}\par\vspace{7pt}
  \noindent\begin{minipage}[t]{0.235\linewidth}\poptile{popblue}{Cours}{complet, illustré, pas à pas}\end{minipage}\hfill
  \begin{minipage}[t]{0.235\linewidth}\poptile{poporange}{Méthodes}{et exercices résolus}\end{minipage}\hfill
  \begin{minipage}[t]{0.235\linewidth}\poptile{poppink}{Exercices}{type examen + corrigés}\end{minipage}\hfill
  \begin{minipage}[t]{0.235\linewidth}\poptile{popgreen}{Fiche bilan}{l'essentiel à retenir}\end{minipage}\par}

% Sommaire
\newtcolorbox{sommaire}{enhanced,colback=white,colframe=popink,boxrule=2.2pt,arc=10pt,
  drop shadow=popink,left=18pt,right=18pt,top=13pt,bottom=13pt,before skip=6pt,after skip=20pt,
  before upper={\poppill{Sommaire}{poppurple}{white}\par\vspace{9pt}\popsemi\fontsize{10.6}{14.5}\selectfont\color{popink}}}

% Signature de fin de cours — poussée tout en bas de la dernière page
\newcommand{\findecours}{%
  \par\vfill\nopagebreak
  \begin{center}\popsquiggle{poporange}\end{center}\vspace{4pt}
  \signatureonbuch
}

%% FIGFIX-BEGIN v5 — correctifs globaux des figures (docs/onbuch-generation/tools/figfix)
\usepackage{adjustbox}\usepackage{etoolbox}\tcbuselibrary{raster}
% toute figure TikZ/pgfplots plus large que la ligne est réduite à la largeur disponible
\BeforeBeginEnvironment{tikzpicture}{\begin{adjustbox}{max width=\linewidth}}
\AfterEndEnvironment{tikzpicture}{\end{adjustbox}}
% légendes pgfplots lisibles par-dessus les courbes
\pgfplotsset{every axis/.append style={legend style={fill=white,fill opacity=0.92,text opacity=1,draw=black!15,inner sep=3pt}}}
% étiquettes d'arêtes (syntaxe edge["..."]) avec fond blanc
\tikzset{every edge quotes/.append style={fill=white,inner sep=1.2pt}}
%% FIGFIX-END
\begin{document}

\pagegarde{Interaction électrostatique et champ électrique}{Physique}{Tle D}{Module 2 : Mouvements et interactions — évolutions temporelles des systèmes mécaniques}{P03}{4 h}{Cette leçon étudie l'interaction électrostatique : nature de la charge électrique, modes d'électrisation, loi de Coulomb et notion de champ électrostatique. Elle aboutit à l'étude du champ uniforme entre les armatures d'un condensateur plan, aux notions de potentiel et de travail de la force électrique, outils indispensables pour la suite du programme (mouvement dans un champ électrique).
\vspace{4pt}
\begin{multicols}{2}\small
\begin{itemize}
\item À la fin de cette leçon, je sais décrire et interpréter les trois modes d'électrisation (frottement, contact, influence)
\item À la fin de cette leçon, je sais énoncer la conservation et la quantification de la charge électrique
\item À la fin de cette leçon, je sais exprimer et calculer l'intensité d'une force électrostatique à l'aide de la loi de Coulomb
\item À la fin de cette leçon, je sais exprimer et représenter le champ électrostatique créé par une charge ponctuelle
\item À la fin de cette leçon, je sais appliquer le principe de superposition pour déterminer le champ créé par plusieurs charges
\item À la fin de cette leçon, je sais tracer et interpréter des lignes de champ
\end{itemize}\end{multicols}}

\begin{sommaire}
\begin{multicols}{2}
\begin{enumerate}
\item Charge électrique et électrisation
\item Loi de Coulomb
\item Champ électrostatique créé par des charges ponctuelles
\item Champ électrique uniforme et condensateur plan
\item Travail de la force électrique, potentiel et d.d.p.
\item Synthèse et méthodes de résolution
\item Activité d'intégration
\item Méthodes et exercices résolus
\item Exercices
\item Corrigés des exercices
\item Fiche bilan
\end{enumerate}
\end{multicols}
\end{sommaire}

\begin{objectifs}
\begin{itemize}
\item À la fin de cette leçon, je sais décrire et interpréter les trois modes d'électrisation (frottement, contact, influence)
\item À la fin de cette leçon, je sais énoncer la conservation et la quantification de la charge électrique
\item À la fin de cette leçon, je sais exprimer et calculer l'intensité d'une force électrostatique à l'aide de la loi de Coulomb
\item À la fin de cette leçon, je sais exprimer et représenter le champ électrostatique créé par une charge ponctuelle
\item À la fin de cette leçon, je sais appliquer le principe de superposition pour déterminer le champ créé par plusieurs charges
\item À la fin de cette leçon, je sais tracer et interpréter des lignes de champ
\item À la fin de cette leçon, je sais décrire un condensateur plan et relier le champ uniforme à la tension par E = U/d
\item À la fin de cette leçon, je sais définir le potentiel électrique et calculer le travail de la force électrique dans un champ uniforme
\end{itemize}
\end{objectifs}

\begin{prerequis}
\begin{itemize}
\item Notion de force et représentation vectorielle (Seconde/Première)
\item Loi de gravitation universelle et champ de gravitation (Première) — analogie utile
\item Structure de l'atome : noyau, électrons, ions (Seconde)
\item Tension électrique et intensité du courant (notions de collège et Première)
\item Produit scalaire et travail d'une force (Première)
\end{itemize}
\end{prerequis}

\newpage

\coursec{Charge électrique et électrisation}

En pleine saison sèche à Maroua, Amina remarque qu'en retirant son pagne synthétique, ses cheveux se dressent et crépitent, et que la poussière colle à l'écran de son téléphone. Plus impressionnant encore : pendant les orages de la saison des pluies à Yaoundé, la foudre frappe parfois les poteaux électriques de la SONATREL et coupe le courant dans tout le quartier. Comment expliquer que des objets, sans même se toucher, s'attirent ou se repoussent ? Quelle est cette force mystérieuse qui agit à distance, et comment la calculer ? Cette leçon lève le voile sur l'interaction électrostatique, phénomène à la base du fonctionnement des photocopieuses, des peintures électrostatiques et des paratonnerres.

\courssub{Expériences d'électrisation et mise en évidence de deux espèces d'électricité}

Commençons par des expériences simples, reproductibles à la maison ou en classe.

\begin{experience}[Électrisation par frottement -- règle et papier]
\textbf{Matériel :} une règle en plastique (ou un stylo à bille), de petits morceaux de papier (confettis), un ballon en caoutchouc, un mur peint.
\textbf{Protocole :}
\begin{enumerate}
    \item Frottte énergiquement la règle en plastique sur un tissu en laine (ou sur tes cheveux secs).
    \item Approche la règle des confettis sans les toucher.
    \item Frottte le ballon sur tes cheveux, puis colle-le au mur.
\end{enumerate}
\textbf{Observations :}
\begin{itemize}
    \item La règle attire les confettis ; ils sautent vers elle, s'y collent un instant, puis retombent.
    \item Le ballon reste collé au mur plusieurs minutes.
    \item Si tu approches deux règles frottées l'une de l'autre, elles se repoussent.
    \item Si tu approches une règle frottée (plastique) d'un bâton de verre frotté avec de la soie, ils s'attirent.
\end{itemize}
\textbf{Interprétation :} Le frottement a transféré des charges électriques. Il existe \textbf{deux espèces d'électricité} : on les appelle \emph{charge positive} et \emph{charge négative}. 
\begin{itemize}
    \item Deux charges de \textbf{même signe} se \textbf{repoussent}.
    \item Deux charges de \textbf{signes opposés} s'\textbf{attirent}.
\end{itemize}
\end{experience}

\begin{definition}[Charge électrique]
La \cle{charge électrique} est une propriété intrinsèque de la matière qui se manifeste par des forces d'attraction ou de répulsion entre corps électrisés. Elle est notée $q$ (ou $Q$) et son unité SI est le \cle{coulomb} (symbole \textbf{C}).
\end{definition}

\courssub{Interprétation microscopique : structure de l'atome et transfert d'électrons}

Pour comprendre l'origine de ces charges, regardons la structure de la matière.

\begin{popfigure}
\begin{tikzpicture}[scale=1.2, every node/.style={font=\small}]
% Noyau
\draw[popdark, thick, fill=poporange!30] (0,0) circle (0.6);
\node at (0,0) {\textbf{Noyau}};
\node at (0,-0.35) {\scriptsize protons $(+)$ + neutrons $(0)$};

% Électrons sur orbites
\foreach \angle/\r in {30/1.5, 150/1.5, 270/1.5, 90/2.2, 210/2.2, 330/2.2} {
    \draw[popblue, thick, fill=popblue!20] (\angle:\r) circle (0.15);
    \node at (\angle:\r) {\scriptsize $-$};
    \draw[popmuted, dashed] (0,0) circle (\r);
}

% Légende
\node[align=center] at (4.5,1.5) {\textbf{Atome neutre}\\$Z$ protons $(+e)$\\$Z$ électrons $(-e)$\\Charge totale $=0$};

% Flèche de transfert
\draw[poporange, very thick, ->] (4.5,0) -- (5.5,0) node[midway, above] {Frottement};

% Atome après perte d'électron
\begin{scope}[xshift=7cm]
\draw[popdark, thick, fill=poporange!30] (0,0) circle (0.6);
\node at (0,0) {\textbf{Noyau}};
\node at (0,-0.35) {\scriptsize protons $(+)$ + neutrons $(0)$};
\foreach \angle/\r in {150/1.5, 270/1.5, 90/2.2, 210/2.2, 330/2.2} {
    \draw[popblue, thick, fill=popblue!20] (\angle:\r) circle (0.15);
    \node at (\angle:\r) {\scriptsize $-$};
    \draw[popmuted, dashed] (0,0) circle (\r);
}
% Électron parti
\draw[popblue, thick, fill=popblue!20] (3.5,1.5) circle (0.15);
\node at (3.5,1.5) {\scriptsize $-$};
\draw[popblue, thick, ->] (30:1.5) -- (3.5,1.5);
\node[align=center] at (3.5,-1.5) {\textbf{Ion positif (cation)}\\$Z$ protons, $Z-1$ électrons\\Charge $= +e$};
\end{scope}

% Atome après gain d'électron
\begin{scope}[xshift=7cm, yshift=-4.5cm]
\draw[popdark, thick, fill=poporange!30] (0,0) circle (0.6);
\node at (0,0) {\textbf{Noyau}};
\node at (0,-0.35) {\scriptsize protons $(+)$ + neutrons $(0)$};
\foreach \angle/\r in {30/1.5, 150/1.5, 270/1.5, 90/2.2, 210/2.2, 330/2.2, 60/2.2} {
    \draw[popblue, thick, fill=popblue!20] (\angle:\r) circle (0.15);
    \node at (\angle:\r) {\scriptsize $-$};
    \draw[popmuted, dashed] (0,0) circle (\r);
}
\node[align=center] at (3.5,-1.5) {\textbf{Ion négatif (anion)}\\$Z$ protons, $Z+1$ électrons\\Charge $= -e$};
\end{scope}
\end{tikzpicture}
\legende{Modèle simplifié de l'atome : le frottement arrache des électrons périphériques (faiblement liés) à un atome pour les donner à un autre. Seuls les électrons se déplacent ; les protons restent dans le noyau.}
\end{popfigure}

\begin{propriete}[Structure de la matière et charge électrique]
La matière est constituée d'atomes. Chaque atome contient :
\begin{itemize}
    \item un \textbf{noyau} (protons de charge $+e$ et neutrons de charge $0$), massif et immobile à l'échelle macroscopique ;
    \item un \textbf{nuage électronique} (électrons de charge $-e$), léger et mobile.
\end{itemize}
Un atome \textbf{neutre} a autant de protons que d'électrons ($Z$ protons, $Z$ électrons). 
\begin{itemize}
    \item Un corps \textbf{chargé positivement} a un \textbf{déficit d'électrons} (il a \emph{perdu} des électrons).
    \item Un corps \textbf{chargé négativement} a un \textbf{excès d'électrons} (il a \emph{gagné} des électrons).
\end{itemize}
La charge élémentaire est $e = \SI{1,602 \times 10^{-19}}{C}$.
\end{propriete}

\begin{attention}[Piège classique : « Un corps chargé positivement a gagné des protons »]
\textbf{FAUX !} Les protons sont enfermés dans le noyau, liés par l'interaction forte. Seuls les \textbf{électrons} (et parfois des ions entiers dans les solutions) peuvent se déplacer d'un corps à l'autre lors d'une électrisation macroscopique. Un corps positif a \emph{perdu} des électrons, pas gagné des protons.
\end{attention}

\courssub{Conducteurs et isolants : la mobilité des charges}

Tous les matériaux ne réagissent pas de la même façon face aux charges.

\begin{definition}[Conducteurs et isolants]
\begin{itemize}
    \item Un \cle{conducteur} (métaux : cuivre, aluminium, fer ; corps humain ; solutions ioniques comme l'eau du robinet ou le sang) possède des \textbf{charges libres} (électrons de conduction dans les métaux, ions en solution) qui peuvent se déplacer \textbf{à l'intérieur du volume} du matériau.
    \item Un \cle{isolant} (verre, plastique, caoutchouc, bois sec, air sec, porcelaine) ne possède \textbf{pas de charges libres mobiles} à l'échelle macroscopique. Les charges y restent localisées là où on les a déposées.
\end{itemize}
\end{definition}

\begin{savaistu}[Le corps humain : un conducteur... parfois dangereux]
Ton corps est un conducteur (eau + sels minéraux $\approx$ solution ionique). C'est pourquoi tu peux servir de « fil » pour décharger un objet électrisé (le « choc » en touchant une portière de voiture par temps sec). C'est aussi pourquoi les électriciens d'ENEO portent des gants isolants et des bottes en caoutchouc : pour \emph{casser} la conduction et éviter l'électrocution. L'air sec est un isolant, mais s'il est trop humide ou si le champ électrique est trop fort (foudre), il s'ionise et devient conducteur.
\end{savaistu}

\courssub{Les trois modes d'électrisation}

\begin{popfigure}
\begin{tikzpicture}[scale=0.9, every node/.style={font=\small}]
% --- 1. Frottement ---
\begin{scope}
\node[align=center, popdark] at (0,3.5) {\textbf{1. Par frottement}};
% Baguette neutre + laine neutre
\draw[thick] (-2,2) rectangle (-0.5,3) node[midway, align=center] {Baguette\\plastique};
\draw[thick, fill=popgray!20] (0.5,2) rectangle (2,3) node[midway] {Laine};
\node at (0.75,1.5) {Avant : neutres};
% Flèche
\draw[poporange, thick, <->] (-0.5,2.5) -- (0.5,2.5) node[midway, above] {Frottement};
% Après
\draw[thick, fill=popblue!10] (-2,0) rectangle (-0.5,1) node[midway, align=center] {$-\!-\!-\!-\!-$\\Baguette $-$};
\draw[thick, fill=poporange!10] (0.5,0) rectangle (2,1) node[midway, align=center] {$+\!+\!+\!+$\\Laine $+$};
\node at (0.75,-0.5) {Après : charges opposées};
\end{scope}

% --- 2. Contact ---
\begin{scope}[xshift=7cm]
\node[align=center, popdark] at (0,3.5) {\textbf{2. Par contact}};
% Baguette chargée + sphère neutre
\draw[thick, fill=popblue!10] (-2,2) rectangle (-0.5,3) node[midway, align=center] {$-\!-\!-\!-\!-$\\Baguette $-$};
\draw[thick] (0.5,2) circle (0.75) node[align=center] {Sphère\\métallique\\neutre};
\draw[poporange, thick, ->] (-0.5,2.5) -- (0.5,2.5) node[midway, above] {Contact};
% Après
\draw[thick, fill=popblue!10] (-2,0) rectangle (-0.5,1) node[midway, align=center] {$-\!-\!-$\\Baguette $-$};
\draw[thick, fill=popblue!10] (0.5,0) circle (0.75) node[align=center] {$-\!-\!-\!-\!-$\\Sphère $-$};
\node at (0.75,-0.5) {Après : même signe};
\end{scope}

% --- 3. Influence ---
\begin{scope}[xshift=14cm]
\node[align=center, popdark] at (0,3.5) {\textbf{3. Par influence}};
% Baguette chargée + sphère neutre (sans contact)
\draw[thick, fill=popblue!10] (-2,2) rectangle (-0.5,3) node[midway, align=center] {$-\!-\!-\!-\!-$\\Baguette $-$};
\draw[thick] (0.5,2) circle (0.75) node[align=center] {Sphère\\métallique\\neutre};
\draw[poporange, thick, ->, dashed] (-0.5,2.5) -- (0.5,2.5) node[midway, above] {Rapprochement};
% Charges induites
\draw[poporange, thick, fill=poporange!20] (0.5,2.75) circle (0.2) node {$+$};
\draw[popblue, thick, fill=popblue!20] (0.5,1.25) circle (0.2) node {$-$};
\draw[popmuted, <->] (0.5,2.95) -- (0.5,1.45) node[midway, right] {Séparation};
% Après mise à la terre (fil vers le bas)
\draw[thick, popgreen] (0.5,1) -- (0.5,0) node[below] {Terre};
\draw[popblue, thick, ->] (0.5,1.25) -- (0.5,0.5) node[midway, right] {$e^-$ partent};
% Résultat final
\draw[thick, fill=poporange!10] (0.5,-1) circle (0.75) node[align=center] {$+\!+\!+\!+$\\Sphère $+$};
\node at (0.75,-2) {Après : signe opposé};
\end{scope}
\end{tikzpicture}
\legende{Les trois modes d'électrisation. (1) Frottement : transfert d'électrons entre deux isolants. (2) Contact : partage des charges entre un conducteur et un corps chargé. (3) Influence : déplacement des charges libres dans un conducteur sous l'action d'un champ extérieur, puis mise à la terre pour figer la charge.}
\end{popfigure}

\begin{methode}[Prévoir le sens du transfert d'électrons et le signe final]
Pour savoir quel corps devient positif/négatif lors d'un frottement :
\begin{enumerate}
    \item Consulte la \textbf{série triboélectrique} (liste des matériaux classés selon leur tendance à perdre ou gagner des électrons). Exemple simplifié : 
    \[
    \text{(+) Tendance à perdre $e^-$ :} \quad \text{Verre} > \text{Nylon} > \text{Laine} > \text{Soie} > \text{Papier} > \text{Coton} > \text{Bois} > \text{Ambre} > \text{Caoutchouc} > \text{Polystyrène} > \text{Plastique (PVC, Téflon)} \quad \text{(-) Tendance à gagner $e^-$}
    \]
    \item Le matériau \textbf{plus haut} dans la série (plus « positif ») \textbf{perd des électrons} $\rightarrow$ devient \textbf{chargé positivement}.
    \item Le matériau \textbf{plus bas} (plus « négatif ») \textbf{gagne des électrons} $\rightarrow$ devient \textbf{chargé négativement}.
    \item \textbf{Exemple :} Plastique (bas) frotté avec Laine (haut) $\rightarrow$ Plastique $-$ / Laine $+$.
\end{enumerate}
\end{methode}

\courssub{Conservation et quantification de la charge électrique}

\begin{propriete}[Conservation de la charge -- Loi admise, vérifiée expérimentalement]
Dans un \textbf{système isolé} (aucun échange de matière avec l'extérieur), la \textbf{charge électrique totale} (somme algébrique de toutes les charges) \textbf{se conserve} au cours de tout processus (frottement, contact, influence, réaction chimique, désintégration nucléaire).
\[
\sum_i q_i = \text{constante}
\]
\textbf{Exemple :} Frottement règle/laine. Avant : $q_{\text{règle}}=0$, $q_{\text{laine}}=0$, $Q_{\text{tot}}=0$. Après : $q_{\text{règle}}=-Q$, $q_{\text{laine}}=+Q$, $Q_{\text{tot}}=0$.
\end{propriete}

\begin{propriete}[Quantification de la charge -- Loi admise]
Toute charge électrique $q$ est un \textbf{multiple entier} de la \cle{charge élémentaire} $e$ :
\[
q = n \cdot e \qquad \text{avec } n \in \mathbb{Z} \text{ (entier relatif) et } e = \SI{1,602 \times 10^{-19}}{C}
\]
$n>0$ : excès d'électrons (charge négative) ; $n<0$ : déficit d'électrons (charge positive).
\textbf{Conséquence :} La charge est \emph{discrète} (quantifiée), pas continue. À l'échelle macroscopique ($q \sim \SI{e-6}{C}$), $n \sim 10^{13}$ : l'effet « grain » est invisible, on traite $q$ comme une variable continue.
\end{propriete}

\begin{exemplebox}[Ordre de grandeur : charges macroscopiques]
\begin{itemize}
    \item \textbf{Ballon frotté :} $q \approx \SI{-1e-8}{C}$ $\Rightarrow$ $n \approx 6 \times 10^{10}$ électrons en excès.
    \item \textbf{Éclair (foudre) :} $q \approx \SI{5}{C} \text{ à } \SI{20}{C}$ $\Rightarrow$ $n \approx 3 \times 10^{19} \text{ à } 1 \times 10^{20}$ électrons transférés en quelques millisecondes ! C'est l'équivalent du courant d'une centrale nucléaire pendant une fraction de seconde.
    \item \textbf{Condensateur d'alimentation (ex. chargeur téléphone) :} $q \approx \SI{1e-3}{C}$.
\end{itemize}
\end{exemplebox}

\begin{savaistu}[La foudre au Cameroun : un spectacle... et un danger]
Le Cameroun est l'un des pays les plus foudroyés au monde (région de la Sanaga, Bamiléké, Adamaoua). Un éclair transporte $\sim \SI{10}{C}$ de charge sous une tension de $\sim \SI{100}{MV}$ et un courant de $\sim \SI{30}{kA}$ en $\sim \SI{30}{\mu s}$. L'énergie libérée ($\sim \SI{1}{GJ}$) chauffe l'air à $\SI{30000}{K}$ (5 fois la surface du Soleil) $\rightarrow$ onde de choc = tonnerre. Les paratonnerres (pointe métallique reliée à la terre) protègent les bâtiments en offrant un chemin préférentiel aux charges : ils \emph{ne attirent pas} la foudre, ils \emph{la guident} sans dommage. Pense à débrancher ta télévision pendant les orages à Yaoundé ou Douala !
\end{savaistu}

\begin{exoresolu}[Calcul du nombre d'électrons transférés]
Énoncé : Une règle en plastique acquiert une charge $q = \SI{-3,2 \times 10^{-9}}{C}$ après frottement sur de la laine. Combien d'électrons a-t-elle gagnés ? Quelle est la charge de la laine ?
\vspace{0.5em}
\tcblower
\textbf{Solution :}
\begin{enumerate}
    \item \textbf{Données :} $q = \SI{-3,2 \times 10^{-9}}{C}$, $e = \SI{1,6 \times 10^{-19}}{C}$.
    \item \textbf{Nombre d'électrons gagnés :} $n = \dfrac{|q|}{e} = \dfrac{3,2 \times 10^{-9}}{1,6 \times 10^{-19}} = 2,0 \times 10^{10}$ électrons.
    \item \textbf{Conservation de la charge :} Système isolé $\{\text{règle, laine}\}$. $Q_{\text{initial}} = 0 = q_{\text{règle}} + q_{\text{laine}}$.
    \item $q_{\text{laine}} = -q_{\text{règle}} = \SI{+3,2 \times 10^{-9}}{C}$.
    \item La laine a \emph{perdu} $2,0 \times 10^{10}$ électrons (charge positive).
\end{enumerate}
\textbf{Réponse :} La règle a gagné $2,0 \times 10^{10}$ électrons ; la laine porte une charge de $\SI{+3,2 \times 10^{-9}}{C}$.
\end{exoresolu}

\begin{exercice}[Application : électroscope à feuilles]
Un électroscope à feuilles d'or (conducteur) est neutre. On approche (sans contact) une baguette de plastique chargée négativement.
\begin{enumerate}
    \item Décris le déplacement des charges libres dans l'électroscope. Dessine un schéma.
    \item Quel est le signe des charges sur les feuilles ? Pourquoi s'écartent-elles ?
    \item On met l'électroscope à la terre (contact avec le doigt) tout en maintenant la baguette proche, puis on retire le doigt, puis la baguette. Quel est le signe final de la charge de l'électroscope ? Justifie.
\end{enumerate}
\end{exercice}

\begin{corrige}[Exercice : électroscope à feuilles]
\begin{enumerate}
    \item \textbf{Influence :} Les électrons libres (charges négatives) du métal sont repoussés par la baguette $-$ $\rightarrow$ ils migrent vers le bas (feuilles). Le haut (support) se charge positivement (déficit d'électrons), le bas (feuilles) négativement (excès).
    \item Les feuilles portent des charges de \textbf{même signe (négatif)} $\rightarrow$ elles se \textbf{repoussent} (force de Coulomb) $\rightarrow$ écartement.
    \item \textbf{Mise à la terre :} Les électrons en excès sur les feuilles peuvent fuir vers la Terre (réservoir infini de charges). Le doigt fait contact $\rightarrow$ les électrons partent vers la Terre. On retire le doigt (coupe le chemin) \textbf{avant} de retirer la baguette. L'électroscope a maintenant un \textbf{déficit d'électrons} global $\rightarrow$ charge \textbf{positive}. En retirant la baguette, les charges positives se redistribuent uniformément. L'électroscope reste chargé positivement.
\end{enumerate}
\end{corrige}

\begin{aretenir}[Bilan : Charge électrique et électrisation]
\begin{itemize}
    \item \textbf{Charge électrique} $q$ (C) : propriété de la matière, deux signes ($+$, $-$). Règle : même signe $\rightarrow$ répulsion, signes contraires $\rightarrow$ attraction.
    \item \textbf{Microspocique :} Seuls les \textbf{électrons} ($-e$) se déplacent macroscopiquement. Protons ($+e$) bloqués dans le noyau.
    \item \textbf{Conducteur} : charges libres mobiles (métaux, corps humain, solutions). \textbf{Isolant} : charges bloquées (plastique, verre, air sec).
    \item \textbf{3 modes :} Frottement (transfert $e^-$ entre isolants), Contact (partage sur conducteurs), Influence (déplacement sans contact + mise à la terre pour figer).
    \item \textbf{Conservation :} $\sum q = \text{cste}$ (système isolé).
    \item \textbf{Quantification :} $q = n \cdot e$, $e = \SI{1,6 \times 10^{-19}}{C}$, $n \in \mathbb{Z}$.
\end{itemize}
\end{aretenir}

\coursec{Loi de Coulomb}

Après avoir identifié la charge électrique comme la propriété fondamentale à l'origine des phénomènes d'électrisation, nous devons maintenant quantifier l'interaction entre deux corps électrisés. Comment la force que deux charges s'exercent mutuellement dépend-elle de leurs valeurs et de leur distance ? C'est à cette question que Charles-Augustin de Coulomb a répondu en 1785 grâce à une expérience d'une grande ingéniosité : la balance de torsion.

\courssub{L'expérience historique : la balance de torsion de Coulomb}

\begin{experience}[Balance de torsion de Coulomb (1785)]
\textbf{Matériel :} Une aiguille légère (paille ou tige de métal) suspendue par un fil de soie ou un cheveu tendu (fil de torsion). Une petite sphère conductrice (souvent en étain ou en laiton) est fixée à une extrémité de l'aiguille. Une seconde sphère identique, montée sur un support isolant, peut être approchée de la première. On dispose d'un corps électrisé (par frottement) pour communiquer des charges aux sphères par contact.

\textbf{Protocole :}
\begin{enumerate}
    \item On électrise les deux sphères de manière identique (même signe de charge) par contact avec le même corps électrisé. Elles se repoussent.
    \item La répulsion fait tourner l'aiguille, ce qui tord le fil de suspension. Le fil exerce un couple de rappel proportionnel à l'angle de torsion $\theta$ : $\mathcal{M}_{\text{rappel}} = -k_{\text{fil}}\,\theta$.
    \item À l'équilibre, le couple électrostatique de répulsion équilibre le couple de rappel du fil. En mesurant l'angle d'équilibre $\theta_{\text{eq}}$, on en déduit la force de répulsion $F$ (connaissant la longueur de l'aiguille $L$ et la constante de torsion $k_{\text{fil}}$ : $F = k_{\text{fil}}\,\theta_{\text{eq}} / L$).
    \item On varie la distance $r$ entre les centres des sphères (en déplaçant la sphère mobile) et on note l'angle $\theta_{\text{eq}}$ correspondant.
    \item On renouvelle l'expérience en divisant la charge d'une sphère par 2 (mise en contact avec une sphère neutre identique), puis par 4, etc., pour étudier l'influence du produit des charges.
\end{enumerate}

\textbf{Observations et interprétations :}
\begin{itemize}
    \item Pour des charges fixées, la force $F$ est proportionnelle à $1/r^2$ : le tracé de $F$ en fonction de $1/r^2$ est une droite passant par l'origine.
    \item Pour une distance fixée, la force $F$ est proportionnelle au produit des charges $q_A q_B$.
    \item La force est attractive si les charges sont de signes contraires, répulsive si elles sont de même signe.
    \item La force suit la droite joignant les centres des sphères (force centrale).
\end{itemize}
\end{experience}

\begin{popfigure}
\begin{tikzpicture}[scale=1.2, >=Stealth]
% Fil de suspension
\draw[thick, popdark] (0,3) -- (0,2.5);
% Support supérieur
\draw[fill=popgray!30] (-0.5,3) rectangle (0.5,3.2);
% Aiguille
\draw[thick, popblue] (-2,2.5) -- (2,2.5);
% Sphère fixe (gauche) - charge q1
\draw[fill=poporange!80] (-2,2.5) circle (0.15);
\node[poporange!80!black] at (-2,2.1) {$q_A$};
% Sphère mobile (droite) - charge q2
\draw[fill=popgreen!70!black] (1.5,2.5) circle (0.15);
\node[popgreen!70!black] at (1.5,2.1) {$q_B$};
% Distance r
\draw[<->, popmuted] (-1.85,2.2) -- (1.35,2.2);
\node[popmuted] at (-0.25,1.9) {$r$};
% Forces - Cas répulsif (mêmes signes)
\draw[->, thick, popred] (-2,2.5) -- (-2.7,2.5) node[left] {$\vec{F}_{A/B}$};
\draw[->, thick, popred] (1.5,2.5) -- (2.2,2.5) node[right] {$\vec{F}_{B/A}$};
% Légende
\node[align=center, popdark] at (0,3.6) {\textbf{Balance de torsion -- Cas répulsif ($q_A q_B > 0$)}};
\end{tikzpicture}
\legende{Schéma de principe de la balance de torsion. Deux sphères de même signe de charge se repoussent, tordant le fil de suspension. L'angle de torsion mesure l'intensité de la force.}
\end{popfigure}

\courssub{Énoncé de la loi de Coulomb}

Les résultats expérimentaux de Coulomb, confirmés depuis avec une précision extrême, se résument dans la loi fondamentale de l'électrostatique.

\begin{propriete}[Loi de Coulomb -- Force entre deux charges ponctuelles]
Deux charges ponctuelles $q_A$ et $q_B$, placées dans le vide (ou l'air, à très bonne approximation), séparées par une distance $r$, s'exercent l'une sur l'autre des forces \textbf{opposées}, de \textbf{même direction} (la droite $(AB)$), de \textbf{même intensité} donnée par :
\[
F = k\,\frac{|q_A\,q_B|}{r^2}
\]
où $k = \num{9,0} \times 10^9\ \SI{}{N.m^2.C^{-2}}$ est la \textbf{constante de Coulomb} (dans le vide/air).
\end{propriete}

\begin{aretenir}[Conditions d'application et vocabulaire]
\begin{itemize}
    \item \textbf{Charges ponctuelles :} La formule est rigoureusement exacte pour des charges ponctuelles. Pour des sphères conductrices chargées, elle s'applique en prenant $r$ comme la distance entre leurs \textbf{centres}, à condition que $r$ soit grand devant leurs rayons (ou qu'elles soient suffisamment éloignées pour que la répartition de charge reste sphérique).
    \item \textbf{Valeur de $k$ :} $k = \num{9,0} \times 10^9\ \SI{}{N.m^2.C^{-2}}$. Au Bac, cette valeur est fournie ou admise.
    \item \textbf{Lien avec $\varepsilon_0$ (complément Bac) :} $k = \dfrac{1}{4\pi\varepsilon_0}$, où $\varepsilon_0 = \num{8,85} \times 10^{-12}\ \SI{}{F.m^{-1}}$ (permittivité diélectrique du vide). Cette forme $\frac{1}{4\pi\varepsilon_0}$ simplifie les équations de Maxwell et le théorème de Gauss.
    \item \textbf{Unités :} Charges en \textbf{coulombs (C)}, distance en \textbf{mètres (m)}, force en \textbf{newtons (N)}. \textbf{Attention :} $1\ \mu C = 10^{-6}\ C$ ; $1\ cm = 10^{-2}\ m$.
\end{itemize}
\end{aretenir}

\courssub{Expression vectorielle et caractère attractif/répulsif}

Pour manipuler les forces dans l'espace (somme vectorielle, projection sur axes), nous avons besoin de l'expression vectorielle. Introduisons le vecteur unitaire $\vec{u}_{AB} = \dfrac{\vec{AB}}{AB}$ dirigé de $A$ vers $B$.

\begin{propriete}[Force de Coulomb -- Forme vectorielle]
La force $\vec{F}_{B/A}$ exercée par la charge $q_A$ (en $A$) sur la charge $q_B$ (en $B$) s'écrit :
\[
\vec{F}_{B/A} = k\,\frac{q_A\,q_B}{r^2}\,\vec{u}_{AB}
\]
La force $\vec{F}_{A/B}$ exercée par $q_B$ sur $q_A$ est :
\[
\vec{F}_{A/B} = k\,\frac{q_A\,q_B}{r^2}\,\vec{u}_{BA} = -\vec{F}_{B/A}
\]
\end{propriete}

\begin{aretenir}[Règle des signes -- Attraction ou Répulsion ?]
Le signe du produit $q_A q_B$ détermine la nature de l'interaction :
\begin{itemize}
    \item \textbf{$q_A q_B > 0$ (mêmes signes) :} $\vec{F}_{B/A}$ a le même sens que $\vec{u}_{AB}$ (de $A$ vers $B$). Les charges se \textbf{repoussent}.
    \item \textbf{$q_A q_B < 0$ (signes contraires) :} $\vec{F}_{B/A}$ a le sens opposé à $\vec{u}_{AB}$ (de $B$ vers $A$). Les charges s'\textbf{attirent}.
\end{itemize}
C'est l'avantage majeur de la forme vectorielle : \textbf{pas de valeur absolue}, le signe algébrique des charges donne directement le sens de la force.
\end{aretenir}

\begin{popfigure}
\begin{tikzpicture}[scale=1.1, >=Stealth]
% Cas Répulsif (haut)
\begin{scope}[yshift=4cm]
\node[align=center, popdark] at (0,1.5) {\textbf{Répulsion ($q_A q_B > 0$)}};
\draw[fill=poporange!80] (-1.5,0) circle (0.2);
\node[poporange!80!black] at (-1.5,-0.5) {$q_A > 0$};
\draw[fill=poporange!80] (1.5,0) circle (0.2);
\node[poporange!80!black] at (1.5,-0.5) {$q_B > 0$};
\draw[->, thick, popred] (-1.5,0) -- (-2.3,0) node[left] {$\vec{F}_{A/B}$};
\draw[->, thick, popred] (1.5,0) -- (2.3,0) node[right] {$\vec{F}_{B/A}$};
\draw[<->, popmuted] (-1.3,0.4) -- (1.3,0.4) node[midway, above] {$r$};
\draw[dashed, popmuted] (-1.5,0) -- (1.5,0);
\node at (0,0.6) {$\vec{u}_{AB}$};
\draw[->, popmuted] (0,0.6) -- (0.5,0.6);
\end{scope}
% Cas Attractif (bas)
\begin{scope}
\node[align=center, popdark] at (0,1.5) {\textbf{Attraction ($q_A q_B < 0$)}};
\draw[fill=popblue!80] (-1.5,0) circle (0.2);
\node[popblue!80!black] at (-1.5,-0.5) {$q_A > 0$};
\draw[fill=popred!80] (1.5,0) circle (0.2);
\node[popred!80!black] at (1.5,-0.5) {$q_B < 0$};
\draw[->, thick, popgreen!70!black] (-1.5,0) -- (-0.7,0);
\draw[->, thick, popgreen!70!black] (1.5,0) -- (0.7,0);
\node[popgreen!70!black] at (-1.1,0.3) {$\vec{F}_{A/B}$};
\node[popgreen!70!black] at (1.1,0.3) {$\vec{F}_{B/A}$};
\draw[<->, popmuted] (-1.3,-0.4) -- (1.3,-0.4) node[midway, below] {$r$};
\draw[dashed, popmuted] (-1.5,0) -- (1.5,0);
\node at (0,-0.6) {$\vec{u}_{AB}$};
\draw[->, popmuted] (0,-0.6) -- (0.5,-0.6);
\end{scope}
\end{tikzpicture}
\legende{Forces de Coulomb : répulsion pour charges de même signe, attraction pour charges de signes contraires. Les forces sont opposées (3\textsuperscript{e} loi de Newton).}
\end{popfigure}

\courssub{Principe des actions réciproques (3\textsuperscript{e} loi de Newton)}

La loi de Coulomb satisfait naturellement le principe des actions réciproques :
\[
\vec{F}_{A/B} = -\vec{F}_{B/A}
\]
Les deux forces ont même intensité, même direction, sens opposés. Elles s'exercent sur \textbf{deux corps différents} ($A$ et $B$). C'est une interaction \textbf{à distance} (sans contact), médiatisée par le champ électrostatique (notion développée en Section 3).

\courssub{Comparaison avec l'interaction gravitationnelle}

L'analogie formelle entre la loi de Coulomb et la loi de Newton sur la gravitation universelle est frappante.

\begin{center}
\begin{tabularx}{\linewidth}{|l|X|X|}\hline
\textbf{Propriété} & \textbf{Interaction Gravitationnelle} & \textbf{Interaction Électrostatique} \\ \hline
Expression scalaire & $F_g = G\,\dfrac{m_A m_B}{r^2}$ & $F_e = k\,\dfrac{|q_A q_B|}{r^2}$ \\ \hline
Constante & $G = \num{6,67} \times 10^{-11}\ \SI{}{N.m^2.kg^{-2}}$ & $k = \num{9,0} \times 10^9\ \SI{}{N.m^2.C^{-2}}$ \\ \hline
Source & Masse $m$ (toujours positive) & Charge $q$ (positive ou négative) \\ \hline
Nature & \textbf{Toujours attractive} & Attractive ou \textbf{répulsive} \\ \hline
Écran & Impossible & Possible (cage de Faraday, conducteurs) \\ \hline
Ordre de grandeur (électron-proton) & $F_g \approx 10^{-47}\ \SI{}{N}$ & $F_e \approx 10^{-8}\ \SI{}{N}$ \\ \hline
\end{tabularx}
\end{center}

\begin{savaistu}[La force électrique domine le monde microscopique]
Pour un couple électron-proton (atome d'hydrogène), le rapport des forces est :
\[
\frac{F_e}{F_g} = \frac{k\,e^2}{G\,m_e m_p} \approx \frac{(\num{9}\times10^9)(\num{1,6}\times10^{-19})^2}{(\num{6,67}\times10^{-11})(\num{9,11}\times10^{-31})(\num{1,67}\times10^{-27})} \approx \mathbf{2,3 \times 10^{39}}
\]
La force électrique est \textbf{39 ordres de grandeur} plus intense que la gravité ! C'est pourquoi la structure de la matière (atomes, molécules, solides, chimie, biologie) est gouvernée par l'électromagnétisme. La gravité ne l'emporte qu'à l'échelle macroscopique (planètes, étoiles) car la matière est globalement neutre (compensation des charges), alors que les masses s'additionnent toujours.
\end{savaistu}

\courssub{Applications numériques et méthode de résolution}

\begin{methode}[Calculer une force électrostatique entre deux charges ponctuelles]
\begin{enumerate}
    \item \textbf{Schéma :} Représenter les charges $A$ et $B$, leurs signes, la distance $r$ (entre centres), le vecteur unitaire $\vec{u}_{AB}$.
    \item \textbf{Conversion SI :} Convertir \textbf{impérativement} les charges en \textbf{coulombs (C)} et la distance en \textbf{mètres (m)}. ($1\ \mu C = 10^{-6}\ C$, $1\ cm = 10^{-2}\ m$).
    \item \textbf{Calcul de l'intensité :} Appliquer $F = k\,\dfrac{|q_A q_B|}{r^2}$.
    \item \textbf{Sens de la force :} Utiliser la forme vectorielle $\vec{F}_{B/A} = k\,\dfrac{q_A q_B}{r^2}\,\vec{u}_{AB}$ (le signe du produit donne le sens par rapport à $\vec{u}_{AB}$) ou la règle « même signe = répulsion, signes contraires = attraction ».
    \item \textbf{Vérification :} Analyse dimensionnelle ($[k] = \SI{}{N.m^2.C^{-2}}$, $[q^2/r^2] = \SI{}{C^2.m^{-2}}$, produit = $\SI{}{N}$). Ordre de grandeur cohérent.
\end{enumerate}
\end{methode}

\begin{exemplebox}[Force entre deux charges ponctuelles]
\textbf{Énoncé :} Deux charges ponctuelles $q_1 = +\num{2,0}\ \mu C$ et $q_2 = -\num{3,0}\ \mu C$ sont placées dans l'air à une distance $r = \num{10,0}\ cm$ l'une de l'autre. Déterminer la force que $q_1$ exerce sur $q_2$ (valeur et direction).

\textbf{Solution :}
\begin{enumerate}
    \item \textbf{Données SI :}
    $q_1 = +\num{2,0} \times 10^{-6}\ C$ ; $q_2 = -\num{3,0} \times 10^{-6}\ C$ ; $r = \num{0,100}\ m$.
    \item \textbf{Intensité :}
    \[
    F = k\,\frac{|q_1 q_2|}{r^2} = \num{9,0} \times 10^9 \times \frac{|(\num{2,0} \times 10^{-6})(\num{-3,0} \times 10^{-6})|}{(\num{0,100})^2}
    \]
    \[
    F = \num{9,0} \times 10^9 \times \frac{\num{6,0} \times 10^{-12}}{\num{1,0} \times 10^{-2}} = \num{9,0} \times 10^9 \times \num{6,0} \times 10^{-10} = \mathbf{5,4\ N}
    \]
    \item \textbf{Sens :} $q_1 q_2 < 0$ (signes contraires) $\Rightarrow$ \textbf{Attraction}.
    La force $\vec{F}_{2/1}$ (exercée par $q_1$ sur $q_2$) est dirigée \textbf{de $q_2$ vers $q_1$}.
    \item \textbf{Expression vectorielle :} En choisissant $\vec{u}_{12}$ dirigé de $q_1$ vers $q_2$ :
    \[
    \vec{F}_{2/1} = k\,\frac{q_1 q_2}{r^2}\,\vec{u}_{12} = \num{9,0} \times 10^9 \times \frac{-\num{6,0} \times 10^{-12}}{\num{0,01}}\,\vec{u}_{12} = -\num{5,4}\,\vec{u}_{12}\ \SI{}{N}
    \]
    Le signe $-$ confirme que la force sur $q_2$ est opposée à $\vec{u}_{12}$, c'est-à-dire vers $q_1$.
\end{enumerate}
\textbf{Résultat :} $\vec{F}_{2/1} = \num{5,4}\ \SI{}{N}$, dirigée de $q_2$ vers $q_1$ (attractive).
\end{exemplebox}

\begin{attention}[Pièges classiques à éviter au Bac]
\begin{itemize}
    \item \textbf{Oublier les conversions :} Utiliser $\mu C$ et $cm$ directement dans la formule donne un résultat faux par un facteur $10^{10}$ ou $10^4$. \textbf{Toujours} convertir en C et m \textbf{avant} de calculer.
    \item \textbf{Oublier le carré sur $r$ :} $F \propto 1/r^2$, pas $1/r$. Si $r$ double, $F$ est divisée par 4.
    \item \textbf{Confondre $\vec{u}_{AB}$ et $\vec{u}_{BA}$ :} $\vec{u}_{AB} = -\vec{u}_{BA}$. Bien définir son repère avant d'écrire la forme vectorielle.
    \item \textbf{Valeur absolue mal placée :} Dans la forme scalaire $F = k|q_A q_B|/r^2$, la valeur absolue assure une intensité positive. Dans la forme vectorielle, \textbf{pas de valeur absolue} : le signe de $q_A q_B$ porte l'information du sens.
    \item \textbf{Charges non ponctuelles :} Si les sphères sont proches, la répartition de charge n'est plus uniforme, la formule $k q_A q_B / r^2$ (avec $r$ distance entre centres) devient approximative. Au Bac, on suppose toujours l'hypothèse « charges ponctuelles » ou « sphères éloignées » validée.
\end{itemize}
\end{attention}

\courssub{Représentations graphiques de la loi de Coulomb}

La dépendance en $1/r^2$ se visualise utilement sous deux formes graphiques.

\begin{popfigure}
\begin{tikzpicture}
\begin{axis}[
    width=\linewidth, height=7cm,
    axis lines=middle,
    xlabel={$1/r^2\ (\SI{}{m^{-2}})$}, ylabel={$F\ (\SI{}{N})$},
    xmin=0, xmax=110, ymin=0, ymax=5.5,
    xtick={0,25,50,75,100}, ytick={0,1,2,3,4,5},
    grid=major, grid style={popline},
    title style={yshift=0.5cm},
    title={\textbf{Variation de $F$ en fonction de $1/r^2$ (charges fixes)}}
]
% Données pour q1=2µC, q2=-3µC => k|q1q2| = 5.4e-2 N.m^2
% F = 0.054 * (1/r^2)
\addplot[popcurve, domain=0:100, samples=50] {0.054*x};
\addplot[only marks, mark=*, popred] coordinates {(25, 1.35) (100, 5.4)};
\node[popred, anchor=south west] at (axis cs:100,5.4) {$r=0,10\ m$};
\node[popred, anchor=north east] at (axis cs:25,1.35) {$r=0,20\ m$};
\end{axis}
\end{tikzpicture}
\legende{Pour des charges fixées, $F$ est proportionnelle à $1/r^2$ : droite passant par l'origine. C'est le tracé linéarisé utilisé par Coulomb pour vérifier la loi.}
\end{popfigure}

\begin{popfigure}
\begin{tikzpicture}
\begin{axis}[
    width=\linewidth, height=7cm,
    axis lines=middle,
    xlabel={$r\ (\SI{}{m})$}, ylabel={$F\ (\SI{}{N})$},
    xmin=0, xmax=0.5, ymin=0, ymax=6,
    xtick={0,0.1,0.2,0.3,0.4,0.5}, ytick={0,1,2,3,4,5},
    grid=major, grid style={popline},
    title style={yshift=0.5cm},
    title={\textbf{Variation de $F$ en fonction de $r$ (charges fixes)}}
]
% F = 0.054 / r^2
\addplot[popcurve, domain=0.05:0.5, samples=200] {0.054/x^2};
\addplot[only marks, mark=*, popblue] coordinates {(0.1, 5.4) (0.2, 1.35) (0.3, 0.6)};
\node[popblue, anchor=south] at (axis cs:0.1,5.4) {$5,4\ N$};
\node[popblue, anchor=north] at (axis cs:0.2,1.35) {$1,35\ N$};
\end{axis}
\end{tikzpicture}
\legende{Courbe $F(r) = k|q_1 q_2|/r^2$ : décroissance hyperbolique rapide. La force devient négligeable dès que la distance augmente modérément.}
\end{popfigure}

\courssub{Exercice résolu : Détermination d'une charge inconnue}

\begin{exoresolu}[Détermination d'une charge par mesure de force]
\textbf{Énoncé :} Une petite sphère conductrice $A$, de charge inconnue $q_A$, est fixée sur un support isolant. On approche une seconde sphère $B$, de charge connue $q_B = +\num{5,0}\ \mu C$, sur la droite de $A$. On mesure la force de répulsion exercée sur $B$ pour différentes distances $r$ entre les centres. Les résultats sont consignés dans le tableau ci-dessous.

\begin{center}
\begin{tabular}{|c|c|c|c|c|c|}\hline
$r\ (\SI{}{cm})$ & 2,0 & 3,0 & 4,0 & 5,0 & 6,0 \\ \hline
$F\ (\SI{}{mN})$ & 56,3 & 25,0 & 14,1 & 9,0 & 6,25 \\ \hline
\end{tabular}
\end{center}

\begin{enumerate}
    \item Vérifier que ces résultats sont compatibles avec la loi de Coulomb.
    \item Déterminer la valeur de la charge $q_A$ (valeur et signe).
\end{enumerate}

\tcblower

\textbf{Solution détaillée :}

\begin{enumerate}
    \item \textbf{Vérification de la loi en $1/r^2$ :}
    On calcule $1/r^2$ (en $\SI{}{m^{-2}}$) pour chaque point.
    \begin{center}
    \begin{tabular}{|c|c|c|c|c|c|}\hline
    $r\ (\SI{}{m})$ & 0,020 & 0,030 & 0,040 & 0,050 & 0,060 \\ \hline
    $1/r^2\ (\times 10^3 \SI{}{m^{-2}})$ & 2,50 & 1,11 & 0,625 & 0,400 & 0,278 \\ \hline
    $F\ (\SI{}{N})$ & $\num{56,3}\times10^{-3}$ & $\num{25,0}\times10^{-3}$ & $\num{14,1}\times10^{-3}$ & $\num{9,0}\times10^{-3}$ & $\num{6,25}\times10^{-3}$ \\ \hline
    \end{tabular}
    \end{center}
    On trace $F$ en fonction de $1/r^2$. Les points s'alignent sur une droite passant par l'origine (proportionnalité). La pente $p$ de cette droite vaut $k\,|q_A q_B|$.
    Calcul de la pente par régression ou prise de deux points extrêmes :
    \[
    p \approx \frac{\num{56,3} \times 10^{-3} - 0}{\num{2,50} \times 10^3 - 0} = \num{22,5} \times 10^{-6}\ \SI{}{N.m^2}
    \]
    (Vérification avec le point $r=6\ cm$ : $p = \num{6,25}\times10^{-3} / \num{0,278}\times10^3 = \num{22,5}\times10^{-6}$). La proportionnalité est parfaite.

    \item \textbf{Détermination de $q_A$ :}
    On a $p = k\,|q_A q_B|$.
    \[
    |q_A| = \frac{p}{k\,|q_B|} = \frac{\num{22,5} \times 10^{-6}}{\num{9,0} \times 10^9 \times \num{5,0} \times 10^{-6}} = \frac{\num{22,5} \times 10^{-6}}{\num{45} \times 10^3} = \num{0,5} \times 10^{-9}\ C = \mathbf{0,50\ nC}
    \]
    \textbf{Signe de $q_A$ :} La force est une \textbf{répulsion} (données de l'énoncé : « force de répulsion »). Or $q_B > 0$. Pour que l'interaction soit répulsive, il faut $q_A q_B > 0$, donc $q_A > 0$.
    \textbf{Résultat final :} $q_A = \mathbf{+\num{0,50}\ nC}$ (ou $+\num{5,0} \times 10^{-10}\ C$).
\end{enumerate}
\end{exoresolu}

\begin{exercice}[Application : Force entre gouttelettes chargées]
Dans un nuage d'orage au-dessus de Yaoundé, deux gouttelettes d'eau sphériques, de rayon $R = \num{1,0}\ mm$, portent chacune une charge $q = -\num{10}\ nC$. Elles sont séparées par une distance $d = \num{2,0}\ cm$ (entre leurs centres).
\begin{enumerate}
    \item Calculer l'intensité de la force électrostatique de répulsion entre ces deux gouttelettes.
    \item Comparer cette force au poids d'une gouttelettes ($\rho_{\text{eau}} = \num{1000}\ \SI{}{kg.m^{-3}}$, $g = \num{9,8}\ \SI{}{N.kg^{-1}}$).
    \item Commenter l'influence de cette force électrique sur la coalescence des gouttelettes (formation de la pluie).
\end{enumerate}
\end{exercice}

\begin{corrige}[Exercice 1]
\begin{enumerate}
    \item \textbf{Force de Coulomb :}
    $q = -\num{10}\ nC = -\num{1,0} \times 10^{-8}\ C$ ; $d = \num{0,020}\ m$.
    \[
    F_e = k\,\frac{q^2}{d^2} = \num{9,0} \times 10^9 \times \frac{(\num{1,0} \times 10^{-8})^2}{(\num{2,0} \times 10^{-2})^2} = \num{9,0} \times 10^9 \times \frac{\num{1,0} \times 10^{-16}}{\num{4,0} \times 10^{-4}} = \num{2,25} \times 10^{-3}\ \SI{}{N}
    \]
    \textbf{$F_e = \num{2,25}\ mN$ (répulsion).}
    \item \textbf{Poids d'une gouttelette :}
    Volume $V = \frac{4}{3}\pi R^3 = \frac{4}{3}\pi (10^{-3})^3 \approx \num{4,19} \times 10^{-9}\ \SI{}{m^3}$.
    Masse $m = \rho V = \num{1000} \times \num{4,19} \times 10^{-9} = \num{4,19} \times 10^{-6}\ \SI{}{kg}$.
    Poids $P = mg = \num{4,19} \times 10^{-6} \times \num{9,8} \approx \num{4,1} \times 10^{-5}\ \SI{}{N} = \num{0,041}\ mN$.
    \item \textbf{Comparaison et commentaire :}
    $F_e / P \approx \num{2,25} / \num{0,041} \approx 55$. La force de répulsion électrique est \textbf{55 fois plus grande} que le poids de la gouttelette.
    Cela crée une barrière électrique qui s'oppose au rapprochement des gouttelettes (coalescence). Pour que la pluie se forme, il faut que les charges soient neutralisées (par capture d'ions de signe contraire, ou par décharge - éclair) ou que les gouttelettes grossissent suffisamment pour que la gravité l'emporte. C'est un facteur clé de la microphysique des nuages.
\end{enumerate}
\end{corrige}

\begin{aretenir}[Bilan : Ce qu'il faut maîtriser sur la loi de Coulomb]
\begin{itemize}
    \item \textbf{Formule scalaire :} $F = k\,\dfrac{|q_A q_B|}{r^2}$ avec $k = \num{9,0} \times 10^9\ \SI{}{N.m^2.C^{-2}}$.
    \item \textbf{Formule vectorielle :} $\vec{F}_{B/A} = k\,\dfrac{q_A q_B}{r^2}\,\vec{u}_{AB}$.
    \item \textbf{Unités SI obligatoires :} $q$ en $C$, $r$ en $m$, $F$ en $N$.
    \item \textbf{Sens :} $q_A q_B > 0 \Rightarrow$ répulsion ; $q_A q_B < 0 \Rightarrow$ attraction.
    \item \textbf{Action-réaction :} $\vec{F}_{A/B} = -\vec{F}_{B/A}$.
    \item \textbf{Graphique :} $F \propto 1/r^2$ (droite passant par l'origine en coordonnées $F$ vs $1/r^2$).
    \item \textbf{Comparaison gravitation :} Même forme $1/r^2$, mais charges signées vs masses positives, et force électrique $\sim 10^{39}$ fois plus forte à l'échelle particulaire.
\end{itemize}
\end{aretenir}

\coursec{Champ électrostatique créé par des charges ponctuelles}

\courssub{De la force à l'action à distance : la notion de champ}

Dans la section précédente, nous avons établi la \textbf{loi de Coulomb} : deux charges ponctuelles $q_1$ et $q_2$ séparées par une distance $r$ s'exercent une force attractive ou répulsive d'intensité $F = k\,\dfrac{|q_1 q_2|}{r^2}$. Cette loi décrit une \emph{action à distance} : la charge $q_1$ agit sur $q_2$ sans contact matériel, instantanément (dans le cadre statique). Pour rendre compte de cette action sans recourir à une « magie » instantanée, le physicien Michael Faraday a introduit au XIX\textsuperscript{e} siècle le concept révolutionnaire de \textbf{champ}.

L'idée est simple et puissante : \emph{une charge source $Q$ modifie les propriétés de l'espace qui l'entoure}. Cet espace modifié est le \textbf{champ électrostatique}. Si l'on place une autre charge $q$ (appelée \textbf{charge test}) en un point $M$ de cet espace, elle « sent » l'influence de $Q$ et subit une force $\vec{F}$. Le champ est la propriété de l'espace \emph{en l'absence de la charge test} ; la force est l'effet du champ \emph{sur} la charge test.

\begin{experience}[Visualiser le champ : l'expérience des limailles de « semoule » ou de graines de ricin]
\textbf{Matériel :} Générateur de haute tension (type Wimshurst ou alimentation 5 kV), deux électrodes (billes ou plaques), cuve en plexiglas, huile de paraffine (ou huile de table), graines de ricin (ou semoule fine, ou limaille de fer pour le magnétostatique).
\textbf{Protocole :}
\begin{enumerate}
    \item Placer les électrodes dans la cuve, les relier au générateur.
    \item Verser l'huile et saupoudrer généreusement les graines de ricin à la surface.
    \item Mettre sous tension (prudence : haute tension ! Ne pas toucher la cuve).
    \item Observer l'orientation des graines.
\end{enumerate}
\textbf{Observations :} Les graines s'alignent le long de courbes bien définies : radiales pour une bille chargée, courbes hyperboliques pour deux charges de signes opposés (dipôle), courbes fuyant l'axe pour deux charges de même signe.
\textbf{Interprétation :} Les graines se polarisent par influence et deviennent de minuscules dipôles électriques. Le couple électrique les oriente \textbf{tangentiellement aux lignes de champ}, dans le sens du vecteur champ $\vec{E}$. Cette expérience rend visible l'invisible : la structure du champ électrostatique.
\end{experience}

\courssub{Définition du vecteur champ électrostatique $\vec{E}$}

Soit une charge source $Q$ fixe. En un point $M$ de l'espace, on place une charge test $q$ (assez petite pour ne pas perturber la distribution des charges sources). La charge $q$ subit une force électrostatique $\vec{F}$. Le \textbf{vecteur champ électrostatique} $\vec{E}(M)$ en $M$ est défini par le rapport :

\begin{definition}[Vecteur champ électrostatique]
Le champ électrostatique en un point $M$ est le vecteur $\vec{E}(M)$ tel que la force électrostatique subie par une charge test $q$ placée en $M$ s'écrive :
\[
\vec{F} = q\,\vec{E}(M)
\]
\begin{itemize}
    \item \textbf{Unité SI :} \SI{1}{N/C} (newton par coulomb), équivalent au \SI{1}{V/m} (volt par mètre).
    \item \textbf{Direction et sens :} $\vec{E}(M)$ a la direction de la force $\vec{F}$ subie par une charge test \textbf{positive} ($q>0$). Si $q<0$, la force est opposée à $\vec{E}$.
    \item \textbf{Indépendance :} $\vec{E}(M)$ ne dépend \textbf{que} des charges sources et de la géométrie, pas de la charge test $q$.
\end{itemize}
\end{definition}

\begin{attention}[Piège classique : le sens du champ]
Ne confonds jamais le sens de la force $\vec{F}$ sur \emph{ta} charge $q$ et le sens du champ $\vec{E}$.
\begin{itemize}
    \item $q > 0$ : $\vec{F}$ et $\vec{E}$ ont \textbf{le même sens}.
    \item $q < 0$ : $\vec{F}$ et $\vec{E}$ ont \textbf{sens opposés}.
\end{itemize}
Par convention, le champ est défini par l'effet sur une charge \textbf{positive}. Une charge positive est \emph{repoussée} par une source positive ($Q>0$) et \emph{attirée} par une source négative ($Q<0$).
\end{attention}

\courssub{Champ créé par une charge ponctuelle}

Considérons une charge ponctuelle $Q$ fixe en $O$. Quelle est la valeur de $\vec{E}(M)$ en un point $M$ tel que $OM = r$ ?
On place une charge test $q$ en $M$. D'après la loi de Coulomb, la force exercée par $Q$ sur $q$ est :
\[
\vec{F} = k\,\frac{Q\,q}{r^2}\,\vec{u}_{OM} \quad \text{avec } \vec{u}_{OM} = \frac{\overrightarrow{OM}}{r} \text{ (vecteur unitaire radial sortant)}
\]
En utilisant la définition $\vec{F} = q\,\vec{E}(M)$, on simplifie par $q$ (non nul) :

\begin{propriete}[Champ d'une charge ponctuelle]
Le champ électrostatique créé en $M$ par une charge ponctuelle $Q$ placée en $O$ est :
\[
\boxed{\vec{E}(M) = k\,\frac{Q}{r^2}\,\vec{u}_{OM}}
\]
où $k = \dfrac{1}{4\pi\varepsilon_0} \approx \SI{9,00e9}{N.m^2/C^2}$, $r = OM$ et $\vec{u}_{OM}$ est le vecteur unitaire \textbf{sortant} de $O$ vers $M$.
\begin{itemize}
    \item \textbf{Direction :} Radiale (droite $OM$).
    \item \textbf{Sens :} \textbf{Centrifuge} (sortant) si $Q > 0$ ; \textbf{centripète} (rentrant) si $Q < 0$.
    \item \textbf{Intensité :} $E(M) = k\,\dfrac{|Q|}{r^2}$. Elle ne dépend que de la distance $r$ (symétrie sphérique).
\end{itemize}
\end{propriete}

\begin{aretenir}[Analyse dimensionnelle]
Vérifions l'unité de $k\,\dfrac{Q}{r^2}$ :
$[k] = \SI{}{N.m^2.C^{-2}}$, $[Q] = \SI{}{C}$, $[r^2] = \SI{}{m^2}$.
$\left[k\,\frac{Q}{r^2}\right] = \SI{}{N.m^2.C^{-2}} \cdot \SI{}{C} \cdot \SI{}{m^{-2}} = \SI{}{N/C}$. C'est correct.
\end{aretenir}

\begin{popfigure}
\begin{tikzpicture}[scale=1.2, >=Stealth]
% Charge positive Q > 0
\begin{scope}[shift={(-4,0)}]
    \fill[popblue!20] (0,0) circle (0.6);
    \node[popblue, font=\bfseries\large] at (0,0) {$+Q$};
    \foreach \angle in {0,30,60,90,120,150,180,210,240,270,300,330} {
        \draw[popblue, thick, ->] (\angle:0.8) -- (\angle:2.5);
    }
    \draw[popmuted, dashed] (0,0) circle (1.5);
    \draw[popmuted, dashed] (0,0) circle (2.2);
    \node[below] at (0,-2.8) {\textbf{Charge $Q>0$ : champ centrifuge}};
    % Vecteur unitaire
    \draw[poporange, very thick, ->] (0,0) -- (1.5,0) node[midway, above] {$\vec{u}_{OM}$};
    \draw[popgreen, very thick, ->] (0,0) -- (2.2,0) node[midway, below] {$\vec{E}$};
\end{scope}

% Charge negative Q < 0
\begin{scope}[shift={(4,0)}]
    \fill[poppink!20] (0,0) circle (0.6);
    \node[poppink, font=\bfseries\large] at (0,0) {$-Q$};
    \foreach \angle in {0,30,60,90,120,150,180,210,240,270,300,330} {
        \draw[poppink, thick, ->] (\angle:2.5) -- (\angle:0.8);
    }
    \draw[popmuted, dashed] (0,0) circle (1.5);
    \draw[popmuted, dashed] (0,0) circle (2.2);
    \node[below] at (0,-2.8) {\textbf{Charge $Q<0$ : champ centripète}};
    % Vecteur unitaire
    \draw[poporange, very thick, ->] (0,0) -- (1.5,0) node[midway, above] {$\vec{u}_{OM}$};
    \draw[popgreen, very thick, ->] (2.2,0) -- (0,0) node[midway, below] {$\vec{E}$};
\end{scope}
\end{tikzpicture}
\legende{Champ électrostatique radial créé par une charge ponctuelle positive (centrifuge) et négative (centripète). Les cercles pointillés sont des surfaces équipotentielles (sphères).}
\end{popfigure}

\courssub{Principe de superposition des champs}

Dans la réalité, un système contient souvent plusieurs charges sources $Q_1, Q_2, \dots, Q_n$. Comment déterminer le champ total $\vec{E}(M)$ en un point $M$ ?

L'expérience (et la théorie relativiste de l'électromagnétisme) montre que les champs s'additionnent \textbf{vectoriellement}. C'est le \textbf{principe de superposition}.

\begin{propriete}[Théorème : Principe de superposition des champs électrostatiques]
Le champ électrostatique $\vec{E}(M)$ créé en un point $M$ par un ensemble de charges ponctuelles $Q_1, Q_2, \dots, Q_n$ est la \textbf{somme vectorielle} des champs créés par chaque charge prise isolément :
\[
\boxed{\vec{E}(M) = \sum_{i=1}^{n} \vec{E}_i(M) = \sum_{i=1}^{n} k\,\frac{Q_i}{r_i^2}\,\vec{u}_i}
\]
où $\vec{E}_i(M)$ est le champ créé par $Q_i$ seule, $r_i = Q_iM$ et $\vec{u}_i$ est le vecteur unitaire \textbf{sortant} de $Q_i$ vers $M$.
\end{propriete}

\begin{attention}[Erreur fatale : additionner les normes !]
$\vec{E} = \vec{E}_1 + \vec{E}_2 \quad\Rightarrow\quad E = \|\vec{E}_1 + \vec{E}_2\|$.
En général, $E \neq E_1 + E_2$ (sauf si $\vec{E}_1$ et $\vec{E}_2$ sont colinéaires et de même sens).
\textbf{Toujours} faire la somme vectorielle (par composantes ou géométrie / parallélogramme) avant de calculer la norme du résultant.
\end{attention}

\courssub{Méthode : Déterminer le champ résultant en un point}

\begin{methode}[Calcul du champ électrostatique résultant en $M$]
\begin{enumerate}
    \item \textbf{Représenter la situation :} Schéma soigné avec les charges sources $Q_i$, le point $M$, les distances $r_i$, les vecteurs unitaires $\vec{u}_i$ (sortants de $Q_i$ vers $M$).
    \item \textbf{Déterminer chaque champ élémentaire $\vec{E}_i$ :}
    \begin{itemize}
        \item Intensité : $E_i = k\,\dfrac{|Q_i|}{r_i^2}$.
        \item Direction : Droite $(Q_i M)$.
        \item Sens : \textbf{Centrifuge} si $Q_i > 0$ (repousse $q>0$), \textbf{centripète} si $Q_i < 0$ (attire $q>0$).
        \item \textbf{Dessiner $\vec{E}_i$ sur le schéma} (origine en $M$).
    \end{itemize}
    \item \textbf{Exploiter les symétries :} Si deux charges sont symétriques par rapport à une droite ou un point, leurs champs ont des composantes qui s'annulent ou s'additionnent simplement. \emph{C'est le secret des calculs rapides au Bac.}
    \item \textbf{Sommer vectoriellement :}
    \begin{itemize}
        \item Soit par la méthode du parallélogramme (construction géométrique).
        \item Soit par projection sur un repère orthonormé $(O, \vec{x}, \vec{y})$ : $E_x = \sum E_{i,x}$, $E_y = \sum E_{i,y}$.
    \end{itemize}
    \item \textbf{Exprimer le résultant :} $\vec{E} = E_x\,\vec{x} + E_y\,\vec{y}$ ; intensité $E = \sqrt{E_x^2 + E_y^2}$ ; angle $\theta$ tel que $\tan\theta = E_y/E_x$.
    \item \textbf{Conclure avec unité et sens.}
\end{enumerate}
\end{methode}

\courssub{Applications : deux charges ponctuelles}

\paragraph{Cas 1 : Point sur la médiatrice du segment $[Q_1 Q_2]$ (charges de même signe).}
Soient $Q_1 = Q_2 = Q > 0$ en $A$ et $B$, $AB = 2a$. Étudions le champ en $M$ sur la médiatrice, à une distance $x$ du milieu $O$.
Par symétrie, les composantes horizontales ($x$) s'annulent. Les composantes verticales ($y$) s'additionnent.
$r = MA = MB = \sqrt{a^2 + x^2}$.
$\vec{E}_A$ et $\vec{E}_B$ ont même intensité $E_0 = k\,\dfrac{Q}{a^2+x^2}$.
Leurs vecteurs unitaires sortants font un angle $\alpha$ avec la verticale tel que $\cos\alpha = \dfrac{x}{r}$.
$\vec{E} = 2 E_0 \cos\alpha\,\vec{y} = 2\,k\,\dfrac{Q}{a^2+x^2}\,\dfrac{x}{\sqrt{a^2+x^2}}\,\vec{y} = \dfrac{2kQx}{(a^2+x^2)^{3/2}}\,\vec{y}$.
Au centre $O$ ($x=0$) : $\vec{E} = \vec{0}$ (point neutre).

\paragraph{Cas 2 : Point sur l'axe $(AB)$ (charges de signes opposés : dipôle).}
Soient $Q_1 = +q$ en $A$, $Q_2 = -q$ en $B$, $AB = 2a$. Point $M$ sur l'axe, à l'extérieur du segment, distance $OM = x > a$.
$\vec{E}_A$ (centrifuge) et $\vec{E}_B$ (centripète) sont \textbf{colinéaires et de même sens} (vers la droite si $M$ est à droite de $B$).
$E = kq\left(\dfrac{1}{(x-a)^2} + \dfrac{1}{(x+a)^2}\right)$.
Au milieu $O$ ($x=0$) : $\vec{E}_A$ et $\vec{E}_B$ sont colinéaires, même sens (vers $B$, la charge négative).
$E(O) = 2 \times k\,\dfrac{q}{a^2} = \dfrac{2kq}{a^2}$ dirigé de $+q$ vers $-q$.

\begin{exemplebox}[Champ au milieu de deux charges opposées]
\textbf{Énoncé :} Deux charges $q_1 = +\SI{4,0}{nC}$ et $q_2 = -\SI{4,0}{nC}$ sont placées en $A$ et $B$ telles que $AB = \SI{20,0}{cm}$. Déterminer le champ électrostatique au milieu $O$ du segment $[AB]$.

\textbf{Solution :}
\begin{enumerate}
    \item \textbf{Données :} $q_1 = +\SI{4,0e-9}{C}$, $q_2 = -\SI{4,0e-9}{C}$, $AO = OB = a = \SI{0,100}{m}$. $k = \SI{9,00e9}{N.m^2/C^2}$.
    \item \textbf{Champ créé par $q_1$ en $O$ :} $q_1 > 0 \Rightarrow$ champ centrifuge.
    $\vec{E}_1 = k\,\dfrac{q_1}{a^2}\,\vec{u}_{AO}$ (vecteur unitaire de $A$ vers $O$, donc vers la droite).
    $E_1 = \SI{9,00e9}{N.m^2/C^2} \times \dfrac{\SI{4,0e-9}{C}}{(\SI{0,100}{m})^2} = \SI{3600}{N/C}$.
    \item \textbf{Champ créé par $q_2$ en $O$ :} $q_2 < 0 \Rightarrow$ champ centripète (vers $q_2$).
    $\vec{E}_2 = k\,\dfrac{|q_2|}{a^2}\,\vec{u}_{OB}$ (vecteur unitaire de $O$ vers $B$, donc vers la droite).
    $E_2 = \SI{3600}{N/C}$.
    \item \textbf{Somme vectorielle :} $\vec{E}_1$ et $\vec{E}_2$ sont \textbf{colinéaires, même direction, même sens} (vers la droite, de $+q$ vers $-q$).
    $\vec{E} = \vec{E}_1 + \vec{E}_2 = (\SI{3600}{N/C} + \SI{3600}{N/C})\,\vec{x} = \SI{7200}{N/C}\,\vec{x}$.
\end{enumerate}
\fbox{\textbf{Au milieu $O$, le champ vaut $\SI{7,2e3}{N/C}$ (ou $\SI{7,2e3}{V/m}$), dirigé de la charge positive vers la charge négative.}}
\end{exemplebox}

\begin{popfigure}
\begin{tikzpicture}[scale=1.5, >=Stealth]
% Charges
\coordinate (A) at (-1.5,0);
\coordinate (O) at (0,0);
\coordinate (B) at (1.5,0);
\coordinate (M) at (0,1.2); % Point sur médiatrice pour démo

% Charges dessin
\fill[popblue!20] (A) circle (0.15);
\node[popblue, font=\bfseries] at (A) {$+q$};
\fill[poppink!20] (B) circle (0.15);
\node[poppink, font=\bfseries] at (B) {$-q$};

% Vecteurs au milieu O
\draw[popblue, very thick, ->] (O) -- ++(0.8,0) node[midway, above] {$\vec{E}_1$};
\draw[poppink, very thick, ->] (O) -- ++(0.8,0) node[midway, below] {$\vec{E}_2$};
\draw[popgreen, ultra thick, ->] (O) -- ++(1.8,0) node[midway, above, popgreen] {$\vec{E} = \vec{E}_1+\vec{E}_2$};

% Vecteurs en M (sur médiatrice) - démo construction parallélogramme
\draw[popmuted, thin, dashed] (A) -- (M) -- (B);
\draw[popblue, thick, ->] (M) -- ++($(A)-(M)$) node[midway, left] {$\vec{E}_A$}; % Centrifuge from A
\draw[poppink, thick, ->] (M) -- ++($(M)-(B)$) node[midway, right] {$\vec{E}_B$}; % Centripete to B
% Parallélogramme
\draw[popmuted, thin, dashed] (M) -- ++($(A)-(M)$) -- ++($(M)-(B)$) -- ++($(B)-(M)$) -- cycle;
\draw[popgreen, ultra thick, ->] (M) -- ++($(A)-(M) + (M)-(B)$) node[midway, above left] {$\vec{E}_{rés}$};
\node[above] at (M) {$M$};
\node[below] at (O) {$O$};
\end{tikzpicture}
\legende{Construction vectorielle du champ résultant. À gauche : au milieu $O$ d'un dipôle, $\vec{E}_1$ et $\vec{E}_2$ s'additionnent colinéairement. À droite : en un point $M$ de la médiatrice, construction par la règle du parallélogramme (symétrie $\Rightarrow$ résultant vertical).}
\end{popfigure}

\courssub{Lignes de champ : représentation graphique du champ}

Le vecteur $\vec{E}(M)$ étant défini en chaque point $M$, on peut tracer des courbes tangentes à ce vecteur à chaque point. Ce sont les \textbf{lignes de champ}.

\begin{definition}[Ligne de champ électrostatique]
Une ligne de champ est une courbe orientée telle qu'en tout point $M$ de cette courbe, la tangente est dirigée suivant le vecteur champ $\vec{E}(M)$ (dans le sens de $\vec{E}$).
\begin{itemize}
    \item \textbf{Propriété 1 :} Les lignes de champ \textbf{ne se coupent jamais} (unicité de la tangente $\Leftrightarrow$ unicité de $\vec{E}$).
    \item \textbf{Propriété 2 :} Elles partent des charges \textbf{positives} (ou de l'infini) et aboutissent aux charges \textbf{négatives} (ou à l'infini).
    \item \textbf{Propriété 3 :} La \textbf{densité} des lignes de champ (nombre par unité de surface perpendiculaire) est proportionnelle à l'intensité $E$ du champ.
\end{itemize}
\end{definition}

\begin{savaistu}[Faraday et les « lignes de force »]
Michael Faraday (1791--1867), fils d'un forgeron, peu instruit en mathématiques mais génie de l'intuition physique, a refusé l'action à distance newtonienne. Il a imaginé l'espace rempli de « lignes de force » tendues comme des élastiques (tension longitudinale) et se repoussant latéralement (pression transversale). Cette vision qualitative a guidé Maxwell vers ses équations. Au Cameroun, les lignes à haute tension de SONATREL (ex: Lom Pangar -- Yaoundé) créent un champ électrique intense autour des câbles ; c'est pourquoi on respecte des distances de sécurité (servitudes) : le champ ionise l'air (effet couronne) et peut perturber les appareils ou la santé.
\end{savaistu}

\paragraph{Cartographies classiques (à connaître par cœur pour le Bac).}

\begin{popfigure}
\begin{tikzpicture}[scale=0.9, >=Stealth]
% Styles
\tikzset{charge+/.style={circle, fill=popblue!20, draw=popblue, thick, minimum size=6mm, inner sep=0}}
\tikzset{charge-/.style={circle, fill=poppink!20, draw=poppink, thick, minimum size=6mm, inner sep=0}}
\tikzset{fieldline/.style={popmuted, thick, ->}}

% 1. Charge unique +
\begin{scope}[shift={(-6,0)}]
\node[charge+] (Q1) at (0,0) {$+Q$};
\foreach \a in {0,30,...,330} {
    \draw[fieldline] (Q1.\a) -- ++(\a:2.5);
}
\node[below] at (0,-3) {\textbf{(a) Charge $+Q$ isolée}};
\end{scope}

% 2. Charge unique -
\begin{scope}[shift={(0,0)}]
\node[charge-] (Q2) at (0,0) {$-Q$};
\foreach \a in {0,30,...,330} {
    \draw[fieldline] ++(\a:2.5) -- (Q2.\a);
}
\node[below] at (0,-3) {\textbf{(b) Charge $-Q$ isolée}};
\end{scope}

% 3. Deux charges même signe ++
\begin{scope}[shift={(6,0)}]
\node[charge+] (Q3) at (-1,0) {$+q$};
\node[charge+] (Q4) at (1,0) {$+q$};
% Lignes radiales sortantes
\foreach \a in {0,30,...,330} {
    \draw[fieldline] (Q3.\a) -- ++(\a:1.5);
    \draw[fieldline] (Q4.\a) -- ++(\a:1.5);
}
% Lignes de fuite latérales (approximation visuelle)
\foreach \y in {-1.5,-1.0,-0.5,0.5,1.0,1.5} {
    \draw[fieldline, popmuted!50] (-3,\y) .. controls (-0.5,\y*0.5) and (0.5,\y*0.5) .. (3,\y);
}
\node[below] at (0,-3) {\textbf{(c) Deux charges $+q$ (répulsion)}};
\end{scope}

% 4. Dipôle +-
\begin{scope}[shift={(-3,-4)}]
\node[charge+] (Q5) at (-1,0) {$+q$};
\node[charge-] (Q6) at (1,0) {$-q$};
% Lignes partant de +q vers -q (axe)
\draw[fieldline] (Q5.0) -- (Q6.180);
% Lignes courbes (approximation par arcs de cercle orthogonaux aux équipotentielles)
% Équation différentielle dy/dx = Ey/Ex. Pour dipôle, lignes = cercles passant par les charges.
\foreach \r in {1.2, 1.6, 2.2, 3.0} {
    \foreach \s in {1,-1} {
        \draw[fieldline] (-1,0) .. controls ++(0.5,\s*0.8) and ++(-0.5,\s*0.8) .. (1,0);
        \draw[fieldline] (-1,0) .. controls ++(0.2,\s*1.5) and ++(-0.2,\s*1.5) .. (1,0);
    }
}
% Lignes allant à l'infini
\foreach \a in {60,120,240,300} {
    \draw[fieldline] (Q5.\a) -- ++(\a:2.5);
}
\foreach \a in {60,120,240,300} {
    \draw[fieldline] ++(\a:2.5) -- (Q6.\a);
}
\node[below] at (0,-3) {\textbf{(d) Dipôle $+q/-q$ (attraction)}};
\end{scope}

% 5. Deux charges même signe --
\begin{scope}[shift={(3,-4)}]
\node[charge-] (Q7) at (-1,0) {$-q$};
\node[charge-] (Q8) at (1,0) {$-q$};
\foreach \a in {0,30,...,330} {
    \draw[fieldline] ++(\a:1.5) -- (Q7.\a);
    \draw[fieldline] ++(\a:1.5) -- (Q8.\a);
}
\foreach \y in {-1.5,-1.0,-0.5,0.5,1.0,1.5} {
    \draw[fieldline, popmuted!50] (3,\y) .. controls (0.5,\y*0.5) and (-0.5,\y*0.5) .. (-3,\y);
}
\node[below] at (0,-3) {\textbf{(e) Deux charges $-q$ (répulsion)}};
\end{scope}
\end{tikzpicture}
\legende{Cartographies des lignes de champ électrostatique. (a) Charge positive : lignes sortantes radiales. (b) Charge négative : lignes entrantes radiales. (c) Deux charges positives : répulsion, point neutre au milieu ($E=0$). (d) Dipôle : lignes partant de $+q$ vers $-q$, densité maximale près des charges. (e) Deux charges négatives : symétrique de (c).}
\end{popfigure}

\courssub{Exercice résumé : superposition et symétrie}

\begin{exoresolu}[Champ sur l'axe d'un quadripôle simple]
\textbf{Énoncé :} Quatre charges sont placées aux sommets d'un carré $ABCD$ de côté $a = \SI{10,0}{cm}$ : $q_A = +q$, $q_B = -q$, $q_C = +q$, $q_D = -q$ avec $q = \SI{2,0}{nC}$. Déterminer le champ électrostatique au centre $O$ du carré.

\textbf{Solution détaillée :}
\begin{enumerate}
    \item \textbf{Géométrie :} $O$ est centre du carré. $OA = OB = OC = OD = r = \dfrac{a}{\sqrt{2}} = \SI{0,0707}{m}$.
    \item \textbf{Champs élémentaires (intensité) :}
    $E_0 = k\,\dfrac{q}{r^2} = \num{9,00e9} \times \dfrac{\num{2,0e-9}}{(\num{0,0707})^2} = \SI{18,0e0}{0,0050} \approx \SI{3600}{N/C}$.
    \item \textbf{Directions et sens (vecteurs unitaires sortants des charges vers $O$) :}
    \begin{itemize}
        \item $A(+q)$ : centrifuge $\Rightarrow$ $\vec{E}_A$ dirigé $\overrightarrow{AO}$ (vers $O$).
        \item $C(+q)$ : centrifuge $\Rightarrow$ $\vec{E}_C$ dirigé $\overrightarrow{CO}$ (vers $O$).
        \item $B(-q)$ : centripète $\Rightarrow$ $\vec{E}_B$ dirigé $\overrightarrow{OB}$ (vers $B$, donc \textbf{vers $O$} aussi car $O$ est entre $B$ et $O$... attention ! Vecteur unitaire sortant de $B$ vers $O$ est $\vec{u}_{BO}$. Champ centripète = vers $B$ = opposé à $\vec{u}_{BO}$ = sens $\overrightarrow{OB}$).
        \item $D(-q)$ : centripète $\Rightarrow$ $\vec{E}_D$ dirigé $\overrightarrow{OD}$ (vers $D$).
    \end{itemize}
    \item \textbf{Symétrie :}
    $\vec{E}_A$ et $\vec{E}_C$ sont opposées ($\overrightarrow{AO} = -\overrightarrow{CO}$) $\Rightarrow$ $\vec{E}_A + \vec{E}_C = \vec{0}$.
    $\vec{E}_B$ et $\vec{E}_D$ sont opposées ($\overrightarrow{OB} = -\overrightarrow{OD}$) $\Rightarrow$ $\vec{E}_B + \vec{E}_D = \vec{0}$.
    \item \textbf{Résultant :} $\vec{E}(O) = \vec{0}$.
\end{enumerate}
\fbox{\textbf{Au centre de cette configuration alternée, le champ est nul par symétrie.}}
\end{exoresolu}

\begin{exercice}[À toi de jouer : Dipôle sur la médiatrice]
Deux charges $q_1 = +\SI{8,0}{nC}$ et $q_2 = -\SI{8,0}{nC}$ sont placées en $A(-10\,\text{cm}, 0)$ et $B(+10\,\text{cm}, 0)$.
\begin{enumerate}
    \item Dessiner le système et les vecteurs $\vec{E}_1$, $\vec{E}_2$ en un point $M(0, y)$ de la médiatrice ($y>0$).
    \item Montrer que le champ résultant est horizontal. Déterminer son sens.
    \item Calculer l'intensité $E$ pour $y = \SI{10,0}{cm}$. Comparer avec le champ au milieu $O$.
\end{enumerate}
\end{exercice}

\begin{corrige}[Exercice : Dipôle sur la médiatrice]
\begin{enumerate}
    \item \textbf{Schéma :} $M(0, y)$. $A(-a,0)$, $B(a,0)$ avec $a=\SI{0,10}{m}$.
    $\vec{E}_1$ (charge $+q$ en $A$) : centrifuge $\Rightarrow$ direction $AM$, sens $A \to M$.
    $\vec{E}_2$ (charge $-q$ en $B$) : centripète $\Rightarrow$ direction $MB$, sens $M \to B$.
    \item \textbf{Symétrie :} Le triangle $AMB$ est isocèle. $\vec{E}_1$ et $\vec{E}_2$ ont même intensité $E_0 = kq/(a^2+y^2)$.
    Leurs composantes verticales ($y$) sont opposées $\Rightarrow$ s'annulent.
    Leurs composantes horizontales ($x$) sont dans le même sens (vers la droite, de $+q$ vers $-q$).
    $\vec{E}$ est \textbf{horizontal, dirigé de la charge positive vers la charge négative}.
    \item \textbf{Calcul pour $y = a = \SI{0,10}{m}$ :}
    $r = \sqrt{a^2+a^2} = a\sqrt{2}$.
    $E_0 = kq/(2a^2)$.
    Composante horizontale de $\vec{E}_1$ : $E_{0x} = E_0 \cos 45^\circ = E_0 \frac{\sqrt{2}}{2} = \frac{kq}{2a^2}\frac{\sqrt{2}}{2} = \frac{kq\sqrt{2}}{4a^2}$.
    Même chose pour $\vec{E}_2$.
    $E = 2 E_{0x} = \frac{kq\sqrt{2}}{2a^2}$.
    Numérique : $E = \frac{\num{9,00e9} \times \num{8,0e-9} \times \sqrt{2}}{2 \times (\num{0,10})^2} = \frac{72 \times 1,414}{0,02} \approx \SI{5091}{N/C}$.
    \textbf{Au milieu $O$ ($y=0$) :} $E(O) = 2 \times kq/a^2 = \SI{14400}{N/C}$.
    Le champ diminue quand on s'éloigne de l'axe sur la médiatrice.
\end{enumerate}
\end{corrige}

\courssub{Synthèse de la section}

\begin{aretenir}[Ce qu'il faut retenir sur le champ de charges ponctuelles]
\begin{itemize}
    \item \textbf{Définition :} $\vec{E}(M) = \vec{F}/q$ (force sur charge test positive unitaire). Unité : \SI{}{N/C} ou \SI{}{V/m}.
    \item \textbf{Charge ponctuelle $Q$ :} $\vec{E}(M) = k\dfrac{Q}{r^2}\vec{u}_{OM}$. Radial. Centrifuge ($Q>0$), centripète ($Q<0$). $E \propto 1/r^2$.
    \item \textbf{Superposition :} $\vec{E} = \sum \vec{E}_i$. \textbf{Somme vectorielle obligatoire.}
    \item \textbf{Méthode :} Schéma $\to$ $\vec{E}_i$ (intensité, direction, sens) $\to$ Symétries $\to$ Composantes $\to$ Résultant.
    \item \textbf{Lignes de champ :} Tangentes à $\vec{E}$, orientées selon $\vec{E}$, ne se coupent pas, partent de $+$ vers $-$, densité $\propto E$.
    \item \textbf{Configurations types :} Charge isolée, dipôle ($+q/-q$), deux charges même signe (point neutre au milieu).
\end{itemize}
\end{aretenir}

\begin{attention}[Check-list avant l'épreuve]
\begin{itemize}
    \item [ ] Ai-je dessiné les vecteurs $\vec{E}_i$ \textbf{au point $M$} (origine en $M$) ?
    \item [ ] Ai-je respecté le sens : \textbf{sortant} pour $Q>0$, \textbf{entrant} pour $Q<0$ ?
    \item [ ] Ai-je utilisé la symétrie pour annuler des composantes avant de calculer ?
    \item [ ] Ai-je additionné les \textbf{composantes} ($x$ et $y$) et non les normes ?
    \item [ ] L'unité finale est-elle en \SI{}{N/C} (ou \SI{}{V/m}) ?
    \item [ ] Pour les lignes de champ : flèches du $+$ vers le $-$, pas de croisement, densité respectée.
\end{itemize}
\end{attention}

\coursec{Champ électrique uniforme et condensateur plan}

Dans la section précédente, nous avons étudié le champ créé par des charges ponctuelles : un champ \emph{radial}, non uniforme, dont la norme décroît comme $1/r^2$. Nous allons maintenant découvrir une configuration très différente, fondamentale en électronique et en industrie : le \textbf{condensateur plan}. C'est le dispositif qui permet de créer, dans une région de l'espace, un champ électrique \emph{constant} — un outil précieux pour accélérer des particules, dévier des faisceaux d'électrons (anciens écrans cathodiques), ou stocker de l'énergie (flash d'appareil photo, alimentation de secours d'hôpital).

\courssub{Description du condensateur plan}

Un \cle{condensateur plan} est constitué de deux \textbf{armatures} métalliques planes, parallèles, de grande surface $S$ par rapport à leur distance de séparation $d$ ($S \gg d^2$), séparées par un \textbf{diélectrique} (isolant : air, vide, papier, mica, céramique, etc.). Les armatures sont reliées aux bornes d'un générateur de tension (pile, secteur, source haute tension).

\begin{popfigure}
\begin{tikzpicture}[scale=0.9, >=Stealth, european]
% Armatures
\draw[fill=popblue!20, draw=popblue, thick] (-3,1.5) rectangle (3,1.8);
\draw[fill=poporange!20, draw=poporange, thick] (-3,-1.8) rectangle (3,-1.5);
% Charges sur armatures
\foreach \x in {-2.5,-1.5,-0.5,0.5,1.5,2.5} {
    \node[popblue, font=\small] at (\x,1.65) {$+$};
    \node[poporange, font=\small] at (\x,-1.65) {$-$};
}
% Générateur
\draw[thick] (4,1.65) -- (5,1.65);
\draw[thick] (4,-1.65) -- (5,-1.65);
\draw (5,1.65) to[battery1, l={$U_{AB}$}, invert] (5,-1.65);
% Lignes de champ
\foreach \x in {-2.5,-1.5,-0.5,0.5,1.5,2.5} {
    \draw[popgreen, thick, ->] (\x,1.4) -- (\x,-1.3);
}
% Légendes plaques
\node[align=center] at (-3.8,1.65) {Armature A\\(potentiel $V_A$)};
\node[align=center] at (-3.8,-1.65) {Armature B\\(potentiel $V_B$)};
\node[align=center] at (0,2.3) {Charges $+Q$};
\node[align=center] at (0,-2.3) {Charges $-Q$};
% Cotes
\draw[|<->|, popdark] (3.3,1.65) -- (3.3,-1.65) node[midway, right] {$d$};
\draw[|<->|, popdark] (-3,2.0) -- (3,2.0) node[midway, above] {$S$ (grande)};
\end{tikzpicture}
\legende{Condensateur plan alimenté par un générateur de tension $U_{AB}$. L'armature A (reliée au $+$ du générateur) porte une charge $+Q$, l'armature B une charge $-Q$. Les lignes de champ (vertes) sont parallèles et équidistantes au centre.}
\end{popfigure}

\begin{definition}[Condensateur plan]
Un \textbf{condensateur plan} est un système de deux conducteurs plans, parallèles, de même surface $S$, séparés par une distance $d$ (diélectrique), portant des charges opposées $+Q$ et $-Q$ et maintenus à une différence de potentiel $U_{AB} = V_A - V_B$.
\end{definition}

\courssub{Cartographie du champ : découverte de l'uniformité}

Si l'on place une petite bille conductrice chargée (ou un champignon de test) entre les armatures, loin des bords, on observe qu'elle subit une force \textbf{constante} en direction, en sens et en norme, quelle que soit sa position. Le tracé des lignes de champ (expérience des graines de semoule dans l'huile de ricin, ou simulation numérique) montre :
\begin{itemize}
    \item Des lignes \textbf{parallèles} entre elles.
    \item Des lignes \textbf{équidistantes} (même densité partout).
    \item Un sens unique : de l'armature positive ($+Q$, potentiel $V_A$) vers l'armature négative ($-Q$, potentiel $V_B$).
\end{itemize}
Aux abords des bords (\emph{effets de bord}), les lignes s'incurvent ; on néglige cette zone en supposant $S \gg d^2$ (condensateur \emph{idéal}).

\begin{definition}[Champ électrique uniforme]
Un champ électrique est dit \textbf{uniforme} (ou constant) dans une région de l'espace si le vecteur champ $\vec{E}$ y a \textbf{la même direction, le même sens et la même norme} en tout point.
\par Dans un condensateur plan idéal, $\vec{E}$ est uniforme entre les armatures, dirigé de $A$ ($+$) vers $B$ ($-$), et \textbf{nul} à l'extérieur (blindage électrostatique).
\end{definition}

\begin{popfigure}
\begin{tikzpicture}[scale=1.1, >=Stealth]
% Armatures
\draw[fill=popblue!10, draw=popblue, very thick] (-3,1) rectangle (3,1.2);
\draw[fill=poporange!10, draw=poporange, very thick] (-3,-1) rectangle (3,-1.2);
% Vecteurs E uniformes
\foreach \x in {-2.5,-1.5,-0.5,0.5,1.5,2.5} {
    \draw[popgreen, very thick, ->] (\x,0.8) -- (\x,-0.8) node[midway, right=2pt] {$\vec{E}$};
}
% Points de mesure
\foreach \y in {-0.5, 0, 0.5} {
    \fill[popdark] (0,\y) circle (1.5pt);
    \node[left] at (-2.8,\y) {$M$};
}
% Cote d
\draw[|<->|, poppurple] (3.4,1) -- (3.4,-1) node[midway, right] {$d$};
% Potentiels
\node[popblue] at (-3.5,1.1) {$V_A$};
\node[poporange] at (-3.5,-1.1) {$V_B$};
\draw[poppurple, <->] (-3.8,1.1) -- (-3.8,-1.1) node[midway, left] {$U_{AB}$};
\end{tikzpicture}
\legende{Champ uniforme $\vec{E}$ entre les armatures. La norme $E$ est constante. $U_{AB} = V_A - V_B > 0$.}
\end{popfigure}

\courssub{Relation fondamentale : $E = \dfrac{U}{d}$}

Dans un champ uniforme, la différence de potentiel entre deux points situés sur une même ligne de champ est proportionnelle à leur distance.

\begin{propriete}[Relation $E = U/d$ dans un champ uniforme]
Dans un champ électrique uniforme $\vec{E}$, la différence de potentiel $U_{AB}$ entre deux points $A$ et $B$ situés sur une même ligne de champ, séparés par la distance $d = AB$, vérifie :
\[
E = \frac{U_{AB}}{d} \qquad \text{avec } E = \|\vec{E}\| \text{ en V$\cdot$m$^{-1}$ (ou N$\cdot$C$^{-1}$), } U_{AB} \text{ en V, } d \text{ en m.}
\]
\par \textbf{Conditions de validité :} champ uniforme ; $A$ et $B$ sur une même ligne de champ ; $d$ mesuré le long de la ligne de champ.
\par \textbf{Analyse dimensionnelle :} $[E] = \frac{[U]}{[d]} = \frac{\text{V}}{\text{m}} = \text{V}\cdot\text{m}^{-1} = \frac{\text{J}}{\text{C}\cdot\text{m}} = \frac{\text{N}\cdot\text{m}}{\text{C}\cdot\text{m}} = \text{N}\cdot\text{C}^{-1}$.
\par \emph{Au Bac :} Cette relation est à connaître et à savoir appliquer. Sa démonstration via le travail de la force électrique (ci-dessous) est exigible.
\end{propriete}

\begin{experience}[Démonstration guidée de $E = U/d$ (via le travail de la force électrique)]
On considère une charge de test $q > 0$ déplacée quasi-statiquement de l'armature $B$ ($-$) vers l'armature $A$ ($+$) le long d'une ligne de champ.
\begin{enumerate}
    \item Force électrique : $\vec{F}_e = q\vec{E}$ (constante, colinéaire au déplacement, orientée de $A$ vers $B$).
    \item Déplacement $\vec{d}_{B\to A}$ : orienté de $B$ vers $A$ (sens inverse de $\vec{E}$), de norme $d$.
    \item Travail de $\vec{F}_e$ : $W_{B\to A}(\vec{F}_e) = \vec{F}_e \cdot \vec{d}_{B\to A} = q E d \cos\pi = -q E d$.
    \item Par définition de la différence de potentiel : $U_{AB} = V_A - V_B = \dfrac{W_{B\to A}^{\text{ext}}}{q} = -\dfrac{W_{B\to A}(\vec{F}_e)}{q}$ (forces extérieures opposées aux forces électriques pour un déplacement quasi-statique).
    \item Donc $U_{AB} = -\dfrac{(-q E d)}{q} = E d \quad\Rightarrow\quad \boxed{E = \dfrac{U_{AB}}{d}}$.
\end{enumerate}
\emph{Remarque :} On retrouve que le champ pointe vers le potentiel décroissant ($V_A > V_B$).
\end{experience}

\courssub{Exploitation : méthode et exemple chiffré}

\begin{methode}[Utiliser la relation $E = U/d$]
\begin{enumerate}
    \item \textbf{Identifier l'armature positive} (reliée au $+$ du générateur, potentiel le plus élevé). Le champ $\vec{E}$ part de cette armature.
    \item \textbf{Orienter $\vec{E}$} : de la plaque $+$ vers la plaque $-$ (direction perpendiculaire aux armatures).
    \item \textbf{Convertir la distance $d$ en mètres} (SI impératif).
    \item \textbf{Appliquer $E = U/d$} pour trouver l'inconnue ($E$, $U$ ou $d$).
    \item \textbf{Vérifier l'homogénéité} : $[E] = \text{V}\cdot\text{m}^{-1} = \text{N}\cdot\text{C}^{-1}$.
\end{enumerate}
\end{methode}

\begin{exemplebox}[Calcul du champ dans un condensateur de flash]
Un condensateur plan d'un appareil photo (flash) est alimenté sous une tension $U = \SI{200}{\V}$. La distance entre les armatures est $d = \SI{5,0}{\mm}$.
\begin{enumerate}
    \item Convertir $d$ : $d = \SI{5,0e-3}{\m}$.
    \item Calculer $E$ : $E = \dfrac{\SI{200}{\V}}{\SI{5,0e-3}{\m}} = \SI{4,0e4}{\V\per\m}$.
\end{enumerate}
\par Le champ vaut $\mathbf{\SI{4,0e4}{\V\per\m}}$ (ou $\text{N}\cdot\text{C}^{-1}$), dirigé de l'armature $+$ vers l'armature $-$.
\end{exemplebox}

\begin{attention}[Pièges classiques — À ne pas commettre !]
\begin{itemize}
    \item \textbf{Sens de $\vec{E}$ :} $\vec{E}$ va \textbf{toujours} du potentiel \textbf{élevé} ($+$) vers le potentiel \textbf{faible} ($-$). JAMAIS l'inverse. (Rappel : $\vec{E} = -\overrightarrow{\text{grad}}\,V$).
    \item \textbf{Unités :} Ne pas confondre $E$ (en $\text{V}\cdot\text{m}^{-1}$) et $U$ (en $\text{V}$). $E$ est un \textbf{gradient de potentiel}.
    \item \textbf{Distance $d$ :} C'est la distance \textbf{entre les armatures} (épaisseur du diélectrique), pas la longueur des plaques. En SI : \textbf{mètres}.
    \item \textbf{Champ non uniforme :} La formule $E=U/d$ ne s'applique \textbf{que} dans un champ uniforme (condensateur plan idéal, loin des bords). Pour une charge ponctuelle, $E = kQ/r^2$ et $U = kQ/r$ : $E \neq U/r$ !
\end{itemize}
\end{attention}

\courssub{Applications concrètes et contextes camerounais}

\begin{savaistu}[Le condensateur au Cameroun : du flash à la dépollution]
\begin{itemize}
    \item \textbf{Flash d'appareil photo / téléphone :} Un petit condensateur (quelques $\mu\text{F}$, $U \sim \SI{300}{\V}$) se charge lentement sur la batterie, puis se décharge brutalement ($< \SI{1}{\ms}$) dans une lampe à xénon : puissance crête $\sim \SI{10}{\kW}$ ! C'est l'application directe du stockage d'énergie $\frac{1}{2} C U^2$ (section 6).
    \item \textbf{Précipitateurs électrostatiques (centrales thermiques, cimenteries comme CIMENCAM) :} Des fils à haut potentiel ($U \sim \SI{50}{\kV}$, $d \sim \SI{10}{\cm}$ $\Rightarrow$ $E \sim \SI{5e5}{\V\per\m}$) ionisent l'air ; les particules de poussière se chargent et sont attirées sur des plaques de collecte. Application directe du champ uniforme et des forces sur charges.
    \item \textbf{Écrans et imprimantes laser :} Déviation de faisceaux d'électrons (anciens tubes cathodiques) ou de toner chargé par des champs uniformes créés par des plaques de déviation.
    \item \textbf{Électrification rurale / SONATREL :} Les condensateurs de puissance (batteries de condensateurs) sont installés dans les postes HTA/BT pour \emph{compenser l'énergie réactive} (améliorer le $\cos\varphi$), réduire les pertes sur le réseau et stabiliser la tension — un enjeu majeur pour la qualité de l'électricité distribuée par ENEO.
\end{itemize}
\end{savaistu}

\courssub{Exercice résolu : Déterminer la tension de claquage}

\begin{exoresolu}[Tension maximale avant claquage de l'air]
L'air sec se comporte comme un isolant tant que le champ électrique ne dépasse pas la \textbf{rigidité diélectrique de l'air} : $E_{\text{max}} \approx \SI{3,0e6}{\V\per\m}$. Au-delà, l'air s'ionise (étincelle, éclair).\\
Un condensateur plan à air a une distance entre armatures $d = \SI{2,0}{\mm}$. Quelle est la tension maximale $U_{\text{max}}$ applicable sans provoquer d'arc électrique ?
\tcblower
\textbf{Solution :}
\begin{enumerate}
    \item Données : $E_{\text{max}} = \SI{3,0e6}{\V\per\m}$ ; $d = \SI{2,0}{\mm} = \SI{2,0e-3}{\m}$.
    \item Relation : $E = U/d \;\Rightarrow\; U = E \times d$.
    \item Calcul : $U_{\text{max}} = (\SI{3,0e6}{\V\per\m}) \times (\SI{2,0e-3}{\m}) = \SI{6,0e3}{\V} = \mathbf{\SI{6,0}{\kV}}$.
\end{enumerate}
\par \textbf{Interprétation :} Une simple étincelle de $\SI{2}{\mm}$ (ex. : en touchant une poignée de porte après avoir marché sur un tapis) correspond à $\sim \SI{6}{\kV}$ ! C'est pourquoi les lignes HT (\SI{90}{\kV}, \SI{225}{\kV}) nécessitent de grandes distances d'isolement (chaînes d'isolateurs en porcelaine ou verre, longueur $\sim \SI{1}{\m}\text{--}\SI{2}{\m}$).
\end{exoresolu}

\begin{aretenir}[Synthèse : Champ uniforme et condensateur plan]
\begin{itemize}
    \item \textbf{Condensateur plan :} 2 armatures $//$, charges $+Q/-Q$, tension $U_{AB}$, distance $d$, diélectrique.
    \item \textbf{Champ $\vec{E}$ :} Uniforme entre armatures (parallèle, équidistant), $\vec{E} \perp$ aux armatures, sens $+\to -$, nul à l'extérieur.
    \item \textbf{Relation clé :} $E = \dfrac{U_{AB}}{d}$ \quad ($E$ en $\text{V}\cdot\text{m}^{-1}$, $U$ en $\text{V}$, $d$ en $\text{m}$).
    \item \textbf{Validité :} $S \gg d^2$ (loin des bords), points sur même ligne de champ.
    \item \textbf{Applications :} Flash, précipitateurs, déviation faisceaux, compensation réactive (réseau ENEO/SONATREL).
\end{itemize}
\end{aretenir}

\coursec{Travail de la force électrique, potentiel et d.d.p.}

Dans la section précédente, nous avons découvert le \cle{condensateur plan} et le \cle{champ électrique uniforme} $\vec{E}$ qui règne entre ses deux plaques, avec la relation $E = \dfrac{U}{d}$. Nous savons donc désormais qu'une charge $q$ placée entre les plaques subit une force $\vec{F} = q\vec{E}$ constante. Une question naturelle se pose alors : si cette charge se déplace sous l'action de cette force — comme un électron qui file d'une plaque vers l'autre dans un tube cathodique —, quel \cle{travail} cette force fournit-elle ? C'est toute la problématique de cette section, qui nous conduira à une grandeur fondamentale de l'électricité : le \cle{potentiel électrique}.

\courssub{Force électrique dans un champ uniforme : une force constante}

\begin{aretenir}[Force électrique dans un champ uniforme]
Une charge ponctuelle $q$ placée dans un champ électrique uniforme $\vec{E}$ subit la force :
\[ \vec{F} = q\,\vec{E} \]
Comme $\vec{E}$ est uniforme (même direction, même sens, même norme en tout point entre les plaques), la force $\vec{F}$ est \textbf{constante} : elle ne dépend pas de la position de la charge.
\begin{itemize}
\item Si $q > 0$ : $\vec{F}$ a même direction et même sens que $\vec{E}$ (de la plaque positive vers la plaque négative).
\item Si $q < 0$ : $\vec{F}$ est opposée à $\vec{E}$ (de la plaque négative vers la plaque positive).
\item Norme : $F = |q|\,E$.
\end{itemize}
\end{aretenir}

Cette situation est remarquable : c'est l'analogue électrique de la chute libre, où le poids $\vec{P} = m\vec{g}$ est constant. De même que le travail du poids ne dépend que de la dénivellation, nous allons voir que le travail de la force électrique ne dépend que d'une « dénivellation électrique » : la différence de potentiel.

\courssub{Travail de la force électrique lors d'un déplacement de A vers B}

\begin{definition}[Travail de la force électrique dans un champ uniforme]
Une charge $q$ se déplace d'un point A vers un point B dans un champ électrique uniforme $\vec{E}$. La force $\vec{F} = q\vec{E}$ étant constante, son travail est le produit scalaire :
\[ W_{AB}(\vec{F}) = \vec{F} \cdot \overrightarrow{AB} = q\,\vec{E} \cdot \overrightarrow{AB} \]
soit, en notant $\alpha$ l'angle entre $\vec{E}$ et $\overrightarrow{AB}$ :
\[ W_{AB}(\vec{F}) = q\,E\,AB\,\cos\alpha \]
Unités : $W$ en joules (J), $q$ en coulombs (C), $E$ en volts par mètre (V·m$^{-1}$), $AB$ en mètres (m).
\end{definition}

Introduisons $d$, la \textbf{projection de $\overrightarrow{AB}$ sur une ligne de champ}, c'est-à-dire la distance algébrique entre les deux plans perpendiculaires à $\vec{E}$ passant par A et B, comptée positivement dans le sens de $\vec{E}$ : $d = AB\cos\alpha$. On obtient alors la forme opérationnelle :

\[ W_{AB}(\vec{F}) = q\,E\,d \]

\begin{attention}[La distance à utiliser est la projection, pas AB !]
L'erreur la plus fréquente au Baccalauréat consiste à remplacer $d$ par la distance $AB$ réellement parcourue. \textbf{Seule compte la projection de $\overrightarrow{AB}$ sur la direction du champ.} Si la charge se déplace perpendiculairement aux lignes de champ ($\alpha = 90^\circ$), alors $d = 0$ et $W_{AB} = 0$ : la force ne travaille pas, exactement comme le poids lors d'un déplacement horizontal.
\end{attention}

\begin{experience}[Le travail ne dépend pas du chemin suivi]
Considérons deux chemins pour aller de A à B entre les plaques :
\begin{itemize}
\item \textbf{Chemin 1} : le segment direct $[AB]$, faisant l'angle $\alpha$ avec $\vec{E}$.
\item \textbf{Chemin 2} : une brisure A $\to$ C $\to$ B, où $[AC]$ est parallèle à $\vec{E}$ et $[CB]$ perpendiculaire à $\vec{E}$.
\end{itemize}
Sur le chemin 2 : $W_{AC} = qE \cdot AC$ (force et déplacement colinéaires) et $W_{CB} = 0$ (force perpendiculaire au déplacement). Or $AC = AB\cos\alpha = d$. Donc :
\[ W_{\text{chemin 2}} = qEd + 0 = qEd = W_{\text{chemin 1}} \]
Tout chemin peut être décomposé en petits tronçons parallèles et perpendiculaires à $\vec{E}$ ; seuls les tronçons parallèles contribuent, et leur somme vaut toujours $d$. \textbf{Le travail de la force électrique ne dépend donc que des points de départ et d'arrivée, pas du chemin suivi.}
\end{experience}

\begin{popfigure}
\begin{center}
\begin{tikzpicture}[scale=1.05, >=Stealth]
% Plaques
\fill[poppinkL] (-0.3,-2.2) rectangle (0.3,2.2);
\fill[popblueL] (5.7,-2.2) rectangle (6.3,2.2);
\node at (0,2.55) {\textbf{Plaque $+$}};
\node at (6,2.55) {\textbf{Plaque $-$}};
\node at (0,-2.55) {$V_+ = +100$ V};
\node at (6,-2.55) {$V_- = 0$ V};
% Champ E
\foreach \y in {-1.5,-0.75,0,0.75,1.5}{
  \draw[->, popgreen, thick] (0.7,\y) -- (1.7,\y);
}
\node[popgreen] at (1.2,1.85) {$\vec{E}$};
% Points A et B
\coordinate (A) at (1.3,-1.2);
\coordinate (B) at (4.8,1.0);
\coordinate (C) at (4.8,-1.2);
\fill[popink] (A) circle (2.2pt); \node[below left, popink] at (A) {A};
\fill[popink] (B) circle (2.2pt); \node[above right, popink] at (B) {B};
\fill[popmuted] (C) circle (1.5pt); \node[below right, popmuted] at (C) {C};
% Chemins
\draw[->, poporange, very thick] (A) -- (B) node[midway, above left, poporange] {chemin 1};
\draw[->, poppurple, very thick] (A) -- (C) node[midway, below, poppurple] {chemin 2};
\draw[->, poppurple, very thick] (C) -- (B);
% Force sur q>0
\draw[->, poppink, very thick] (A) -- ++(1.1,0) node[below, poppink] {$\vec{F}=q\vec{E}$};
% distance d
\draw[<->, popdark, thick] (1.3,-1.75) -- (4.8,-1.75) node[midway, below] {$d$};
% angle alpha
\draw[popmuted] (2.6,-1.2) arc (0:21:1.3);
\node[popmuted] at (2.9,-0.95) {$\alpha$};
\end{tikzpicture}
\end{center}
\legende{Charge $q>0$ se déplaçant de A vers B dans un champ uniforme. Le travail de $\vec{F}$ est identique sur le chemin direct (orange) et sur le chemin brisé (violet) : il ne dépend que de la projection $d$.}
\end{popfigure}

\begin{propriete}[La force électrique est conservative]
Le travail de la force électrique entre deux points A et B est indépendant du chemin suivi : on dit que la force électrique est une \cle{force conservative}. En particulier, son travail sur un trajet fermé (retour au point de départ) est nul.
\end{propriete}

\courssub{Potentiel électrique et différence de potentiel}

Puisque le travail ne dépend que des points A et B, on peut associer à chaque point de l'espace un nombre, le \cle{potentiel électrique}, de sorte que le travail s'exprime comme une différence — exactement comme le travail du poids s'exprime à l'aide de l'altitude.

\begin{definition}[Potentiel électrique et différence de potentiel]
Le \cle{potentiel électrique} $V$ en un point est une grandeur scalaire, exprimée en \textbf{volts} (V), définie à une constante additive près (on choisit une référence, par exemple $V = 0$ à la plaque négative ou à la masse du circuit).

La \cle{différence de potentiel} (d.d.p.) ou \cle{tension} entre A et B est :
\[ U_{AB} = V_A - V_B \]
Contrairement au potentiel, la tension $U_{AB}$ est parfaitement définie, car la constante arbitraire s'élimine dans la différence.
\end{definition}

\begin{attention}[Signe de $U_{AB}$ : attention à l'ordre des indices]
$U_{AB} = V_A - V_B$ et $U_{BA} = V_B - V_A = -U_{AB}$. Inverser les indices change le signe ! Avant tout calcul, identifie clairement quelle plaque est au potentiel le plus haut (la plaque positive) et écris la tension dans le bon ordre. Une erreur de signe sur $U_{AB}$ fausse le signe du travail, donc toute la conclusion énergétique.
\end{attention}

\courssub{Théorème : $W_{AB}(\vec{F}) = q\,U_{AB}$}

\begin{propriete}[Travail de la force électrique et tension — démonstration exigible]
Dans un champ électrique uniforme, le travail de la force électrique subie par une charge $q$ allant de A à B est :
\[ W_{AB}(\vec{F}) = q\,U_{AB} = q\,(V_A - V_B) \]
\emph{Démonstration.} Choisissons un axe $(Ox)$ porté par les lignes de champ, orienté dans le sens de $\vec{E}$ (du potentiel haut vers le potentiel bas). Soit $d$ la projection de $\overrightarrow{AB}$ sur cet axe (distance algébrique entre les plans équipotentiels passant par A et B). Nous avons établi :
\[ W_{AB}(\vec{F}) = q\,E\,d \]
Dans un champ uniforme, le potentiel varie linéairement le long de $(Ox)$ : la différence de potentiel entre les deux plans équipotentiels vaut $V_A - V_B = E\,d$, c'est-à-dire $U_{AB} = E\,d$. En reportant :
\[ W_{AB}(\vec{F}) = q\,E\,d = q\,U_{AB} \qquad \blacksquare \]
\end{propriete}

\begin{aretenir}[Boucle pédagogique : la relation $E = U/d$ retrouvée]
Remarque la cohérence de l'édifice : la relation $E = \dfrac{U}{d}$ admise à la section 4 se \textbf{déduit} du calcul de travail $W = qEd = qU$, qui impose $Ed = U$. Travail, champ et tension sont trois visages d'une même réalité physique. Retiens la chaîne :
\[ \vec{F} = q\vec{E} \quad\Longrightarrow\quad W_{AB} = qEd \quad\Longleftrightarrow\quad W_{AB} = qU_{AB} \quad\Longleftrightarrow\quad E = \frac{U}{d} \]
\end{aretenir}

\begin{methode}[Calculer le travail de la force électrique en 4 étapes]
\begin{enumerate}
\item \textbf{Identifier le signe de $q$} : la charge est-elle positive ou négative (électron : $q = -e$) ?
\item \textbf{Identifier le sens de $\vec{E}$} : il va de la plaque positive (haut potentiel) vers la plaque négative (bas potentiel).
\item \textbf{Déterminer $U_{AB} = V_A - V_B$} (ou la distance projetée $d$ si l'on utilise $W = qEd$), en veillant à l'ordre des indices.
\item \textbf{Conclure sur le signe} : $W = qU_{AB}$.
\begin{itemize}
\item $W > 0$ : travail \textbf{moteur} — la force électrique favorise le déplacement (déplacement « spontané » de la charge).
\item $W < 0$ : travail \textbf{résistant} — il faut un opérateur extérieur pour imposer ce déplacement.
\end{itemize}
\end{enumerate}
\end{methode}

\begin{exemplebox}[Électron accéléré entre les plaques d'un condensateur]
Un électron ($q = -e = \num{-1,6e-19}$ C) part de la plaque négative d'un condensateur plan et arrive sur la plaque positive. La tension entre les plaques vaut $U = 100$ V.

Plaque négative = point A ($V_A = 0$ V), plaque positive = point B ($V_B = 100$ V). Alors :
\[ W_{AB}(\vec{F}) = q\,U_{AB} = q\,(V_A - V_B) = (-\num{1,6e-19})\times(0 - 100) = +\num{1,6e-17}\ \text{J} \]
Le travail est \textbf{positif} : la force électrique est motrice, l'électron est bel et bien attiré par la plaque positive. C'est exactement le principe du canon à électrons des anciens téléviseurs et des oscilloscopes utilisés au laboratoire.
\end{exemplebox}

\courssub{Surfaces équipotentielles (complément Bac)}

\begin{definition}[Surface équipotentielle]
Une \cle{surface équipotentielle} est l'ensemble des points situés au même potentiel $V$. Dans un condensateur plan, ce sont des \textbf{plans parallèles aux plaques}.
\end{definition}

\begin{propriete}[Équipotentielles et lignes de champ]
Les surfaces équipotentielles sont \textbf{perpendiculaires aux lignes de champ}. Les lignes de champ sont orientées dans le sens des potentiels décroissants. Un déplacement le long d'une équipotentielle ($V_A = V_B$, donc $d = 0$) s'accompagne d'un travail nul : $W_{AB} = q(V_A - V_B) = 0$.
\end{propriete}

\begin{popfigure}
\begin{center}
\begin{tikzpicture}[scale=1.0, >=Stealth]
% Plaques
\fill[poppinkL] (-0.25,-2.0) rectangle (0.25,2.0);
\fill[popblueL] (5.75,-2.0) rectangle (6.25,2.0);
\node at (0,2.3) {$+$ ($V = 100$ V)};
\node at (6,2.3) {$-$ ($V = 0$ V)};
% Lignes de champ
\foreach \y in {-1.5,-0.9,-0.3,0.3,0.9,1.5}{
  \draw[->, popgreen, thick] (0.5,\y) -- (5.5,\y);
}
\node[popgreen] at (3,1.8) {lignes de champ $\vec{E}$};
% Equipotentielles
\foreach \x/\v in {1.5/75, 3.0/50, 4.5/25}{
  \draw[poporange, very thick, dashed] (\x,-1.9) -- (\x,1.9);
  \node[poporange, below] at (\x,-1.95) {$V = \v$ V};
}
\node[poporange] at (4.5,2.3) {équipotentielles};
\end{tikzpicture}
\end{center}
\legende{Dans un condensateur plan, les surfaces équipotentielles (traits orange discontinus) sont des plans parallèles aux plaques, perpendiculaires aux lignes de champ (vertes). Le potentiel décroît dans le sens de $\vec{E}$.}
\end{popfigure}

\courssub{Signe du travail et lien avec l'énergie potentielle}

Analysons tous les cas possibles avec $W_{AB} = q(V_A - V_B)$ :

\begin{center}
\begin{tabularx}{\linewidth}{|l|X|X|}
\hline
\rowcolor{popblueL} \textbf{Situation} & \textbf{Signe de $W_{AB}$} & \textbf{Interprétation} \\
\hline
$q>0$ vers les potentiels décroissants ($V_A > V_B$) & $W > 0$ moteur & La charge positive suit « naturellement » le champ. \\
\hline
$q>0$ vers les potentiels croissants ($V_A < V_B$) & $W < 0$ résistant & Un opérateur doit fournir de l'énergie. \\
\hline
$q<0$ vers les potentiels croissants ($V_A < V_B$) & $W > 0$ moteur & L'électron remonte « naturellement » les potentiels. \\
\hline
$q<0$ vers les potentiels décroissants ($V_A > V_B$) & $W < 0$ résistant & Déplacement forcé. \\
\hline
\end{tabularx}
\end{center}

\begin{savaistu}[Pourquoi « haute tension » ?]
Quand ENEO annonce une ligne à \SI{90}{kV} ou \SI{225}{kV} partant du barrage de Nachtigal, il s'agit d'une \textbf{différence de potentiel} : le travail fourni à chaque coulomb de charge transporté vaut $W = qU$, soit \SI{2,25e5}{J} par coulomb ! Transporter l'énergie sous très haute tension permet de réduire l'intensité du courant, donc les pertes par effet Joule dans les câbles sur les longues distances entre la centrale et les villes.
\end{savaistu}

\begin{aretenir}[Vers l'énergie potentielle électrostatique]
Parce que la force électrique est conservative, on peut lui associer une \cle{énergie potentielle} $E_p$ telle que :
\[ W_{AB}(\vec{F}) = -\Delta E_p = E_{pA} - E_{pB} \quad\text{avec}\quad E_p = qV + \text{cte} \]
Un travail moteur ($W>0$) correspond à une diminution d'énergie potentielle, convertie en énergie cinétique : c'est la clé de l'étude du \textbf{mouvement d'une particule chargée dans un champ uniforme}, objet de la leçon suivante.
\end{aretenir}

\begin{exoresolu}[Travail, tension et potentiels dans un condensateur plan]
\textbf{Énoncé.} Entre deux plaques verticales P$_1$ et P$_2$ distantes de $d = \SI{5,0}{cm}$, on applique une tension $U_{12} = V_1 - V_2 = \SI{200}{V}$. Une charge $q = \num{+2,0e-9}$ C se déplace d'un point A situé à \SI{1,0}{cm} de P$_1$ vers un point B situé à \SI{4,0}{cm} de P$_1$, en suivant un arc de cercle quelconque. On prend $V_2 = 0$.
\begin{enumerate}
\item Calculer le champ $E$ régnant entre les plaques.
\item Calculer les potentiels $V_A$ et $V_B$.
\item Calculer le travail $W_{AB}(\vec{F})$ de deux façons différentes.
\end{enumerate}
\tcblower
\textbf{Solution détaillée.}
\begin{enumerate}
\item Champ uniforme entre les plaques :
\[ E = \frac{U_{12}}{d} = \frac{200}{\num{5,0e-2}} = \SI{4,0e3}{V.m^{-1}} \]
dirigé de P$_1$ (potentiel le plus haut) vers P$_2$.
\item Le potentiel décroît linéairement de P$_1$ vers P$_2$ : en un point situé à la distance $x$ de P$_1$, $V = V_1 - Ex$.
\begin{align*}
V_A &= 200 - \num{4,0e3}\times\num{1,0e-2} = \SI{160}{V} \\
V_B &= 200 - \num{4,0e3}\times\num{4,0e-2} = \SI{40}{V}
\end{align*}
\item \textbf{Première méthode} (via la tension) :
\[ W_{AB} = q(V_A - V_B) = \num{2,0e-9}\times(160-40) = \num{+2,4e-7}\ \text{J} \]
\textbf{Deuxième méthode} (via la projection) : la distance projetée sur une ligne de champ vaut $d' = 4,0 - 1,0 = \SI{3,0e-2}{m}$, comptée dans le sens de $\vec{E}$ :
\[ W_{AB} = qEd' = \num{2,0e-9}\times\num{4,0e3}\times\num{3,0e-2} = \num{+2,4e-7}\ \text{J} \]
Les deux résultats coïncident : le travail ne dépend pas du chemin (l'arc de cercle importe peu). $W_{AB} > 0$ : la charge positive descend spontanément les potentiels, la force est motrice.
\end{enumerate}
\end{exoresolu}

\begin{exercice}[À toi de jouer]
Entre les plaques d'un condensateur plan séparées de $d = \SI{10}{cm}$, règne un champ $E = \SI{5,0e3}{V.m^{-1}}$, dirigé de la plaque M vers la plaque N. Un proton ($q = +e = \num{1,6e-19}$ C) est déplacé de N vers M.
\begin{enumerate}
\item Calculer la tension $U_{MN}$.
\item Calculer le travail de la force électrique au cours de ce déplacement. Commenter son signe.
\item Quelle énergie un opérateur extérieur doit-il fournir, au minimum ?
\end{enumerate}
\end{exercice}

\begin{corrige}[Exercice — À toi de jouer]
\begin{enumerate}
\item $U_{MN} = E\,d = \num{5,0e3}\times\num{0,10} = \SI{500}{V}$ (M est au potentiel le plus haut puisque $\vec{E}$ va de M vers N).
\item Le déplacement est de N vers M : $W_{NM} = q\,U_{NM} = q(V_N - V_M) = \num{1,6e-19}\times(-500) = \num{-8,0e-17}$ J. Travail \textbf{résistant} : le proton positif est repoussé par M, il ne peut pas y aller spontanément.
\item L'opérateur doit fournir au minimum $W_{\text{op}} = -W_{NM} = \num{+8,0e-17}$ J.
\end{enumerate}
\end{corrige}

\begin{aretenir}[L'essentiel de la section]
\begin{itemize}
\item Dans un champ uniforme, $\vec{F} = q\vec{E}$ est constante ; son travail est $W_{AB} = q\,\vec{E}\cdot\overrightarrow{AB} = qEd$, où $d$ est la \textbf{projection} de $\overrightarrow{AB}$ sur une ligne de champ.
\item Le travail ne dépend pas du chemin : la force électrique est \textbf{conservative}.
\item Relation fondamentale : $W_{AB} = q\,U_{AB} = q(V_A - V_B)$, avec $U_{AB}$ en volts.
\item $E = \dfrac{U}{d}$ se déduit de cette relation ; les équipotentielles sont perpendiculaires aux lignes de champ.
\item Signe de $W$ : moteur si la charge suit son mouvement spontané ($q>0$ vers les potentiels décroissants, $q<0$ vers les potentiels croissants), résistant sinon.
\end{itemize}
\end{aretenir}

\coursec{Synthèse et méthodes de résolution}

Nous venons de voir, dans la section précédente, que la force électrique est une force \cle{conservative} : son travail ne dépend que des potentiels de départ et d'arrivée, $W_{AB}(\vec{F}) = q\,U_{AB} = q(V_A - V_B)$. Cette dernière brique achève l'édifice du chapitre : nous disposons désormais d'une \emph{chaîne complète de grandeurs} — charge, force, champ, tension, travail — reliées entre elles par des formules simples. L'objectif de cette section est de t'apprendre à \textbf{naviguer dans cette chaîne} sans te tromper, comme le fera un candidat bien entraîné le jour du Baccalauréat.

\courssub{Bilan des grandeurs et de leurs relations}

\begin{aretenir}[Les cinq grandeurs du chapitre]
\begin{center}
\begin{tabularx}{\linewidth}{|l|c|c|X|}
\hline
\rowcolor{popblueL} \textbf{Grandeur} & \textbf{Symbole} & \textbf{Unité SI} & \textbf{Signification} \\
\hline
Charge électrique & $q$ & coulomb (C) & Propriété de la matière, quantifiée : $q = n\,e$ avec $e = \num{1,6e-19}$ C. \\
\hline
Force électrostatique & $F$ & newton (N) & Interaction entre deux charges : $F = k\dfrac{|q_1 q_2|}{d^2}$, $k = \num{9,0e9}$ N$\cdot$m$^2$\cdot$C^{-2}$. \\
\hline
Champ électrique & $E$ & V/m ou N/C & Force par unité de charge : $\vec{F} = q\vec{E}$. \\
\hline
Tension (d.d.p.) & $U$ & volt (V) & Différence de potentiel ; champ uniforme : $E = \dfrac{U}{d}$. \\
\hline
Travail de la force & $W$ & joule (J) & Énergie transférée : $W_{AB} = q\,U_{AB}$. \\
\hline
\end{tabularx}
\end{center}
\end{aretenir}

\begin{propriete}[La chaîne de raisonnement du chapitre]
Toutes les relations du chapitre s'organisent en une chaîne logique :
\[ \text{charges } q \;\xrightarrow{\;F = k\frac{|q_1 q_2|}{d^2}\;}\; \text{force } \vec{F} \;\xrightarrow{\;\vec{E} = \frac{\vec{F}}{q}\;}\; \text{champ } \vec{E} \;\xrightarrow{\;E = \frac{U}{d}\;}\; \text{tension } U \;\xrightarrow{\;W = qU\;}\; \text{travail } W \]
On peut remonter la chaîne en sens inverse : $U = \dfrac{W}{q}$, $E = \dfrac{U}{d}$, $F = qE$, etc.
\end{propriete}

\begin{attention}[Sens physique de chaque relation]
Chaque flèche a un sens précis qu'il faut savoir formuler en français :
\begin{itemize}
\item $\vec{E} = \vec{F}/q$ : le champ est la force \emph{rapportée à la charge} qui la subit ; le champ existe \textbf{même en l'absence de charge test}.
\item $E = U/d$ : valable \textbf{uniquement} dans un champ uniforme (condensateur plan) ; $d$ est la distance entre les plaques, pas une distance quelconque.
\item $W_{AB} = q\,U_{AB}$ : le signe de $W$ dépend du signe de $q$ \emph{et} du sens du déplacement par rapport au champ.
\end{itemize}
\end{attention}

La carte mentale suivante résume l'ensemble du chapitre. Apprends à la reproduire de mémoire : c'est un excellent plan de révision.

\begin{popfigure}
\begin{center}
\begin{tikzpicture}[
  node distance=2.6cm and 3.4cm,
  box/.style={draw=popblue, thick, rounded corners=6pt, fill=popblueL, minimum width=2.6cm, minimum height=1.1cm, align=center, font=\bfseries},
  rel/.style={-{Stealth[length=3mm]}, thick, popdark},
  lbl/.style={fill=popcream, inner sep=2pt, font=\small, align=center}]
\node[box, align=center] (q) {Charge $q$\\ \small (C)};
\node[box, right=of q, align=center] (F) {Force $\vec{F}$\\ \small (N)};
\node[box, right=of F, align=center] (E) {Champ $\vec{E}$\\ \small (V/m)};
\node[box, below=of E, align=center] (U) {Tension $U$\\ \small (V)};
\node[box, below=of q, align=center] (W) {Travail $W$\\ \small (J)};
\draw[rel] (q) -- (F) node[lbl, midway, above] {$F = k\dfrac{|q_1 q_2|}{d^2}$};
\draw[rel] (F) -- (E) node[lbl, midway, above] {$\vec{E} = \dfrac{\vec{F}}{q}$};
\draw[rel] (E) -- (U) node[lbl, midway, right] {$E = \dfrac{U}{d}$};
\draw[rel] (U) -- (W) node[lbl, midway, above] {$W_{AB} = q\,U_{AB}$};
\draw[rel, poporange] (W) -- (q) node[lbl, midway, left] {$q = n\,e$};
\end{tikzpicture}
\end{center}
\legende{Carte mentale du chapitre : les cinq grandeurs et les formules qui les relient.}
\end{popfigure}

\courssub{Méthode générale de résolution d'un exercice}

Face à un exercice d'électrostatique, le candidat qui réussit n'est pas celui qui connaît le plus de formules, mais celui qui suit une \textbf{démarche ordonnée}. Voici la fiche de résolution type Bac.

\begin{methode}[Fiche de résolution type Bac — 5 étapes]
\begin{enumerate}
\item \textbf{Schéma obligatoire.} Représente les charges (signes indiqués), les points remarquables, les vecteurs $\vec{F}$ et $\vec{E}$ (direction ET sens), un repère si nécessaire. Un schéma clair vaut souvent la moitié des points.
\item \textbf{Données converties.} Liste les données et convertis-les \emph{immédiatement} en unités SI : distances en mètres ($1$ cm $= 10^{-2}$ m), charges en coulombs ($1\,\mu\text{C} = 10^{-6}$ C ; $1\,\text{nC} = 10^{-9}$ C), tensions en volts.
\item \textbf{Choix de la relation adaptée.} Demande-toi : « quelle grandeur je connais, quelle grandeur je cherche ? » puis choisis la flèche de la carte mentale qui les relie (Coulomb, $\vec{F} = q\vec{E}$, $E = U/d$, $W = qU$). Pour plusieurs charges, applique le \cle{principe de superposition} : somme \emph{vectorielle}.
\item \textbf{Calcul littéral puis numérique.} Écris d'abord l'expression littérale, vérifie son homogénéité (analyse dimensionnelle), puis fais l'application numérique avec les puissances de dix bien regroupées.
\item \textbf{Vérification critique.} Contrôle le signe (attraction/répulsion, sens de $\vec{E}$), l'ordre de grandeur (un champ de $10^{15}$ V/m au lycée est suspect !) et le nombre de chiffres significatifs. Rédige une phrase de conclusion avec l'unité.
\end{enumerate}
\end{methode}

\begin{attention}[Les 5 erreurs les plus fréquentes au Bac]
\begin{enumerate}
\item \textbf{Oublier de convertir les distances en mètres} : avec $d$ en cm, la force est fausse d'un facteur $10^{4}$ (car $d$ est au carré !).
\item \textbf{Additionner les normes au lieu des vecteurs} : pour deux charges, $\vec{E} = \vec{E}_1 + \vec{E}_2$ est une somme \emph{vectorielle}. Si les champs sont opposés, ils se soustraient.
\item \textbf{Confondre le sens de $\vec{E}$ et celui de $\vec{F}$ pour une charge négative} : si $q < 0$, $\vec{F}$ et $\vec{E}$ sont de \emph{sens contraires}.
\item \textbf{Utiliser $E = U/d$ hors du champ uniforme} : cette relation n'existe qu'entre les plaques d'un condensateur plan ; pour une charge ponctuelle, $E = k|q|/r^2$.
\item \textbf{Se tromper de signe dans le travail} : $W_{AB} = q(V_A - V_B)$. Un électron ($q = -e$) qui va vers le potentiel le plus élevé a un travail moteur, pas résistant.
\end{enumerate}
\end{attention}

\courssub{Analogie gravitation / électrostatique}

Tu as déjà rencontré une loi en $1/r^2$ : la gravitation. Comparer les deux interactions est un classique du Bac et un excellent moyen de mémoriser.

\begin{center}
\begin{tabularx}{\linewidth}{|l|X|X|}
\hline
\rowcolor{popblueL} \textbf{Critère} & \textbf{Gravitation} & \textbf{Électrostatique} \\
\hline
Loi & $F = G\dfrac{m\,m'}{r^2}$ & $F = k\dfrac{|q\,q'|}{r^2}$ \\
\hline
Constante & $G = \num{6,67e-11}$ SI & $k = \num{9,0e9}$ SI \\
\hline
Grandeur source & masse (toujours positive) & charge (positive \emph{ou} négative) \\
\hline
Sens & toujours attractive & attractive ou répulsive \\
\hline
Champ & $\vec{g} = \vec{F}/m$ & $\vec{E} = \vec{F}/q$ \\
\hline
Ordre de grandeur & très faible à l'échelle atomique & $\approx 10^{36}$ fois plus intense que la gravitation entre proton et électron \\
\hline
\end{tabularx}
\end{center}

\begin{exemplebox}[Gravitation contre électricité dans l'atome d'hydrogène]
Entre le proton et l'électron de l'atome d'hydrogène ($r = \num{5,3e-11}$ m) :
\[ F_{\text{élec}} = k\frac{e^2}{r^2} = \num{9,0e9}\times\frac{(\num{1,6e-19})^2}{(\num{5,3e-11})^2} \approx \num{8,2e-8}\ \text{N} \]
\[ F_{\text{grav}} = G\frac{m_p m_e}{r^2} = \num{6,67e-11}\times\frac{\num{1,67e-27}\times\num{9,11e-31}}{(\num{5,3e-11})^2} \approx \num{3,6e-47}\ \text{N} \]
Le rapport vaut environ $\num{2,3e39}$ : la gravitation est \textbf{totalement négligeable} à l'échelle atomique. C'est pourquoi, dans tout le chapitre, on néglige le poids des particules devant la force électrique.
\end{exemplebox}

\courssub{Entraînement guidé : exercice de synthèse}

\begin{exoresolu}[Charge dans un condensateur plan — synthèse complète]
\textbf{Énoncé.} Deux plaques métalliques planes et parallèles, distantes de $d = 10$ cm, sont reliées à un générateur de tension $U = 500$ V. La plaque $A$ est au potentiel le plus élevé.
\begin{enumerate}
\item Déterminer les caractéristiques du champ électrique $\vec{E}$ entre les plaques.
\item Une petite sphère de charge $q = -2{,}0$ nC est placée en un point $M$ entre les plaques. Calculer la force électrique qu'elle subit et préciser son sens.
\item Calculer le travail de la force électrique lorsque la sphère se déplace de la plaque $A$ vers la plaque $B$.
\end{enumerate}
\tcblower
\textbf{Solution détaillée.}
\textbf{Étape 1 — Schéma et données.} On représente les plaques horizontales, $A$ en haut (potentiel élevé), $\vec{E}$ dirigé de $A$ vers $B$ (du $+$ vers le $-$). Conversions : $d = \num{0,10}$ m ; $q = \num{-2,0e-9}$ C.

\textbf{Étape 2 — Champ.} Champ uniforme :
\[ E = \frac{U}{d} = \frac{500}{\num{0,10}} = \num{5,0e3}\ \text{V/m} \]
Direction : perpendiculaire aux plaques ; sens : de $A$ vers $B$ (vers les potentiels décroissants).

\textbf{Étape 3 — Force.} $F = |q|\,E = \num{2,0e-9}\times\num{5,0e3} = \num{1,0e-5}$ N. Comme $q < 0$, $\vec{F}$ est de sens \emph{opposé} à $\vec{E}$ : la sphère est poussée \textbf{vers la plaque $A$}.

\textbf{Étape 4 — Travail.} $W_{AB} = q(V_A - V_B) = q\,U = \num{-2,0e-9}\times 500 = \num{-1,0e-6}$ J.
Le travail est \textbf{résistant} ($W < 0$) : cohérent, car la sphère se déplace de $A$ vers $B$ alors que la force la ramène vers $A$. Il faudrait un opérateur extérieur pour réaliser ce déplacement.

\textbf{Vérification.} Ordres de grandeur plausibles ; signes cohérents avec le schéma ; homogénéité : $qU$ est bien en C$\times$V $=$ J.
\end{exoresolu}

\begin{exercice}[À toi de jouer]
Deux charges $q_1 = +4{,}0\ \mu$C et $q_2 = -1{,}0\ \mu$C sont placées en $A$ et $B$, distants de $AB = 30$ cm. Détermine le champ électrique résultant au point $M$, milieu de $[AB]$ (direction, sens, norme). \emph{Indication : les deux champs créés en $M$ sont-ils de même sens ou de sens opposés ?}
\end{exercice}

\begin{corrige}[Exercice — champ au milieu de deux charges]
En $M$, $AM = BM = 15$ cm $= \num{0,15}$ m. Le champ créé par $q_1 > 0$ pointe de $A$ vers $B$ (fuyant la charge positive) ; celui créé par $q_2 < 0$ pointe aussi de $A$ vers $B$ (vers la charge négative). Ils s'additionnent :
\[ E = k\frac{|q_1|}{AM^2} + k\frac{|q_2|}{BM^2} = \num{9,0e9}\times\frac{\num{4,0e-6} + \num{1,0e-6}}{(\num{0,15})^2} = \num{2,0e6}\ \text{V/m} \]
$\vec{E}$ est dirigé de $A$ vers $B$, de norme $\num{2,0e6}$ V/m.
\end{corrige}

\begin{savaistu}[Le pouvoir des pointes et les paratonnerres]
Près d'une pointe conductrice, le champ électrique devient considérablement plus intense : c'est le \cle{pouvoir des pointes}. Lors des orages — fréquents et violents dans nos régions, de Douala à Bamenda — le champ entre le nuage chargé et le sol peut atteindre des millions de volts par mètre. Le paratonnerre, tige métallique pointue reliée à la terre, concentre le champ à son extrémité, ionise l'air environnant et « offre » à la foudre un chemin privilégié vers le sol, protégeant ainsi le bâtiment. Inventé par Benjamin Franklin en 1752, il équipe aujourd'hui nos églises, nos écoles et nos immeubles : une application directe et quotidienne de l'électrostatique que tu viens d'étudier !
\end{savaistu}

\begin{aretenir}[L'essentiel de la leçon]
\begin{itemize}
\item La charge est quantifiée ($q = ne$) et se conserve ; l'électrisation se fait par frottement, contact ou influence.
\item Loi de Coulomb : $F = k\dfrac{|q_1 q_2|}{d^2}$, attractive ou répulsive selon les signes.
\item Champ : $\vec{E} = \vec{F}/q$ ; principe de superposition \emph{vectorielle} ; champ uniforme : $E = U/d$.
\item Travail : $W_{AB} = q\,U_{AB}$, indépendant du chemin suivi.
\item Méthode : schéma, conversions SI, relation adaptée, calcul littéral, vérification du signe et de l'ordre de grandeur.
\end{itemize}
\end{aretenir}

\newpage

\coursec{Activité d'intégration}

\courssub{Objectifs de l'activité}

À l'issue de cette activité d'intégration, tu seras capable de :

\begin{itemize}
\item modéliser une situation réelle (un orage au-dessus d'un lycée) par le schéma du \cle{condensateur plan} et justifier l'hypothèse du \cle{champ électrique uniforme} ;
\item calculer un champ électrique à partir d'une tension par la relation $E = \dfrac{U}{d}$ et le comparer à une valeur critique (le \cle{champ disruptif} de l'air) ;
\item calculer la \cle{force électrique} subie par une charge placée dans un champ uniforme ($F = |q|E$) et le \cle{travail} de cette force ($W = qU$ ou $W = F\,d$) ;
\item appliquer la \cle{loi de Coulomb} entre deux charges ponctuelles et discuter la validité du modèle choisi ;
\item rédiger une note technique argumentée, à destination d'un décideur non spécialiste, en t'appuyant sur des résultats chiffrés.
\end{itemize}

\courssub{Situation}

\textbf{Le paratonnerre du lycée : champ et tension dans un orage.}

Le lycée bilingue de Nkongsamba, perché sur une colline des monts Manengouba, est chaque année exposé à de violents orages en saison des pluies. Le préfet a demandé au proviseur de justifier l'installation d'un paratonnerre sur le bâtiment des salles de sciences. Le proviseur, qui n'est pas physicien, sollicite ta classe de Terminale C pour produire une \textbf{note technique d'une demi-page} expliquant, chiffres à l'appui, pourquoi la foudre risque de frapper le bâtiment et comment le paratonnerre protège le lycée.

Pour mener l'étude, on adopte la modélisation suivante. La base d'un cumulonimbus chargé négativement est assimilée à une plaque plane horizontale portant la charge totale $Q = -20~\text{C}$, située à l'altitude $h = 500~\text{m}$ au-dessus du sol. Le sol, conducteur, joue le rôle de la seconde plaque : l'ensemble \{nuage -- sol\} forme un \og condensateur géant \fg{} dans lequel le champ électrique $\vec{E}$ est supposé \textbf{uniforme}, vertical, dirigé du sol vers le nuage.

\begin{popfigure}
\begin{tikzpicture}[scale=0.75, >=Stealth]
% Sol (plaque inférieure)
\draw[thick, fill=popgreenL!30] (-5,-1.5) rectangle (5,-1);
\node[align=center] at (0,-1.25) {Sol conducteur\\$V_{\text{sol}}$ (potentiel élevé)};
% Nuage (plaque supérieure)
\draw[thick, fill=popblueL!30] (-5,4) rectangle (5,4.5);
\node[align=center] at (0,4.25) {Base du nuage ($Q<0$)\\$V_{\text{nuage}}$ (potentiel bas)};
% Lignes de champ uniformes
\foreach \x in {-4,-2,0,2,4} {
    \draw[->, popblue, thick] (\x, 3.5) -- (\x, -0.5);
}
% Légende champ
\node[popblue, rotate=90] at (-4.5, 1.5) {$\vec{E}$ uniforme};
% Bâtiment + Paratonnerre
\draw[thick, fill=poporangeL!30] (3,-1) rectangle (4,1.5);
\draw[thick, popdark] (3.5,1.5) -- (3.5,3); % Mât
\draw[thick, fill=poporange] (3.5,3) circle (0.12); % Pointe
\node[popdark, right] at (3.7, 3) {$q=+1\,\mu\text{C}$};
% Force sur la pointe
\draw[->, popred, very thick] (3.5,3) -- (3.5,3.8) node[above, popred] {$\vec{F}=q\vec{E}$};
% Distance d
\draw[<->, dashed, popmuted] (-5.5,-1) -- (-5.5,4) node[midway, left, popmuted] {$d = h = 500\,\text{m}$};
% Tension
\node[popmuted, left] at (-5.5, 4.5) {$U = 10^8\,\text{V}$};
\end{tikzpicture}
\legende{Modélisation du système \{nuage -- sol\} en condensateur plan : champ uniforme $\vec{E}$, concentration du champ au niveau de la pointe du paratonnerre (effet de pointe).}
\end{popfigure}

\textbf{Données numériques :}

\begin{center}
\begin{tabularx}{\linewidth}{|l|X|}
\hline
\rowcolor{popblueL} Grandeur & Valeur \\
\hline
Tension entre le nuage et le sol & $U = 1{,}0\times 10^{8}~\text{V}$ (100 millions de volts) \\
\hline
Distance nuage -- sol & $d = h = 500~\text{m}$ \\
\hline
Champ disruptif de l'air sec & $E_{\text{dis}} = 3\times 10^{6}~\text{V}\cdot\text{m}^{-1}$ \\
\hline
Charge accumulée sur la pointe du paratonnerre & $q = +1~\mu\text{C} = 10^{-6}~\text{C}$ \\
\hline
Charge induite au sommet d'un immeuble voisin & $Q' = +20~\text{C}$ \\
\hline
Constante de Coulomb & $k = 9{,}0\times 10^{9}~\text{N}\cdot\text{m}^{2}\cdot\text{C}^{-2}$ \\
\hline
\end{tabularx}
\end{center}

\courssub{Consignes}

\begin{enumerate}
\item \textbf{Champ entre le nuage et le sol.} En supposant le champ uniforme, calculer la valeur $E$ du champ électrique régnant entre la base du nuage et le sol. Vérifier l'homogénéité de la relation utilisée.
\item \textbf{Risque de foudre.} Comparer $E$ au champ disruptif $E_{\text{dis}}$ de l'air (champ au-delà duquel l'air devient conducteur). Conclure : l'air peut-il s'ioniser ? Que se passe-t-il alors ?
\item \textbf{Force sur la pointe du paratonnerre.} Calculer la force électrostatique $\vec{F}$ subie par la charge $q = +1~\mu\text{C}$ déposée sur la pointe. Préciser sa direction et son sens sur un schéma.
\item \textbf{Travail de la force.} La charge $q$ est \og arrachée \fg{} de la pointe et entraînée vers le nuage sur une distance $\ell = 10~\text{m}$. Calculer le travail $W$ de la force électrique au cours de ce déplacement. Ce travail est-il moteur ou résistant ? Justifier.
\item \textbf{Modèle de la charge ponctuelle.} En traitant le nuage comme une charge ponctuelle $Q = -20~\text{C}$ située à $500~\text{m}$ d'une charge induite $Q' = +20~\text{C}$ au sommet d'un immeuble, calculer la force de Coulomb entre les deux charges. Comparer l'ordre de grandeur de cette force à celle calculée à la question 3.
\item \textbf{Rédaction.} Rédiger la note technique d'une demi-page destinée au proviseur : elle doit expliquer l'origine du danger, le rôle de la pointe du paratonnerre (concentration des charges et du champ) et recommander l'installation, en citant au moins deux résultats chiffrés.
\end{enumerate}

\courssub{Ressources à mobiliser}

\begin{aretenir}[Boîte à outils de l'activité]
\begin{itemize}
\item \textbf{Champ uniforme et tension :} $E = \dfrac{U}{d}$, avec $U$ en volts (V), $d$ en mètres (m) et $E$ en $\text{V}\cdot\text{m}^{-1}$ (équivalent au $\text{N}\cdot\text{C}^{-1}$).
\item \textbf{Force sur une charge dans un champ :} $\vec{F} = q\vec{E}$, donc $F = |q|\,E$. Le sens de $\vec{F}$ est celui de $\vec{E}$ si $q>0$, opposé si $q<0$.
\item \textbf{Travail de la force électrique dans un champ uniforme :} $W_{AB}(\vec{F}) = q\,E\,\ell$ pour un déplacement $\ell$ dans le sens du champ (travail moteur si $qE\ell > 0$).
\item \textbf{Loi de Coulomb :} $F = k\,\dfrac{|q_1 q_2|}{r^{2}}$ ; force attractive entre charges de signes opposés, répulsive entre charges de même signe.
\item \textbf{Principe de superposition :} le champ créé par plusieurs charges est la somme vectorielle des champs créés par chacune.
\end{itemize}
\end{aretenir}

\courssub{Démarche et corrigé commenté}

\begin{exoresolu}[Le paratonnerre du lycée : corrigé complet]
On modélise le système \{nuage -- sol\} par un condensateur plan géant : $Q = -20~\text{C}$ à $h = 500~\text{m}$, $U = 10^{8}~\text{V}$, $E_{\text{dis}} = 3\times 10^{6}~\text{V}\cdot\text{m}^{-1}$, $q = +1~\mu\text{C}$, $Q' = +20~\text{C}$, $k = 9\times 10^{9}$ SI.

\tcblower

\textbf{Question 1 : champ entre le nuage et le sol.}

Dans un champ uniforme, la tension et le champ sont liés par $E = \dfrac{U}{d}$. Analyse dimensionnelle : $[U]/[d] = \text{V}\cdot\text{m}^{-1}$, qui est bien une unité de champ électrique (équivalente au $\text{N}\cdot\text{C}^{-1}$).
\[
E = \frac{U}{d} = \frac{1{,}0\times 10^{8}}{500} = 2{,}0\times 10^{5}~\text{V}\cdot\text{m}^{-1}.
\]
Le champ est vertical, dirigé du sol (potentiel le plus élevé) vers le nuage chargé négativement (potentiel le plus bas) : $\vec{E}$ est orienté dans le sens de la décroissance du potentiel.

\textbf{Question 2 : risque de foudre.}

Comparons :
\[
E = 2{,}0\times 10^{5}~\text{V}\cdot\text{m}^{-1} \quad \text{et} \quad E_{\text{dis}} = 3\times 10^{6}~\text{V}\cdot\text{m}^{-1}.
\]
On a $E \approx \dfrac{E_{\text{dis}}}{15}$ : le champ moyen est \textbf{quinze fois plus faible} que le champ disruptif. L'air, dans ce modèle simplifié, reste isolant : la foudre ne jaillit pas encore. \emph{Mais} le modèle du champ uniforme sous-estime la réalité : le champ est considérablement renforcé au voisinage des pointes et des arêtes (effet de pointe). C'est précisément pourquoi la foudre frappe les objets élevés et pointus, et pourquoi le paratonnerre fonctionne : sa pointe concentre les lignes de champ, $E$ y dépasse localement $E_{\text{dis}}$, l'air s'ionise et le canal conducteur se forme \emph{à l'endroit choisi}, protégeant le reste du bâtiment.

\textbf{Question 3 : force sur la charge de la pointe.}

La charge $q = +1~\mu\text{C}$ placée dans le champ $\vec{E}$ subit :
\[
F = |q|\,E = 10^{-6}\times 2{,}0\times 10^{5} = 0{,}20~\text{N}.
\]
Comme $q>0$, $\vec{F}$ a le même sens que $\vec{E}$ : verticale, dirigée \textbf{vers le nuage} (vers le haut). Cette force de $0{,}20~\text{N}$ (équivalente au poids d'une masse de $20~\text{g}$) tend littéralement à arracher les charges positives de la pointe.

\textbf{Question 4 : travail de la force sur 10 m.}

Le déplacement de $\ell = 10~\text{m}$ s'effectue dans le sens de la force (vers le nuage) : le travail est \textbf{moteur} :
\[
W = q\,E\,\ell = 10^{-6}\times 2{,}0\times 10^{5}\times 10 = 2{,}0~\text{J}.
\]
Vérification par les unités : $\text{C}\times\text{V}\cdot\text{m}^{-1}\times\text{m} = \text{C}\cdot\text{V} = \text{J}$. Ce travail moteur confirme que la force électrique \emph{accélère} la charge vers le nuage : c'est le début de l'avalanche électronique qui débouche sur la foudre.

\textbf{Question 5 : force de Coulomb entre le nuage et l'immeuble.}

En assimilant le nuage et la charge induite à des charges ponctuelles distantes de $r = 500~\text{m}$ :
\[
F = k\,\frac{|Q\,Q'|}{r^{2}} = 9{,}0\times 10^{9}\times\frac{20\times 20}{(500)^{2}}
= 9{,}0\times 10^{9}\times\frac{400}{2{,}5\times 10^{5}}
= 1{,}44\times 10^{7}~\text{N}.
\]
Soit environ \textbf{14 millions de newtons}, force attractive (charges de signes opposés) — l'équivalent du poids d'une masse de $1\,400$ tonnes ! Cette force gigantesque, près de $10^{8}$ fois plus grande que celle de la question 3, illustre l'ampleur des interactions électrostatiques à l'échelle d'un orage et justifie que l'immeuble non protégé soit en danger. (On notera que le modèle de la charge ponctuelle est grossier à cette échelle : il donne un ordre de grandeur, pas une valeur exacte.)

\textbf{Question 6 : note technique (modèle de rédaction).}

\emph{Note technique à l'attention de M.\ le Proviseur.} Lors d'un orage, la base du nuage (charge d'environ $-20~\text{C}$) et le sol forment un condensateur géant soumis à une tension de l'ordre de $10^{8}~\text{V}$. Le champ électrique moyen atteint $2\times 10^{5}~\text{V/m}$ ; il est certes inférieur au seuil de claquage de l'air ($3\times 10^{6}~\text{V/m}$), mais il est fortement amplifié au voisinage des pointes et des parties saillantes des bâtiments (effet de pointe). Les forces électrostatiques mises en jeu sont considérables : de l'ordre de $10^{7}~\text{N}$ entre le nuage et les charges induites sur un immeuble. Le paratonnerre, tige métallique pointue reliée à la terre, concentre volontairement le champ en sa pointe : l'air s'y ionise en premier et le courant de foudre est canalisé vers le sol sans traverser le bâtiment. \textbf{Nous recommandons l'installation d'un paratonnerre sur le bâtiment des salles de sciences avant la prochaine saison des pluies.}
\end{exoresolu}

\begin{attention}[Pièges rencontrés dans cette activité]
\begin{itemize}
\item Ne pas oublier de convertir les microcoulombs : $1~\mu\text{C} = 10^{-6}~\text{C}$.
\item Dans $E = U/d$, $d$ est la distance \emph{entre les plaques} (500 m), pas la longueur du déplacement de la charge (10 m). Confondre les deux est l'erreur la plus fréquente.
\item Le signe du travail se détermine en comparant le sens de $\vec{F}$ et celui du déplacement, pas seulement le signe de $q$.
\item Dans la loi de Coulomb, la distance est \emph{au carré} : oublier ce carré change le résultat d'un facteur 500 ici.
\end{itemize}
\end{attention}

\courssub{Bilan de l'activité}

Cette activité t'a permis de mobiliser, dans une situation authentique de protection contre la foudre, l'ensemble des savoir-faire de la leçon :

\begin{itemize}
\item la relation $E = \dfrac{U}{d}$ a donné un sens concret au champ uniforme et à la notion de tension : $2\times 10^{5}~\text{V}\cdot\text{m}^{-1}$ entre le nuage et le sol ;
\item la comparaison avec le champ disruptif a montré qu'un modèle physique se juge à l'aune de ses prédictions : le modèle du champ uniforme, insuffisant à lui seul, a conduit à introduire l'\textbf{effet de pointe}, clé du fonctionnement du paratonnerre ;
\item les relations $\vec{F} = q\vec{E}$ et $W = qE\ell$ ont permis de quantifier la force sur une charge ($0{,}20~\text{N}$) et l'énergie mise en jeu ($2{,}0~\text{J}$) ;
\item la loi de Coulomb a révélé l'ordre de grandeur colossal des forces électrostatiques d'un orage ($\sim 10^{7}~\text{N}$), tout en invitant à discuter la validité du modèle de la charge ponctuelle ;
\item enfin, la rédaction de la note technique a développé une compétence transversale essentielle : \textbf{communiquer un raisonnement scientifique chiffré à un non-spécialiste pour éclairer une décision}.
\end{itemize}

\begin{savaistu}[La foudre au Cameroun]
Le Cameroun, avec ses reliefs volcaniques (mont Cameroun, monts Manengouba) et sa longue saison des pluies, connaît une activité orageuse intense : certaines régions de l'Ouest enregistrent plus de 100 jours d'orage par an. Les paratonnerres, inventés par Benjamin Franklin en 1752, protègent aujourd'hui les antennes relais de MTN et Orange, les barrages hydroélectriques comme celui de Lom Pangar, et de plus en plus de lycées et d'hôpitaux. Comprendre le champ électrique, c'est comprendre comment l'homme apprivoise la foudre.
\end{savaistu}

\newpage

\coursec{Méthodes et exercices résolus}

\courssub{Fiche méthode 1 : Réaliser et interpréter une expérience d'électrisation}

\begin{methode}[Analyser une expérience d'électrisation]
\begin{enumerate}
\item \textbf{Identifier le mode d'électrisation} : frottement (transfert d'électrons entre deux corps), contact (transfert entre un corps chargé et un corps neutre), influence (répartition des charges sans contact, à distance).
\item \textbf{Se poser la question des porteurs de charges} : dans les solides, seuls les \cle{électrons} se déplacent ; les noyaux (protons) restent fixes. Un corps chargé positivement a \emph{perdu} des électrons, jamais gagné des protons.
\item \textbf{Appliquer la conservation de la charge} : dans un système isolé, la charge totale se conserve : $q_{\text{avant}} = q_{\text{après}}$.
\item \textbf{Appliquer la quantification} : toute charge est un multiple entier de la charge élémentaire : $q = n\,e$ avec $e = \num{1,6e-19}$ C et $n \in \mathbb{Z}$. Le nombre d'électrons transférés est $n = \dfrac{|q|}{e}$.
\item \textbf{Conclure par le signe} : attraction entre charges de signes opposés, répulsion entre charges de même signe. Attention : un corps neutre peut être \emph{attiré} par un corps chargé (électrisation par influence).
\end{enumerate}
\end{methode}

\begin{exoresolu}[Électrisation d'une règle et pendule électrostatique]
Une règle en plastique frottée avec un chiffon de laine acquiert une charge $q = -\num{3,2e-8}$ C. On l'approche d'une petite boule de sureau suspendue à un fil, initialement neutre : la boule est attirée, touche la règle, puis est repoussée.
\begin{enumerate}
\item Quel est le mode d'électrisation de la règle ? A-t-elle gagné ou perdu des électrons ? Combien ?
\item Expliquer pourquoi la boule neutre est d'abord attirée.
\item Expliquer pourquoi la boule est repoussée après le contact.
\end{enumerate}
\tcblower
\textbf{1. Électrisation de la règle.}

Le frottement contre la laine arrache des électrons à celle-ci : la règle s'électrise \textbf{par frottement}. Sa charge étant négative ($q < 0$), elle a \textbf{gagné} des électrons.

Par quantification de la charge : $q = -n\,e$, d'où
\[
n = \frac{|q|}{e} = \frac{\num{3,2e-8}}{\num{1,6e-19}} = \num{2,0e11}\ \text{électrons}.
\]
Vérification : $n$ est bien un entier (ordre de grandeur $10^{11}$), cohérent avec la quantification.

\textbf{2. Attraction de la boule neutre.}

La règle chargée négativement, approchée de la boule neutre, repousse les électrons libres de la boule vers sa face opposée : la face proche devient positive, la face éloignée négative. C'est l'\textbf{électrisation par influence}. La distance entre la règle et la face positive étant plus petite, la force attractive l'emporte sur la force répulsive : la boule est attirée.

\textbf{3. Répulsion après contact.}

Au contact, des électrons passent de la règle vers la boule (\textbf{électrisation par contact}) : la boule devient chargée négativement, comme la règle. Deux charges de même signe se repoussent : la boule s'écarte.

\textbf{Conclusion :} la règle a gagné $\num{2,0e11}$ électrons par frottement ; la boule neutre est attirée par influence, puis repoussée après s'être chargée négativement par contact.
\end{exoresolu}

\courssub{Fiche méthode 2 : Calculer une force électrostatique (loi de Coulomb)}

\begin{methode}[Appliquer la loi de Coulomb]
\begin{enumerate}
\item \textbf{Schématiser} : placer les charges $q_A$ et $q_B$, la distance $r = AB$, et le vecteur unitaire $\vec{u}_{AB}$ dirigé de $A$ vers $B$.
\item \textbf{Écrire la loi de Coulomb vectoriellement} :
\[
\vec{F}_{A/B} = k\,\frac{q_A q_B}{r^2}\,\vec{u}_{AB}, \qquad k = \num{9,0e9}\ \text{N·m}^2\text{·C}^{-2}.
\]
Le signe du produit $q_A q_B$ donne le sens : $q_A q_B > 0$ : répulsion ; $q_A q_B < 0$ : attraction.
\item \textbf{Convertir les unités avant tout calcul} : distances en mètres (cm $\to$ $10^{-2}$ m), charges en coulombs ($\mu$C $\to$ $10^{-6}$ C, nC $\to$ $10^{-9}$ C).
\item \textbf{Calculer la norme} avec les valeurs absolues : $F = k\,\dfrac{|q_A q_B|}{r^2}$, puis préciser direction et sens sur le schéma.
\item \textbf{Comparer si demandé} à d'autres forces (poids $P = mg$, force gravitationnelle $F_g = G\dfrac{mm'}{r^2}$) en calculant un rapport.
\end{enumerate}
\end{methode}

\begin{exoresolu}[Interaction entre deux billes chargées]
Deux petites billes ponctuelles $A$ et $B$, de masses $m = 0,50$ g chacune, portent les charges $q_A = +\num{4,0e-8}$ C et $q_B = -\num{2,0e-8}$ C. Elles sont distantes de $r = 10$ cm. On donne $k = \num{9,0e9}$ SI et $g = 10$ N/kg.
\begin{enumerate}
\item Représenter les forces électrostatiques et calculer leur intensité commune.
\item Comparer cette force au poids d'une bille. Conclure.
\end{enumerate}
\tcblower
\textbf{1. Forces électrostatiques.}

Schéma : les charges étant de signes opposés, les forces sont \textbf{attractives}, dirigées selon la droite $(AB)$.

\begin{popfigure}
\begin{tikzpicture}[scale=1.1]
\fill[popblue] (0,0) circle (0.09) node[above=4pt] {$A$};
\fill[poporange] (3,0) circle (0.09) node[above=4pt] {$B$};
\node[below=4pt, popblue] at (0,0) {$q_A>0$};
\node[below=4pt, poporange] at (3,0) {$q_B<0$};
\draw[-{Stealth}, very thick, popblue] (0.15,0) -- (1.5,0) node[midway, above] {$\vec{F}_{B/A}$};
\draw[-{Stealth}, very thick, poporange] (2.85,0) -- (1.6,0) node[midway, above] {};
\draw[<->, popmuted] (0,-0.7) -- (3,-0.7) node[midway, below] {$r = \SI{10}{cm}$};
\end{tikzpicture}
\legende{Forces attractives entre deux charges de signes opposés.}
\end{popfigure}

Conversion : $r = \SI{10}{cm} = \SI{0,10}{m}$. Intensité :
\[
F = k\,\frac{|q_A q_B|}{r^2} = \num{9,0e9}\times\frac{\num{4,0e-8}\times\num{2,0e-8}}{(\num{0,10})^2}.
\]
Calcul : numérateur $= \num{9,0e9}\times\num{8,0e-16} = \num{7,2e-6}$ ; dénominateur $= \num{1,0e-2}$. Donc
\[
\boxed{F = \SI{7,2e-4}{N}}.
\]
D'après le principe des actions réciproques, $\vec{F}_{A/B} = -\vec{F}_{B/A}$ : les deux forces ont même norme.

\textbf{2. Comparaison au poids.}

\[
P = mg = \num{0,50e-3}\times 10 = \SI{5,0e-3}{N}, \qquad \frac{P}{F} = \frac{\num{5,0e-3}}{\num{7,2e-4}} \approx 7.
\]

\textbf{Conclusion :} la force électrostatique vaut $\SI{7,2e-4}{N}$ ; elle est du même ordre de grandeur que le poids (environ 7 fois plus faible), ce qui explique qu'elle puisse faire dévier ou soulever de petits objets légers.
\end{exoresolu}

\courssub{Fiche méthode 3 : Déterminer le champ créé par une charge ponctuelle}

\begin{methode}[Champ électrostatique d'une charge ponctuelle]
\begin{enumerate}
\item \textbf{Écrire l'expression vectorielle} du champ créé en $M$ par la charge $q$ placée en $O$ :
\[
\vec{E}(M) = k\,\frac{q}{r^2}\,\vec{u}, \qquad r = OM,
\]
où $\vec{u}$ est le vecteur unitaire dirigé de $O$ vers $M$.
\item \textbf{Raisonner sur le sens} : si $q > 0$, $\vec{E}$ est \emph{centrifuge} (pointe vers l'extérieur) ; si $q < 0$, $\vec{E}$ est \emph{centripète} (pointe vers la charge).
\item \textbf{Calculer la norme} : $E = k\,\dfrac{|q|}{r^2}$, en volts par mètre (V/m, équivalent à N/C).
\item \textbf{Exploiter la relation fondamentale} : une charge $q_0$ placée en $M$ subit $\vec{F} = q_0\,\vec{E}$. Réciproquement, si l'on connaît $\vec{F}$ et $q_0$, alors $\vec{E} = \vec{F}/q_0$.
\item \textbf{Décrire les lignes de champ} : radiales, sortantes pour $q>0$, entrantes pour $q<0$ ; elles ne se coupent jamais.
\end{enumerate}
\end{methode}

\begin{exoresolu}[Champ créé par un ion et force subie par un électron]
Un ion aluminium $\ce{Al^{3+}}$, assimilé à une charge ponctuelle $q = +\num{4,8e-19}$ C, est placé en $O$. Un point $M$ est situé à $r = 3,0$ cm de $O$.
\begin{enumerate}
\item Déterminer les caractéristiques du champ $\vec{E}$ créé en $M$.
\item On place en $M$ un électron (charge $-e = -\num{1,6e-19}$ C). Déterminer la force qu'il subit.
\end{enumerate}
\tcblower
\textbf{1. Champ en $M$.}

La charge $q$ est positive : le champ est \textbf{centrifuge}, dirigé de $O$ vers $M$ (radial, sortant). Sa norme :
\[
E = k\,\frac{|q|}{r^2} = \num{9,0e9}\times\frac{\num{4,8e-19}}{(\num{0,030})^2}
= \num{9,0e9}\times\frac{\num{4,8e-19}}{\num{9,0e-4}}.
\]
\[
E = \num{9,0e9}\times\num{5,33e-16} \approx \boxed{\SI{4,8e-6}{V/m}}.
\]
Analyse dimensionnelle : $k q / r^2$ s'exprime en N·C$^{-2}$·C·m$^{-2}$ = N/C = V/m. Correct.

\textbf{2. Force sur l'électron.}

\[
\vec{F} = q_0\,\vec{E} = -e\,\vec{E}.
\]
Comme $q_0 < 0$, $\vec{F}$ est \textbf{opposée à} $\vec{E}$ : dirigée de $M$ vers $O$ (attraction vers l'ion positif). Norme :
\[
F = e\,E = \num{1,6e-19}\times\num{4,8e-6} \approx \boxed{\SI{7,7e-25}{N}}.
\]

\textbf{Conclusion :} le champ en $M$ vaut $\SI{4,8e-6}{V/m}$, radial sortant ; l'électron subit une force attractive de norme $\SI{7,7e-25}{N}$, colinéaire et de sens opposé à $\vec{E}$.
\end{exoresolu}

\courssub{Fiche méthode 4 : Superposer les champs de plusieurs charges}

\begin{methode}[Principe de superposition]
\begin{enumerate}
\item \textbf{Énoncer le principe} : le champ total en $M$ est la somme vectorielle des champs créés par chaque charge :
\[
\vec{E}(M) = \vec{E}_1(M) + \vec{E}_2(M) + \cdots
\]
\item \textbf{Dessiner séparément} chaque champ $\vec{E}_i$ en $M$ : direction = droite joignant la charge à $M$ ; sens selon le signe de la charge.
\item \textbf{Exploiter les symétries} : pour deux charges égales, les composantes perpendiculaires à la médiatrice (ou à l'axe) s'annulent souvent — cela divise le travail par deux.
\item \textbf{Projeter sur des axes} et additionner les composantes :
\[
E_x = \sum_i E_{i,x}, \qquad E_y = \sum_i E_{i,y}, \qquad E = \sqrt{E_x^2 + E_y^2}.
\]
\item \textbf{Pour trouver un point de champ nul}, écrire $E_1 = E_2$ en normes et raisonner sur les positions possibles (les champs doivent être opposés : point situé sur la droite des charges).
\end{enumerate}
\end{methode}

\begin{exoresolu}[Champ au sommet d'un triangle équilatéral]
Deux charges $q_A = q_B = +\num{2,0e-8}$ C sont placées aux sommets $A$ et $B$ d'un triangle équilatéral $ABC$ de côté $a = 6,0$ cm. Déterminer le champ électrostatique total au sommet $C$ (direction, sens, norme).
\tcblower
\textbf{Mise en place.} Les deux charges sont positives : $\vec{E}_A$ est dirigé selon $(AC)$ en s'éloignant de $A$, $\vec{E}_B$ selon $(BC)$ en s'éloignant de $B$. Chaque champ fait un angle de $60°$ avec l'autre.

\begin{popfigure}
\begin{tikzpicture}[scale=1.2]
\coordinate (A) at (0,0);
\coordinate (B) at (3,0);
\coordinate (C) at (1.5,2.598);
\fill[popblue] (A) circle (0.07) node[below left] {$A\ (+q)$};
\fill[popblue] (B) circle (0.07) node[below right] {$B\ (+q)$};
\fill[popdark] (C) circle (0.05) node[above] {$C$};
\draw[popmuted] (A)--(B)--(C)--cycle;
\draw[-{Stealth}, thick, poporange] (C) -- ++(120:1.3) node[above left] {$\vec{E}_A$};
\draw[-{Stealth}, thick, popgreen] (C) -- ++(60:1.3) node[above right] {$\vec{E}_B$};
\draw[-{Stealth}, very thick, poppink] (C) -- ++(90:1.3) node[right] {$\vec{E}$};
\draw[dashed, popmuted] (C) -- (1.5,0);
\end{tikzpicture}
\legende{Superposition des champs en $C$ : par symétrie, $\vec{E}$ est porté par la médiatrice de $[AB]$.}
\end{popfigure}

\textbf{Normes des champs partiels.} Comme $CA = CB = a$ et $q_A = q_B$ :
\[
E_A = E_B = k\,\frac{q}{a^2} = \num{9,0e9}\times\frac{\num{2,0e-8}}{(\num{0,060})^2}
= \num{9,0e9}\times\frac{\num{2,0e-8}}{\num{3,6e-3}} = \SI{5,0e4}{V/m}.
\]

\textbf{Projection.} Par symétrie, les composantes horizontales de $\vec{E}_A$ et $\vec{E}_B$ s'annulent. Les composantes verticales s'ajoutent ; chacune vaut $E_A\cos 30°$ (angle entre $\vec{E}_A$ et la verticale : $30°$) :
\[
E = 2\,E_A\cos 30° = 2\times\num{5,0e4}\times\frac{\sqrt{3}}{2} = \num{5,0e4}\sqrt{3}.
\]
\[
\boxed{E \approx \SI{8,7e4}{V/m}}.
\]

\textbf{Conclusion :} le champ total en $C$ est porté par la médiatrice de $[AB]$, dirigé vers le haut (en s'éloignant des charges positives), de norme $E \approx \SI{8,7e4}{V/m}$.
\end{exoresolu}

\courssub{Fiche méthode 5 : Exploiter un champ uniforme et le condensateur plan}

\begin{methode}[Champ uniforme entre deux plaques]
\begin{enumerate}
\item \textbf{Identifier la situation} : deux plaques planes, parallèles, de grandes dimensions, soumises à une tension $U = V_A - V_B$, distantes de $d$ $\Rightarrow$ champ \cle{uniforme} entre les plaques.
\item \textbf{Appliquer la relation fondamentale} :
\[
E = \frac{U}{d} \qquad \text{avec } U \text{ en volts (V), } d \text{ en mètres (m), } E \text{ en V/m}.
\]
\item \textbf{Orienter correctement} : $\vec{E}$ est perpendiculaire aux plaques, dirigé de la plaque de \emph{plus haut potentiel} vers celle de \emph{plus bas potentiel} (dans le sens des potentiels décroissants). Les lignes de champ sont parallèles et équidistantes.
\item \textbf{Relier champ et force} : une charge $q$ placée entre les plaques subit $\vec{F} = q\vec{E}$, de norme constante $F = |q|\,E = |q|\,\dfrac{U}{d}$.
\item \textbf{Attention aux conversions} : $d$ est souvent donnée en cm ou mm — convertir impérativement en mètres.
\end{enumerate}
\end{methode}

\begin{exoresolu}[Condensateur plan d'un dispositif de déviation]
Deux plaques métalliques planes et parallèles $P_1$ et $P_2$, distantes de $d = 5,0$ cm, sont reliées à un générateur qui impose la tension $U = V_{P_1} - V_{P_2} = +\num{2,0e3}$ V.
\begin{enumerate}
\item Déterminer les caractéristiques du champ $\vec{E}$ entre les plaques.
\item Un proton (charge $+e = \num{1,6e-19}$ C, masse $m = \num{1,67e-27}$ kg) est placé entre les plaques. Comparer la force électrique subie à son poids ($g = 10$ N/kg).
\end{enumerate}
\tcblower
\textbf{1. Champ entre les plaques.}

Le champ est uniforme, perpendiculaire aux plaques. Comme $V_{P_1} > V_{P_2}$, $\vec{E}$ est dirigé de $P_1$ vers $P_2$ (sens des potentiels décroissants). Norme :
\[
E = \frac{U}{d} = \frac{\num{2,0e3}}{\num{5,0e-2}} = \boxed{\SI{4,0e4}{V/m}}.
\]

\textbf{2. Force sur le proton et comparaison au poids.}

Le proton porte une charge positive : $\vec{F} = e\vec{E}$ a même sens que $\vec{E}$, donc dirigée de $P_1$ vers $P_2$.
\[
F = e\,E = \num{1,6e-19}\times\num{4,0e4} = \SI{6,4e-15}{N}.
\]
Poids du proton :
\[
P = mg = \num{1,67e-27}\times 10 = \SI{1,67e-26}{N}.
\]
Rapport :
\[
\frac{F}{P} = \frac{\num{6,4e-15}}{\num{1,67e-26}} \approx \num{3,8e11}.
\]

\textbf{Conclusion :} le champ vaut $\SI{4,0e4}{V/m}$, dirigé de $P_1$ vers $P_2$ ; la force électrique sur le proton ($\SI{6,4e-15}{N}$) est environ $10^{11}$ fois plus grande que son poids : \textbf{le poids est toujours négligeable devant la force électrique pour les particules élémentaires}.
\end{exoresolu}

\courssub{Fiche méthode 6 : Travail de la force électrique, potentiel et d.d.p.}

\begin{methode}[Travail, potentiel et énergie dans un champ uniforme]
\begin{enumerate}
\item \textbf{Retenir la définition du potentiel} : la d.d.p. entre deux points est liée au champ par
\[
U_{AB} = V_A - V_B = E\,d
\]
où $d$ est la distance mesurée \emph{le long des lignes de champ} entre les équipotentielles passant par $A$ et $B$.
\item \textbf{Calculer le travail de la force électrique} sur une charge $q$ allant de $A$ à $B$ :
\[
W_{AB}(\vec{F}) = q\,(V_A - V_B) = q\,U_{AB}.
\]
Le travail ne dépend pas du chemin suivi (force conservative).
\item \textbf{Respecter les signes} : $W > 0$ : travail moteur ; $W < 0$ : travail résistant. Le signe dépend à la fois de $q$ et de $U_{AB}$.
\item \textbf{Utiliser le théorème de l'énergie cinétique} si la particule est en mouvement :
\[
\Delta E_c = \tfrac{1}{2}mv_B^2 - \tfrac{1}{2}mv_A^2 = q\,U_{AB}
\]
(le poids étant négligeable). Une charge positive est accélérée vers les potentiels décroissants ; une charge négative vers les potentiels croissants.
\item \textbf{Vérifier les unités} : $W$ en joules (J), $q$ en C, $U$ en V ; $1\ \text{eV} = \num{1,6e-19}$ J.
\end{enumerate}
\end{methode}

\begin{exoresolu}[Accélération d'un électron dans un canon à électrons]
Dans un oscilloscope, un électron (charge $-e = -\num{1,6e-19}$ C, masse $m = \num{9,11e-31}$ kg) quitte la cathode $C$ avec une vitesse négligeable et arrive sur l'anode $A$. Entre $C$ et $A$ règne un champ uniforme, la tension étant $U_{AC} = V_A - V_C = +\num{1,0e3}$ V.
\begin{enumerate}
\item Calculer le travail de la force électrique subie par l'électron entre $C$ et $A$.
\item En déduire la vitesse d'arrivée de l'électron sur l'anode. Exprimer aussi son énergie en électronvolts.
\end{enumerate}
\tcblower
\textbf{1. Travail de la force électrique.}

\[
W_{CA}(\vec{F}) = q\,(V_C - V_A) = (-e)\times(-U_{AC}) = e\,U_{AC}.
\]
\[
W_{CA} = \num{1,6e-19}\times\num{1,0e3} = \boxed{\SI{1,6e-16}{J}}.
\]
Le travail est positif (moteur) : l'électron, charge négative, est bien attiré vers l'anode de potentiel plus élevé.

\textbf{2. Vitesse d'arrivée.}

Le poids de l'électron ($P = mg \approx \SI{9,11e-30}{N}$) est négligeable devant la force électrique. Théorème de l'énergie cinétique entre $C$ et $A$ :
\[
\frac{1}{2}mv_A^2 - \frac{1}{2}mv_C^2 = W_{CA}.
\]
Avec $v_C \approx 0$ :
\[
v_A = \sqrt{\frac{2\,W_{CA}}{m}} = \sqrt{\frac{2\times\num{1,6e-16}}{\num{9,11e-31}}}
= \sqrt{\num{3,5126e14}} \approx \boxed{\SI{1,87e7}{m/s}}.
\]
Analyse dimensionnelle : $\sqrt{\text{J/kg}} = \sqrt{\text{m}^2\text{·s}^{-2}} = $ m/s. Correct.

En électronvolts : l'électron franchit une d.d.p. de 1000 V, donc
\[
E_c = e\times 1000\ \text{V} = 1000\ \text{eV} = 1,0\ \text{keV}.
\]

\textbf{Conclusion :} le travail de la force électrique vaut $\SI{1,6e-16}{J}$ ; l'électron atteint l'anode avec la vitesse $v_A \approx \SI{1,87e7}{m/s}$ et une énergie cinétique de 1,0 keV. C'est le principe du canon à électrons des oscilloscopes et des anciens téléviseurs à tube.
\end{exoresolu}

\courssub{Synthèse et méthodes de résolution}

\begin{aretenir}[Formulaires clés du chapitre]
\begin{itemize}
\item \textbf{Quantification de la charge :} $q = n e \quad (e = \SI{1,6e-19}{C})$.
\item \textbf{Loi de Coulomb (vectorielle) :} $\vec{F}_{A/B} = k \dfrac{q_A q_B}{r^2} \vec{u}_{AB} \quad (k = \SI{9,0e9}{N.m^2.C^{-2}})$.
\item \textbf{Champ d'une charge ponctuelle :} $\vec{E}(M) = k \dfrac{q}{r^2} \vec{u}$ (centrifuge si $q>0$, centripète si $q<0$).
\item \textbf{Principe de superposition :} $\vec{E} = \sum \vec{E}_i$ (somme \textbf{vectorielle}).
\item \textbf{Champ uniforme (condensateur) :} $E = \dfrac{U}{d}$, $\vec{E} \perp$ aux plaques, sens $V_{\text{haut}} \to V_{\text{bas}}$.
\item \textbf{Force électrique :} $\vec{F} = q \vec{E}$.
\item \textbf{Travail de la force électrique :} $W_{AB} = q (V_A - V_B) = q U_{AB}$.
\item \textbf{Théorème de l'énergie cinétique (particule chargée) :} $\Delta E_c = q U_{AB}$ (poids négligeable pour $e^-$, $p^+$, ions).
\end{itemize}
\end{aretenir}

\begin{methode}[Check-list « Bac » pour un exercice d'électrostatique]
\begin{enumerate}
\item \textbf{Unités :} Toute distance en \textbf{mètres}, toute charge en \textbf{coulombs}. Écrire la conversion explicitement.
\item \textbf{Vecteurs :} Dessiner $\vec{E}_i$ et $\vec{F}_i$ avant de calculer. Ne jamais additionner des normes sauf colinéarité prouvée.
\item \textbf{Sens :} $q>0 \Rightarrow \vec{E}$ sortante, $\vec{F}$ même sens que $\vec{E}$ ; $q<0 \Rightarrow \vec{E}$ entrante, $\vec{F}$ sens opposé à $\vec{E}$. Entre plaques : $\vec{E}$ va du $+$ vers le $-$.
\item \textbf{Signes du travail :} $W_{AB} = q U_{AB}$ avec $U_{AB} = V_A - V_B$. Vérifier que le signe de $W$ correspond à l'intuition (moteur/résistant).
\item \textbf{Négligeabilité du poids :} Pour une particule élémentaire, calculer $F/P$ et écrire « $F \gg P$ donc poids négligeable ». Pour un objet macroscopique, comparer avant de conclure.
\end{enumerate}
\end{methode}

\begin{attention}[Les 5 erreurs qui coûtent des points au Bac]
\begin{enumerate}
\item \textbf{Oublier les conversions d'unités.} La distance en cm non convertie dans $F = k|q_1 q_2|/r^2$ fausse le résultat d'un facteur $10^4$ ! De même, $d$ dans $E = U/d$ doit être en mètres. Écris systématiquement la conversion sur ta copie.
\item \textbf{Additionner les normes au lieu des vecteurs.} Pour plusieurs charges, le champ total est une somme \emph{vectorielle} : $\vec{E} = \vec{E}_1 + \vec{E}_2$. Écrire $E = E_1 + E_2$ sans justification (colinéarité, symétrie) est une erreur de raisonnement sanctionnée.
\item \textbf{Se tromper de sens pour $\vec{E}$ et $\vec{F}$.} $\vec{E}$ pointe vers l'extérieur pour $q>0$, vers la charge pour $q<0$ ; entre deux plaques, $\vec{E}$ va du \emph{plus haut} vers le \emph{plus bas} potentiel. Et $\vec{F} = q\vec{E}$ : si $q<0$, la force est \emph{opposée} au champ. Un schéma avec flèches évite tout.
\item \textbf{Mal gérer les signes dans le travail.} $W_{AB} = q(V_A - V_B) = q\,U_{AB}$ : attention à l'ordre des indices ! Confondre $U_{AB}$ et $U_{BA}$ change le signe du travail, donc la conclusion (moteur/résistant) et le signe de $\Delta E_c$.
\item \textbf{Conserver le poids d'une particule élémentaire.} Pour un électron, un proton ou un ion, $F/P \sim 10^{10}$ : le poids est négligeable et il faut le \emph{justifier quantitativement} (calcul du rapport), pas l'affirmer gratuitement. À l'inverse, pour une bille de laboratoire, comparer $F$ et $P$ avant de conclure.
\end{enumerate}
\end{attention}

\newpage

\coursec{Exercices}

\courssub{Je vérifie mes connaissances}

\begin{exercice}[QCM : une seule bonne réponse]
Pour chaque question, choisir la bonne réponse et justifier brièvement.
\begin{enumerate}
\item Deux corps frottés l'un contre l'autre s'électrisent. Le verre frotté avec la laine se charge positivement car :
\begin{enumerate}
\item il a gagné des protons ;
\item il a perdu des électrons ;
\item il a gagné des électrons ;
\item il a perdu des neutrons.
\end{enumerate}
\item La force d'interaction entre deux charges ponctuelles $q_1$ et $q_2$ distantes de $r$ est proportionnelle à :
\begin{enumerate}
\item $r$ ; \item $r^2$ ; \item $\dfrac{1}{r}$ ; \item $\dfrac{1}{r^2}$.
\end{enumerate}
\item Si l'on double la distance entre deux charges ponctuelles, la norme de la force électrostatique est :
\begin{enumerate}
\item divisée par 2 ; \item multipliée par 2 ; \item divisée par 4 ; \item inchangée.
\end{enumerate}
\item Entre les armatures d'un condensateur plan soumis à une tension $U = \SI{200}{V}$ et distantes de $d = \SI{5,0}{cm}$, le champ électrique vaut :
\begin{enumerate}
\item $E = \SI{10}{V.m^{-1}}$ ; \item $E = \SI{4000}{V.m^{-1}}$ ; \item $E = \SI{1000}{V.m^{-1}}$ ; \item $E = \SI{40}{V.m^{-1}}$.
\end{enumerate}
\item Les lignes de champ électrostatique :
\begin{enumerate}
\item sortent des charges négatives et entrent dans les charges positives ;
\item sortent des charges positives et entrent dans les charges négatives ;
\item peuvent se couper en un point ;
\item sont toujours rectilignes.
\end{enumerate}
\end{enumerate}
\end{exercice}

\begin{exercice}[Vrai ou faux, avec justification]
Répondre par vrai ou faux en justifiant chaque réponse.
\begin{enumerate}
\item La charge électrique d'un corps peut prendre n'importe quelle valeur réelle.
\item Deux charges de signes contraires se repoussent.
\item Le champ électrostatique créé en un point par plusieurs charges est la somme vectorielle des champs créés par chacune des charges.
\item Dans un champ électrique uniforme, le travail de la force électrique dépend du chemin suivi entre deux points.
\item La force électrostatique subie par une charge négative placée dans un champ $\vec{E}$ est de même sens que $\vec{E}$.
\end{enumerate}
\end{exercice}

\begin{exercice}[Définitions et restitution de cours]
\begin{enumerate}
\item Définir : électrisation par influence ; charge électrique élémentaire ; ligne de champ électrostatique.
\item Énoncer la loi de Coulomb en précisant la signification et l'unité SI de chaque grandeur.
\item Donner l'expression vectorielle du champ électrostatique créé par une charge ponctuelle $q$ placée en $O$ en un point $M$.
\item Rappeler la relation entre le champ électrique uniforme $E$, la tension $U$ et la distance $d$ entre les armatures d'un condensateur plan. Vérifier l'homogénéité de cette relation par analyse dimensionnelle.
\end{enumerate}
\end{exercice}

\begin{exercice}[Calculs directs]
\begin{enumerate}
\item Calculer la norme de la force électrostatique entre deux charges $q_1 = +\SI{2,0}{\micro C}$ et $q_2 = -\SI{5,0}{\micro C}$ distantes de $r = \SI{10}{cm}$. Préciser si elle est attractive ou répulsive. On donne $k = \SI{9,0e9}{N.m^2.C^{-2}}$.
\item Une charge $q = +\SI{4,0}{nC}$ est placée en un point $O$. Calculer la norme du champ électrostatique qu'elle crée en un point $M$ situé à $d = \SI{20}{cm}$ de $O$, et préciser sa direction et son sens.
\item Un électron ($q = -e = -\SI{1,6e-19}{C}$) est placé dans un champ uniforme $E = \SI{2,0e3}{V.m^{-1}}$. Calculer la norme de la force électrique qu'il subit et la comparer à son poids ($m_e = \SI{9,1e-31}{kg}$, $g = \SI{9,8}{N.kg^{-1}}$). Conclure.
\end{enumerate}
\end{exercice}

\courssub{Je m'entraîne}

\begin{exercice}[Pendule électrostatique]
Une petite sphère de masse $m = \SI{0,50}{g}$, portant une charge $q$, est suspendue à un fil inextensible de longueur $\ell = \SI{20}{cm}$. On approche de la sphère une baguette chargée qui crée au niveau de la sphère un champ horizontal uniforme de norme $E = \SI{3,0e4}{V.m^{-1}}$. À l'équilibre, le fil fait un angle $\alpha = \SI{15}{\degree}$ avec la verticale.
\begin{enumerate}
\item Faire le bilan des forces appliquées à la sphère et représenter ces forces sur un schéma.
\item Déterminer la valeur algébrique de la charge $q$ (on précisera son signe sachant que la sphère est attirée vers la baguette positive).
\item On donne $g = \SI{9,8}{N.kg^{-1}}$.
\end{enumerate}
\end{exercice}

\begin{exercice}[Champ créé par deux charges]
Deux charges ponctuelles $q_A = +\SI{3,0}{\micro C}$ et $q_B = -\SI{1,0}{\micro C}$ sont placées en deux points $A$ et $B$ distants de $AB = \SI{30}{cm}$.
\begin{enumerate}
\item Représenter sur un schéma les vecteurs champs $\vec{E}_A$ et $\vec{E}_B$ créés au point $M$, milieu du segment $[AB]$.
\item Calculer les normes $E_A$ et $E_B$, puis la norme du champ résultant $\vec{E}$ en $M$. Préciser sa direction et son sens.
\item On place en $M$ une charge $q = -\SI{2,0}{nC}$. Déterminer les caractéristiques (direction, sens, norme) de la force électrostatique qu'elle subit.
\end{enumerate}
On donne $k = \SI{9,0e9}{N.m^2.C^{-2}}$.
\end{exercice}

\begin{exercice}[Condensateur plan et travail de la force électrique]
Un condensateur plan a ses armatures distantes de $d = \SI{2,0}{cm}$ et est soumis à une tension $U = \SI{400}{V}$.
\begin{enumerate}
\item Calculer la norme du champ électrique uniforme régnant entre les armatures.
\item Un électron se déplace de la plaque négative vers la plaque positive. Calculer le travail de la force électrique au cours de ce déplacement. Ce travail est-il moteur ou résistant ?
\item En admettant que l'électron part de la plaque négative avec une vitesse nulle et que seule la force électrique travaille, déterminer sa vitesse d'arrivée sur la plaque positive.
\end{enumerate}
On donne : $e = \SI{1,6e-19}{C}$ ; $m_e = \SI{9,1e-31}{kg}$.
\end{exercice}

\begin{exercice}[Lecture d'une carte de lignes de champ]
On donne la carte des lignes de champ électrostatique créées par deux charges ponctuelles $q_1$ et $q_2$ placées en deux points $A$ et $B$ : 12 lignes partent de $A$ et 4 lignes arrivent en $B$ ; les lignes sont resserrées au voisinage de $A$.
\begin{enumerate}
\item Déterminer les signes des charges $q_1$ et $q_2$. Justifier.
\item Comparer les valeurs absolues de $q_1$ et $q_2$ et donner la valeur du rapport $\dfrac{|q_1|}{|q_2|}$.
\item Indiquer, en justifiant, la zone de la carte où le champ est le plus intense.
\item Existe-t-il un point de champ nul entre $A$ et $B$ ? Justifier la réponse.
\end{enumerate}
\end{exercice}

\courssub{Je me prépare au Bac}

\begin{exercice}[Type Bac : champ et force entre deux charges]
Deux charges ponctuelles $q_A = +\SI{3,0}{\micro C}$ et $q_B = -\SI{1,0}{\micro C}$ sont fixées en deux points $A$ et $B$ d'un axe horizontal, avec $AB = \SI{30}{cm}$. On note $M$ le milieu du segment $[AB]$. On donne $k = \SI{9,0e9}{N.m^2.C^{-2}}$.
\begin{enumerate}
\item \textbf{Étude du champ en $M$.}
\begin{enumerate}
\item Donner l'expression vectorielle du champ électrostatique créé par une charge ponctuelle, puis celle du champ résultant en $M$ (principe de superposition).
\item Représenter sur un schéma à l'échelle les vecteurs $\vec{E}_A$ et $\vec{E}_B$ en $M$.
\item Calculer $E_A$ et $E_B$, puis déterminer les caractéristiques complètes du champ résultant $\vec{E}(M)$.
\end{enumerate}
\item \textbf{Force subie par une charge placée en $M$.} On place en $M$ une charge $q = -\SI{2,0}{nC}$.
\begin{enumerate}
\item Donner l'expression de la force électrostatique $\vec{F}$ subie par $q$ en fonction de $\vec{E}(M)$.
\item Calculer la norme de $\vec{F}$ et préciser sa direction et son sens.
\item Retrouver ce résultat en calculant directement les forces de Coulomb exercées par $q_A$ et $q_B$ sur $q$.
\end{enumerate}
\item \textbf{Point de champ nul.} Montrer qu'il existe un point $P$ de l'axe $(AB)$ où le champ électrostatique total est nul. Préciser, en le justifiant, si ce point se trouve entre $A$ et $B$, à gauche de $A$ ou à droite de $B$, puis calculer la distance $BP$.
\end{enumerate}
\end{exercice}

\begin{exercice}[Type Bac : trois charges aux sommets d'un triangle équilatéral]
Trois charges ponctuelles identiques $q = +\SI{2,0}{\micro C}$ sont placées aux sommets $A$, $B$ et $C$ d'un triangle équilatéral de côté $a = \SI{30}{cm}$. On note $G$ le centre de gravité du triangle. On donne $k = \SI{9,0e9}{N.m^2.C^{-2}}$.
\begin{enumerate}
\item Montrer que la distance entre un sommet et le point $G$ vaut $r = \dfrac{a}{\sqrt{3}}$. Calculer sa valeur numérique.
\item Calculer la norme du champ $\vec{E}_A$ créé par la charge placée en $A$ au point $G$. Que valent $E_B$ et $E_C$ ?
\item Construire sur un schéma soigné les trois vecteurs $\vec{E}_A$, $\vec{E}_B$ et $\vec{E}_C$ en $G$.
\item En exploitant les propriétés de symétrie de la figure, montrer que le champ résultant en $G$ est nul.
\item On remplace la charge en $C$ par une charge $q' = -\SI{2,0}{\micro C}$. Déterminer alors les caractéristiques (direction, sens, norme) du champ résultant en $G$.
\item On place en $G$ un électron ($q = -e$, $e = \SI{1,6e-19}{C}$) dans la situation de la question 5. Calculer la norme de la force qu'il subit.
\end{enumerate}
\end{exercice}

\begin{exercice}[Type Bac : expérience de Millikan et quantification de la charge]
Entre les armatures horizontales d'un condensateur plan distantes de $d = \SI{1,0}{cm}$, on établit une tension $U = \SI{500}{V}$, l'armature supérieure étant au potentiel le plus élevé. Une gouttelette d'huile de masse $m = \SI{1,3e-14}{kg}$, chargée négativement, est introduite entre les armatures. Pour la valeur $U$ ci-dessus, elle reste immobile. On donne $g = \SI{9,8}{N.kg^{-1}}$ et $e = \SI{1,6e-19}{C}$.
\begin{enumerate}
\item Représenter le condensateur, le sens du champ $\vec{E}$ entre les armatures, et faire le bilan des forces appliquées à la gouttelette.
\item Calculer la norme du champ électrique $E$ régnant entre les armatures.
\item Écrire la condition d'équilibre de la gouttelette et en déduire la valeur absolue $|q|$ de sa charge. Justifier le signe de $q$.
\item Calculer le rapport $n = \dfrac{|q|}{e}$. Que représente $n$ ? Enoncer la loi de quantification de la charge électrique et conclure.
\item La tension est maintenant portée à $U' = \SI{600}{V}$. Préciser, en le justifiant, le sens du mouvement de la gouttelette, puis calculer la norme de la force résultante qu'elle subit.
\item Calculer le travail de la force électrique lorsque la gouttelette se déplace de la plaque négative vers la plaque positive sous la tension $U'$. Ce travail est-il moteur ou résistant ?
\end{enumerate}
\end{exercice}

\newpage

\coursec{Corrigés des exercices}

\begin{corrige}[Exercice 1 — QCM : une seule bonne réponse]
\begin{enumerate}
\item \textbf{Bonne réponse : 2.} Le verre perd des électrons au profit de la laine ; il se retrouve en déficit d'électrons, donc chargé positivement. Les protons, liés dans les noyaux, ne se déplacent pas dans un solide.
\item \textbf{Bonne réponse : 4.} Loi de Coulomb : $F = k\dfrac{|q_1 q_2|}{r^2}$, donc $F \propto \dfrac{1}{r^2}$.
\item \textbf{Bonne réponse : 3.} Si $r \to 2r$, alors $F \to \dfrac{F}{(2)^2} = \dfrac{F}{4}$ : la force est divisée par 4.
\item \textbf{Bonne réponse : 2.} Champ uniforme : $E = \dfrac{U}{d} = \dfrac{200}{0,050} = \SI{4,0e3}{V.m^{-1}}$.
\item \textbf{Bonne réponse : 2.} Par convention, les lignes de champ sortent des charges positives (sources) et entrent dans les charges négatives (puits). Elles ne se coupent jamais (le champ en un point a une direction unique) et sont courbes en général.
\end{enumerate}
\end{corrige}

\begin{corrige}[Exercice 2 — Vrai ou faux, avec justification]
\begin{enumerate}
\item \textbf{Faux.} La charge est \cle{quantifiée} : $q = n e$ avec $n \in \mathbb{Z}$ et $e = \SI{1,60e-19}{C}$. Elle ne peut pas prendre n'importe quelle valeur réelle.
\item \textbf{Faux.} Deux charges de signes contraires s'\textbf{attirent} ; ce sont les charges de même signe qui se repoussent.
\item \textbf{Vrai.} C'est le \cle{principe de superposition} : $\vec{E}(M) = \sum_i \vec{E}_i(M)$.
\item \textbf{Faux.} La force électrique est \cle{conservative} : son travail ne dépend que des points de départ et d'arrivée : $W_{AB} = q(V_A - V_B)$.
\item \textbf{Faux.} $\vec{F} = q\vec{E}$ : si $q < 0$, $\vec{F}$ est de sens \textbf{opposé} à $\vec{E}$.
\end{enumerate}
\end{corrige}

\begin{corrige}[Exercice 3 — Définitions et restitution de cours]
\begin{enumerate}
\item \textbf{Électrisation par influence :} redistribution des charges libres d'un conducteur neutre sous l'action d'un corps chargé approché, \emph{sans contact} ; le conducteur reste globalement neutre mais présente des charges induites de signes opposés sur ses faces.
\textbf{Charge élémentaire :} $e = \SI{1,602e-19}{C}$, valeur absolue de la charge de l'électron ; toute charge s'écrit $q = n e$ avec $n \in \mathbb{Z}$.
\textbf{Ligne de champ :} courbe orientée tangente en chacun de ses points au vecteur $\vec{E}$ ; elle sort des charges positives et entre dans les charges négatives.
\item \textbf{Loi de Coulomb :} $\vec{F}_{1/2} = k\dfrac{q_1 q_2}{r^2}\,\vec{u}_{12}$, avec $k = \SI{9,0e9}{N.m^2.C^{-2}}$, $q_1, q_2$ en \si{C}, $r$ en \si{m}, $\vec{u}_{12}$ vecteur unitaire dirigé de 1 vers 2. Attraction si $q_1 q_2 < 0$, répulsion si $q_1 q_2 > 0$.
\item $\vec{E}(M) = k\dfrac{q}{r^2}\,\vec{u}_r$ avec $\vec{u}_r = \dfrac{\overrightarrow{OM}}{r}$.
\item $E = \dfrac{U}{d}$. \textbf{Analyse dimensionnelle :} $[U] = \si{J.C^{-1}} = \si{N.m.C^{-1}}$, donc $[U/d] = \si{N.m.C^{-1}.m^{-1}} = \si{N.C^{-1}} = [E]$. La relation est homogène.
\end{enumerate}
\end{corrige}

\begin{corrige}[Exercice 4 — Calculs directs]
\begin{enumerate}
\item $F = k\dfrac{|q_1 q_2|}{r^2} = 9,0\times 10^{9}\times\dfrac{2,0\times 10^{-6}\times 5,0\times 10^{-6}}{(0,10)^2} = 9,0\times 10^{9}\times\dfrac{1,0\times 10^{-11}}{1,0\times 10^{-2}} = \SI{9,0}{N}$.
Signes contraires $\Rightarrow$ force \textbf{attractive}.
\item $E = k\dfrac{|q|}{d^2} = 9,0\times 10^{9}\times\dfrac{4,0\times 10^{-9}}{(0,20)^2} = \dfrac{36}{0,04} = \SI{9,0e2}{V.m^{-1}}$.
Direction : la droite $(OM)$ ; sens : \textbf{de $O$ vers $M$} (champ sortant, car $q>0$).
\item $F_e = eE = 1,60\times 10^{-19}\times 2,0\times 10^{3} = \SI{3,2e-16}{N}$.
$P = m_e g = 9,11\times 10^{-31}\times 9,8 \approx \SI{8,9e-30}{N}$.
$\dfrac{F_e}{P} \approx 3,6\times 10^{13}$ : la force électrique est environ $10^{13}$ fois plus intense que le poids. \textbf{Conclusion :} le poids des particules chargées est négligeable devant la force électrique (canons à électrons, oscilloscopes, accélérateurs).
\end{enumerate}
\end{corrige}

\begin{corrige}[Exercice 5 — Pendule électrostatique]
\begin{enumerate}
\item \textbf{Bilan des forces} appliquées à la sphère (système : sphère ; référentiel terrestre galiléen) :
\begin{itemize}
\item le poids $\vec{P} = m\vec{g}$, vertical, vers le bas ;
\item la tension $\vec{T}$ du fil, dirigée selon le fil vers le point de suspension ;
\item la force électrique $\vec{F}_e = q\vec{E}$, horizontale.
\end{itemize}
\item À l'équilibre : $\vec{T} + \vec{P} + \vec{F}_e = \vec{0}$. En projetant sur l'horizontale et la verticale :
\[ \begin{cases} T\sin\alpha = |q|E \\ T\cos\alpha = mg \end{cases} \Rightarrow \tan\alpha = \frac{|q|E}{mg} \Rightarrow |q| = \frac{mg\tan\alpha}{E}. \]
Application numérique ($m = \SI{5,0e-4}{kg}$, $\tan 15^\circ \approx 0,268$) :
\[ |q| = \frac{5,0\times 10^{-4}\times 9,8\times 0,268}{3,0\times 10^{4}} \approx \SI{4,4e-8}{C} = \SI{44}{nC}. \]
La sphère est attirée vers la baguette positive : elle porte donc une charge \textbf{négative} :
\[ \boxed{q \approx -\SI{4,4e-8}{C}} \]
\end{enumerate}
\textbf{Point de vigilance :} convertir la masse en kilogrammes ; la condition d'équilibre (somme vectorielle nulle) doit être écrite avant toute projection.
\end{corrige}

\begin{corrige}[Exercice 6 — Champ créé par deux charges]
\begin{enumerate}
\item $M$ milieu de $[AB]$ : $AM = MB = \SI{0,15}{m}$.
$\vec{E}_A$ : $q_A > 0$, champ sortant : en $M$, dirigé de $A$ vers $B$.
$\vec{E}_B$ : $q_B < 0$, champ entrant : en $M$, dirigé de $M$ vers $B$.
Les deux vecteurs ont \textbf{même direction et même sens} (vers $B$).
\item $E_A = k\dfrac{|q_A|}{AM^2} = 9,0\times 10^{9}\times\dfrac{3,0\times 10^{-6}}{0,0225} = \SI{1,2e6}{V.m^{-1}}$.
$E_B = k\dfrac{|q_B|}{MB^2} = 9,0\times 10^{9}\times\dfrac{1,0\times 10^{-6}}{0,0225} = \SI{4,0e5}{V.m^{-1}}$.
Par superposition (vecteurs colinéaires de même sens) :
\[ E(M) = E_A + E_B = \SI{1,6e6}{V.m^{-1}}, \quad \text{direction } (AB), \text{ sens } A \to B. \]
\item $\vec{F} = q\vec{E}(M)$ avec $q = -\SI{2,0}{nC}$ :
$F = |q|E = 2,0\times 10^{-9}\times 1,6\times 10^{6} = \SI{3,2e-3}{N}$.
Direction $(AB)$ ; sens \textbf{opposé à $\vec{E}$} (car $q<0$), donc de $B$ vers $A$.
\end{enumerate}
\end{corrige}

\begin{corrige}[Exercice 7 — Condensateur plan et travail de la force électrique]
\begin{enumerate}
\item $E = \dfrac{U}{d} = \dfrac{400}{0,020} = \SI{2,0e4}{V.m^{-1}}$.
\item L'électron va de la plaque $(-)$ vers la plaque $(+)$ : le déplacement est de sens opposé à $\vec{E}$. La force $\vec{F}_e = -e\vec{E}$ est donc de \textbf{même sens que le déplacement} :
\[ W = F_e\, d = eEd = eU = 1,60\times 10^{-19}\times 400 = \SI{6,4e-17}{J}. \]
$W > 0$ : travail \textbf{moteur}.
\item Théorème de l'énergie cinétique (seule la force électrique travaille, le poids étant négligeable) :
\[ W = \tfrac{1}{2}m_e v^2 - 0 \Rightarrow v = \sqrt{\frac{2W}{m_e}} = \sqrt{\frac{2\times 6,4\times 10^{-17}}{9,11\times 10^{-31}}} \approx \SI{1,2e7}{m.s^{-1}}. \]
($v/c \approx 0,04$ : le cadre newtonien reste valable.)
\end{enumerate}
\end{corrige}

\begin{corrige}[Exercice 8 — Lecture d'une carte de lignes de champ]
\begin{enumerate}
\item Les lignes \textbf{sortent} de $A$ donc $q_1 > 0$ ; elles \textbf{arrivent} en $B$ donc $q_2 < 0$.
\item Le nombre de lignes est proportionnel à la valeur absolue de la charge :
\[ \frac{|q_1|}{|q_2|} = \frac{12}{4} = 3 \quad\Rightarrow\quad |q_1| = 3|q_2|. \]
\item Le champ est le plus intense là où les lignes sont les plus resserrées : \textbf{au voisinage de $A$}.
\item Entre $A$ et $B$, $\vec{E}_A$ et $\vec{E}_B$ ont même direction et même sens (de $A$ vers $B$) : leur somme ne peut pas s'annuler. \textbf{Aucun point de champ nul n'existe entre $A$ et $B$} ; il se situe sur la droite $(AB)$, à l'extérieur du segment, du côté de la charge la plus faible en valeur absolue (au-delà de $B$).
\end{enumerate}
\end{corrige}

\begin{corrige}[Exercice 9 — Type Bac : champ et force entre deux charges]
\textbf{Barème indicatif (sur 10) :} 1.a) 1 pt ; 1.b) 1 pt ; 1.c) 2 pts ; 2.a) 0,5 pt ; 2.b) 1 pt ; 2.c) 1,5 pt ; 3) 3 pts (discussion des zones 1,5 pt ; calcul 1,5 pt).
\begin{enumerate}
\item \textbf{Champ en $M$.}
\begin{enumerate}
\item $\vec{E}(M) = k\dfrac{q}{OM^2}\vec{u}_{OM}$ ; par superposition : $\vec{E}(M) = \vec{E}_A(M) + \vec{E}_B(M)$.
\item Schéma : $\vec{E}_A$ et $\vec{E}_B$ colinéaires, tous deux dirigés de $A$ vers $B$ (voir figure du cours).
\item $AM = MB = \SI{0,15}{m}$.
$E_A = 9,0\times 10^{9}\times\dfrac{3,0\times 10^{-6}}{0,0225} = \SI{1,2e6}{V.m^{-1}}$ ; $E_B = 9,0\times 10^{9}\times\dfrac{1,0\times 10^{-6}}{0,0225} = \SI{4,0e5}{V.m^{-1}}$.
$\vec{E}(M)$ : direction $(AB)$, sens $A\to B$, norme $E(M) = E_A + E_B = \SI{1,6e6}{V.m^{-1}}$.
\end{enumerate}
\item \textbf{Force sur $q = -\SI{2,0}{nC}$.}
\begin{enumerate}
\item $\vec{F} = q\vec{E}(M)$.
\item $F = |q|E = 2,0\times 10^{-9}\times 1,6\times 10^{6} = \SI{3,2e-3}{N}$ ; direction $(AB)$, sens $B \to A$ (car $q<0$).
\item \textbf{Vérification directe :} $q_A q < 0$ : attraction, $\vec{F}_A$ dirigée vers $A$ :
$F_A = 9,0\times 10^{9}\times\dfrac{3,0\times 10^{-6}\times 2,0\times 10^{-9}}{0,0225} = \SI{2,4e-3}{N}$.
$q_B q > 0$ : répulsion, $\vec{F}_B$ dirigée aussi vers $A$ :
$F_B = 9,0\times 10^{9}\times\dfrac{1,0\times 10^{-6}\times 2,0\times 10^{-9}}{0,0225} = \SI{0,8e-3}{N}$.
$F = F_A + F_B = \SI{3,2e-3}{N}$ vers $A$. \textbf{Résultat confirmé.}
\end{enumerate}
\item \textbf{Point de champ nul.} Soit $P$ d'abscisse $x$ sur l'axe ($A$ en $0$, $B$ en $0,30$ m).
\begin{itemize}
\item Entre $A$ et $B$ : $\vec{E}_A$ et $\vec{E}_B$ ont même sens (vers $B$) : annulation impossible.
\item À gauche de $A$ : $|q_A| > |q_B|$ et $P$ plus proche de $A$ que de $B$, donc $E_A > E_B$ partout : annulation impossible.
\item À droite de $B$ : les champs sont de sens opposés et on peut avoir $E_A = E_B$ : \textbf{seule zone possible}.
\end{itemize}
Condition : $\dfrac{3,0\times 10^{-6}}{x^2} = \dfrac{1,0\times 10^{-6}}{(x-0,30)^2} \Rightarrow \sqrt{3}\,(x - 0,30) = x$.
\[ x = \frac{0,30\sqrt{3}}{\sqrt{3}-1} = \frac{0,30\times 1,732}{0,732} \approx \SI{0,71}{m}. \]
Donc $BP = x - 0,30 \approx \SI{0,41}{m}$ : \textbf{$P$ est à droite de $B$, à $BP \approx \SI{41}{cm}$}.
\end{enumerate}
\textbf{Points de vigilance :} citer le principe de superposition ; justifier la zone du point de champ nul par une discussion sur les sens des champs (exigée au Bac) ; garder les valeurs absolues dans le calcul des normes.
\end{corrige}

\begin{corrige}[Exercice 10 — Type Bac : trois charges aux sommets d'un triangle équilatéral]
\textbf{Barème indicatif (sur 10) :} 1) 1,5 pt ; 2) 1,5 pt ; 3) 1 pt ; 4) 1,5 pt ; 5) 2,5 pts ; 6) 2 pts.
\begin{enumerate}
\item La médiane du triangle équilatéral vaut $h = \dfrac{a\sqrt{3}}{2}$ ; $G$ est situé aux $\dfrac{2}{3}$ de la médiane à partir du sommet :
\[ r = AG = \frac{2}{3}\times\frac{a\sqrt{3}}{2} = \frac{a}{\sqrt{3}} = \frac{0,30}{1,732} \approx \SI{0,173}{m}. \]
Vérification utile : $r^2 = \dfrac{a^2}{3} = \SI{0,030}{m^2}$.
\item $E_A = k\dfrac{q}{r^2} = 9,0\times 10^{9}\times\dfrac{2,0\times 10^{-6}}{0,030} = \SI{6,0e5}{V.m^{-1}}$.
Par symétrie : $E_B = E_C = \SI{6,0e5}{V.m^{-1}}$.
\item Schéma : trois vecteurs de même norme, dirigés de chaque sommet vers l'extérieur en passant par $G$, faisant entre eux des angles de $120^\circ$.
\item \textbf{Symétrie :} la figure est invariante par rotation de $120^\circ$ autour de $G$ ; le champ résultant doit donc être invariant par cette rotation. Le seul vecteur invariant est le vecteur nul : $\vec{E}(G) = \vec{0}$.
\item Avec $q' = -\SI{2,0}{\micro C}$ en $C$ : $\vec{E}_C$ change de sens (dirigé de $G$ vers $C$), norme inchangée. Par symétrie par rapport à la médiane $(CG)$, les composantes perpendiculaires à $(CG)$ de $\vec{E}_A$ et $\vec{E}_B$ s'annulent ; leurs composantes sur $(CG)$ valent chacune $E_A\cos 30^\circ$ et pointent vers $C$ :
\[ E = 2E_A\cos 30^\circ + E_C = 2\times 6,0\times 10^{5}\times\frac{\sqrt{3}}{2} + 6,0\times 10^{5} \approx 1,04\times 10^{6} + 0,60\times 10^{6} = \SI{1,64e6}{V.m^{-1}}. \]
Direction : la médiane $(CG)$ ; sens : \textbf{de $G$ vers $C$}.
\item Force sur l'électron : $F = eE = 1,60\times 10^{-19}\times 1,64\times 10^{6} \approx \SI{2,6e-13}{N}$.
Direction $(CG)$ ; sens \textbf{opposé à $\vec{E}$} (car $q = -e < 0$), donc de $C$ vers $G$.
\end{enumerate}
\textbf{Points de vigilance :} ne pas oublier que $\cos 30^\circ = \sqrt{3}/2$ ; le raisonnement de symétrie doit être rédigé explicitement ; pour une charge négative, inverser le sens du champ avant de composer.
\end{corrige}

\begin{corrige}[Exercice 11 — Type Bac : expérience de Millikan et quantification de la charge]
\textbf{Barème indicatif (sur 10) :} 1) 1,5 pt ; 2) 1 pt ; 3) 2 pts ; 4) 2 pts ; 5) 2 pts ; 6) 1,5 pt.
\begin{enumerate}
\item Armature supérieure au potentiel le plus élevé : $\vec{E}$ est vertical, dirigé \textbf{vers le bas} (des hauts vers les bas potentiels). Forces sur la gouttelette : le poids $\vec{P} = m\vec{g}$ (vers le bas) et la force électrique $\vec{F}_e = q\vec{E}$ (vers le haut, car $q<0$).
\item $E = \dfrac{U}{d} = \dfrac{500}{0,010} = \SI{5,0e4}{V.m^{-1}}$.
\item Équilibre : $\vec{F}_e + \vec{P} = \vec{0} \Rightarrow |q|E = mg$ :
\[ |q| = \frac{mg}{E} = \frac{1,3\times 10^{-14}\times 9,8}{5,0\times 10^{4}} \approx \SI{2,5e-18}{C}. \]
Le poids étant vers le bas, $\vec{F}_e$ doit être vers le haut ; comme $\vec{E}$ est vers le bas, il faut $q<0$ : la gouttelette est bien chargée \textbf{négativement}.
\item $n = \dfrac{|q|}{e} = \dfrac{2,5\times 10^{-18}}{1,60\times 10^{-19}} \approx 16$.
$n$ est le nombre de charges élémentaires en excès : la gouttelette porte 16 électrons supplémentaires ($q = -16e$).
\textbf{Loi de quantification :} toute charge électrique est un multiple entier de la charge élémentaire : $q = ne$, $n\in\mathbb{Z}$. Le résultat entier confirme expérimentalement cette loi.
\item $E' = \dfrac{U'}{d} = \SI{6,0e4}{V.m^{-1}}$. Nouvelle force électrique (vers le haut) :
$F_e' = |q|E' = 2,5\times 10^{-18}\times 6,0\times 10^{4} = \SI{1,5e-13}{N}$, supérieure au poids $P = mg \approx \SI{1,3e-13}{N}$.
La force résultante est donc dirigée \textbf{vers le haut} (vers l'armature positive) :
\[ F = F_e' - P = (1,5 - 1,3)\times 10^{-13} \approx \SI{2,5e-14}{N}. \]
\item Déplacement de la plaque $(-)$ vers la plaque $(+)$ : $\vec{F}_e'$ et le déplacement sont de même sens :
\[ W = F_e'\, d = |q|E'd = |q|U' = 2,5\times 10^{-18}\times 600 \approx \SI{1,5e-15}{J}. \]
$W > 0$ : travail \textbf{moteur}.
\end{enumerate}
\textbf{Points de vigilance :} justifier le signe de $q$ par le sens de $\vec{F}_e$ (argument attendu au Bac) ; pour le travail, on peut aussi écrire $W = q(V_A - V_B)$ en faisant attention aux signes ; arrondir $n$ à l'entier le plus proche et commenter ce résultat.
\end{corrige}

\newpage

\coursec{Fiche bilan}

\begin{popfigure}
\begin{tikzpicture}[
  mindmap,
  grow cyclic,
  every node/.style={concept, font=\small},
  root concept/.append style={concept color=popblue, font=\bfseries\small, minimum size=3.2cm},
  level 1/.append style={level distance=3.4cm, sibling angle=60, font=\scriptsize},
  concept color=popblue
]
\node[root concept]{Interactions électrostatiques}
  child[concept color=poporange]{ node[align=center] {Charges et électrisation\\ $e=\num{1,6e-19}$ C\\ $q=ne$} }
  child[concept color=popgreen]{ node[align=center] {Loi de Coulomb\\ $F=k\dfrac{|q_1 q_2|}{r^2}$} }
  child[concept color=poppurple]{ node[align=center] {Champ $\vec{E}$\\ $\vec{E}=\dfrac{\vec{F}}{q}$\\ $E=k\dfrac{|Q|}{r^2}$} }
  child[concept color=poppink]{ node[align=center] {Superposition\\ $\vec{E}=\sum \vec{E}_i$\\ lignes de champ} }
  child[concept color=popgold]{ node[align=center] {Champ uniforme\\ $E=\dfrac{U}{d}$\\ condensateur plan} }
  child[concept color=popblue!60]{ node[align=center] {Travail et potentiel\\ $W_{AB}=qU_{AB}$\\ $U_{AB}=V_A-V_B$} };
\end{tikzpicture}
\legende{Carte mentale de la leçon : les six piliers de l'électrostatique.}
\end{popfigure}

\begin{bilan}
\textbf{Définitions clés}
\begin{itemize}
  \item \cle{Charge électrique} $q$ : grandeur algébrique (en coulombs, C), \cle{quantifiée} : $q = n\,e$ avec $n \in \mathbb{Z}$ et $e = \num{1,602e-19}$ C ; elle se \cle{conserve} dans tout système isolé.
  \item \cle{Électrisation} : par \emph{frottement} (transfert d'électrons), par \emph{contact} (répartition des charges), par \emph{influence} (sans contact, redistribution).
  \item \cle{Champ électrostatique} $\vec{E}$ en un point $M$ : $\vec{E}(M) = \dfrac{\vec{F}}{q}$, où $\vec{F}$ est la force subie par une charge témoin $q$ placée en $M$. Unité : V/m (ou N/C).
  \item \cle{Potentiel électrique} $V$ (en volts) : défini à une constante près ; $U_{AB} = V_A - V_B$.
\end{itemize}

\textbf{Formules à connaître par c\oe ur}
\begin{center}
\begin{tabularx}{\linewidth}{|l|X|l|}
\hline
\rowcolor{popblueL} \textbf{Grandeur} & \textbf{Formule} & \textbf{Unités SI} \\
\hline
Force de Coulomb & $F = k\,\dfrac{|q_1 q_2|}{r^2}$, \quad $k = \num{9,0e9}$ & N ; $q$ en C, $r$ en m \\
\hline
Champ d'une charge ponctuelle & $E(M) = k\,\dfrac{|Q|}{r^2}$, radial, sortant si $Q>0$ & V/m \\
\hline
Principe de superposition & $\vec{E}(M) = \vec{E}_1(M) + \vec{E}_2(M) + \cdots$ & V/m \\
\hline
Champ uniforme (condensateur plan) & $E = \dfrac{U}{d}$ & $U$ en V, $d$ en m \\
\hline
Travail de la force électrique & $W_{AB}(\vec{F}) = q\,U_{AB} = q(V_A - V_B)$ & J \\
\hline
Lien champ--potentiel (uniforme) & $U_{AB} = E \cdot d$ (le long d'une ligne de champ) & V \\
\hline
\end{tabularx}
\end{center}

\textbf{Méthodes express}
\begin{itemize}
  \item \textbf{Force entre deux charges} : convertir $q$ en C et $r$ en m, appliquer Coulomb pour la \emph{valeur}, puis conclure sur le \emph{sens} : charges de même signe $\Rightarrow$ répulsion ; signes opposés $\Rightarrow$ attraction.
  \item \textbf{Champ résultant en un point} : dessiner chaque $\vec{E}_i$ (sortant de $Q_i>0$, rentrant vers $Q_i<0$), calculer les normes, projeter sur des axes, additionner les composantes.
  \item \textbf{Point de champ nul} : chercher sur la droite joignant les charges ; entre deux charges de même signe, à l'extérieur sinon.
  \item \textbf{Condensateur plan} : $E = U/d$ ; le champ est uniforme, perpendiculaire aux plaques, dirigé de la plaque positive vers la plaque négative ; les lignes de champ vont dans le sens des potentiels décroissants.
  \item \textbf{Travail} : $W_{AB} = qU_{AB}$ ; vérifier le signe (travail moteur si $qU_{AB}>0$) ; le travail ne dépend pas du chemin suivi (force conservative).
\end{itemize}
\end{bilan}

\begin{aretenir}[Checklist avant le Bac]
\begin{itemize}
  \item $\square$ Je sais citer les trois modes d'électrisation et les illustrer par une expérience simple.
  \item $\square$ Je connais la valeur de la charge élémentaire $e = \num{1,6e-19}$ C et j'utilise $q = ne$.
  \item $\square$ J'énonce la loi de Coulomb avec la valeur de $k = \num{9,0e9}$ N\,m$^2$\,C$^{-2}$ et ses conditions d'application (charges ponctuelles).
  \item $\square$ Je détermine le sens (attraction/répulsion) d'une force électrostatique à partir des signes des charges.
  \item $\square$ Je convertis systématiquement les distances en mètres et les charges en coulombs avant tout calcul.
  \item $\square$ Je trace le champ créé par une charge ponctuelle (radial, sortant si $Q>0$) et j'applique le principe de superposition vectorielle.
  \item $\square$ Je sais qu'une ligne de champ est tangente à $\vec{E}$ et orientée dans le sens des potentiels décroissants.
  \item $\square$ Je décris un condensateur plan et j'utilise $E = U/d$ avec $d$ en mètres.
  \item $\square$ Je calcule $W_{AB} = qU_{AB}$ en veillant au signe et je sais que ce travail est indépendant du chemin.
  \item $\square$ Je vérifie l'homogénéité de chaque résultat et je l'accompagne de son unité SI.
\end{itemize}
\end{aretenir}

\findecours

\end{document}
